\documentclass[a4paper,12pt]{article}
\usepackage[utf8]{inputenc}
\usepackage[T1]{fontenc}
\usepackage[italian]{babel}
\usepackage[sfdefault]{atkinson}
\renewcommand*\familydefault{\sfdefault}
\usepackage{float}
\usepackage{microtype}
\usepackage{geometry}
\usepackage{setspace}
\usepackage{enumitem}
\usepackage{titlesec}
\usepackage{chngpage}
\usepackage{tocloft}
\usepackage{longtable}
\usepackage{ltablex}
\usepackage{array}
\usepackage{graphicx}
\usepackage{fancyhdr}
\usepackage[table]{xcolor}
\usepackage{color,soul}
\usepackage[most]{tcolorbox}
\usepackage{nameref}
\usepackage[colorlinks=true]{hyperref}

% --- INSERISCI QUESTO NEL PREAMBOLO DEL TUO DOCUMENTO ---
\usepackage{amssymb} % per il simbolo \checkmark
\usepackage{pifont}

% Macro per semplificare e alleggerire il codice della tabella
\newcommand{\ok}{\textcolor{green!60!black}{$\checkmark$ Implementato}}
\newcommand{\ko}{\textcolor{red!80!black}{$\times$ Mancante}}
\newcommand{\warn}[1]{\textcolor{orange!90!black}{\textbf{!} Multiplo (#1)}}
% Use URL-style line breaking to keep long test file names inside table borders.
\newcommand{\f}[1]{\nolinkurl{#1}\newline}
\newcommand{\q}[1]{\textit{#1}}
\newcommand{\sep}{\newline \vspace{2mm} \hrule \vspace{2mm}}
% --------------------------------------------------------

% Colori ZPUS - Verde, Nero, Bianco
\definecolor{zpusgreen}{RGB}{4, 138, 55}
\definecolor{zpusdarkgreen}{RGB}{0, 100, 0}
\definecolor{zpusblack}{RGB}{0, 0, 0}
\definecolor{zpuswhite}{RGB}{255, 255, 255}
\definecolor{zpuslightgray}{RGB}{245, 245, 245}

\hypersetup{
    linkcolor=black,
    urlcolor=blue
}

\definecolor{lightblack}{gray}{0.35}
\newcommand{\glossario}[1]{\textit{#1}\textsubscript{\textbf{\textit{\textcolor{lightblack}{G}}}}}
\newcommand{\ped}[1]{\textsubscript{#1}}

\pagestyle{fancy}
\setlength{\headwidth}{\textwidth}
\fancyhfoffset[L,R]{0pt}
\lhead{\rightmark}
\rhead{7-ZPUs}
\lfoot{Piano di Qualifica}
\rfoot{\thepage}
\cfoot{}
\renewcommand{\headrulewidth}{0.8pt}
\renewcommand{\footrulewidth}{0.8pt}

\renewcommand{\contentsname}{Indice}
\renewcommand{\listoftables}{Elenco delle tabelle}
\renewcommand{\listoffigures}{Elenco delle immagini}

\geometry{margin=2.5cm}
\setstretch{1.1}

\titleclass{\subsubsubsection}{straight}[\subsection]

\newcounter{subsubsubsection}[subsubsection]
\renewcommand\thesubsubsubsection{\thesubsubsection.\arabic{subsubsubsection}}
\renewcommand\theparagraph{\thesubsubsubsection.\arabic{paragraph}} % optional; useful if paragraphs are to be numbered

\titleformat{\subsubsubsection}
  {\normalfont\normalsize\bfseries}{\thesubsubsubsection}{1em}{}
\titlespacing*{\subsubsubsection}
{0pt}{3.25ex plus 1ex minus .2ex}{1.5ex plus .2ex}

\makeatletter
\renewcommand\paragraph{\@startsection{paragraph}{5}{\z@}%
  {3.25ex \@plus1ex \@minus.2ex}%
  {-1em}%
  {\normalfont\normalsize\bfseries}}
\renewcommand\subparagraph{\@startsection{subparagraph}{6}{\parindent}%
  {3.25ex \@plus1ex \@minus .2ex}%
  {-1em}%
  {\normalfont\normalsize\bfseries}}
\def\toclevel@subsubsubsection{4}
\def\toclevel@paragraph{5}
%\def\toclevel@paragraph{6}
\def\toclevel@subparagraph{6}
\def\l@subsubsubsection{\@dottedtocline{4}{7em}{4em}}
\def\l@paragraph{\@dottedtocline{5}{10em}{5em}}
\def\l@subparagraph{\@dottedtocline{6}{14em}{6em}}
\makeatother

\setcounter{secnumdepth}{4}
\setcounter{tocdepth}{4}

\titleformat{\section}{\LARGE\bfseries}{\thesection}{1em}{}
\titleformat{\subsection}{\Large\bfseries}{\thesubsection}{1em}{}
\titleformat{\subsubsection}{\large\bfseries}{\thesubsubsection}{1em}{}
\titleformat{\paragraph}{\large\bfseries}{\theparagraph}{1em}{}
\titleformat{\subparagraph}{\normalsize\bfseries}{\thesubparagraph}{1em}{}

\setlength\parindent{0pt}

\begin{document}

\begin{center}
    \includegraphics[width=9.5cm]{../../assets/logo7ZPUs.jpg}\\
    \small\hspace{10cm} 7zpus.swe@gmail.com\\
    \vspace{0.5cm}
    \LARGE \textbf{Appendice - Piano di Qualifica}\\
    \vspace{0.3cm}
    \normalsize \textbf{Versione:} 2.0.0\\
    \textbf{Responsabile:} Georgescu Diana \\
    \textbf{Destinatari:} Sanmarco Informatica S.p.A., Prof. Tullio Vardanega, Prof. Riccardo Cardin, Team 7-ZPUs \\
\end{center}

\vspace{0.3cm}
\hrule
\vspace{0.5cm}


\section{Test di unità}

\setlength{\LTleft}{-0.4cm}
\renewcommand{\arraystretch}{1.2}
\rowcolors{2}{gray!15}{white}
\begin{longtable}{|p{2cm}|p{3.5cm}|p{9.5cm}|}
\hline
\rowcolor{zpusgreen!30}
\textbf{ID Test} & \textbf{Stato} & \textbf{Risultato atteso} \\
\hline
\endfirsthead
\rowcolor{zpusgreen!30}
\textbf{ID Test} & \textbf{Stato} & \textbf{Risultato atteso} \\
\hline
\endhead
TU-1 & \ok & saved.getId() non deve essere null \\
\hline
TU-2 & \ok & saved.getUuid() deve essere "dip-1" \\
\hline
TU-3 & \ok & saved.getIntegrityStatus() deve essere IntegrityStatusEnum.UNKNOWN \\
\hline
TU-4 & \ok & count deve essere 1 \\
\hline
TU-5 & \ok & second.getId() deve essere first.getId() \\
\hline
TU-6 & \ok & second.getIntegrityStatus() deve essere IntegrityStatusEnum.UNKNOWN \\
\hline
TU-7 & \ok & byId?.getUuid() deve essere "dip-3" \\
\hline
TU-8 & \ok & byUuid?.getId() deve essere saved.getId() \\
\hline
TU-9 & \ok & valid.map((d) => d.getUuid()) deve essere uguale a ["dip-a"] \\
\hline
TU-10 & \ok & invalid.map((d) => d.getUuid()) deve essere uguale a ["dip-b"] \\
\hline
TU-11 & \ok & saved.getId() non deve essere null \\
\hline
TU-12 & \ok & found?.getUuid() deve essere "dc-uuid" \\
\hline
TU-13 & \ok & found?.getDipUuid() deve essere "dip-uuid" \\
\hline
TU-14 & \ok & rows deve avere lunghezza 1 \\
\hline
TU-15 & \ok & rows[0].name deve essere "Nome 2" \\
\hline
TU-16 & \ok & dao.getByDipId(dipId) deve avere lunghezza 2 \\
\hline
TU-17 & \ok & dao.getByStatus(IntegrityStatusEnum.VALID) deve avere lunghezza 1 \\
\hline
TU-18 & \ok & dao.search("Contr") deve avere lunghezza 1 \\
\hline
TU-19 & \ok & dao.search("NON-EXISTING") deve avere lunghezza 0 \\
\hline
TU-20 & \ok & integritya deve essere IntegrityStatusEnum.VALID \\
\hline
TU-21 & \ok & integrityb deve essere IntegrityStatusEnum.VALID \\
\hline
TU-22 & \ok & integritya2 deve essere IntegrityStatusEnum.INVALID \\
\hline
TU-23 & \ok & integrityb2 deve essere IntegrityStatusEnum.INVALID \\
\hline
TU-24 & \ok & found non deve essere null \\
\hline
TU-25 & \ok & found?.getUuid() deve essere "doc-a" \\
\hline
TU-26 & \ok & found?.getProcessId() deve essere processId \\
\hline
TU-27 & \ok & found?.getMetadata().findNodeByName("Titolo")?.getStringValue() deve essere "A" \\
\hline
TU-28 & \ok & dao.getByProcessId(processId) deve avere lunghezza 2 \\
\hline
TU-29 & \ok & dao.getByStatus(IntegrityStatusEnum.VALID) deve avere lunghezza 1 \\
\hline
TU-30 & \ok & dao.getByStatus(IntegrityStatusEnum.INVALID) deve avere lunghezza 1 \\
\hline
TU-31 & \ok & results deve avere lunghezza 2 \\
\hline
TU-32 & \ok & results[0].document.getUuid() deve essere "doc-sem-a" \\
\hline
TU-33 & \ok & results[0].score deve essere maggiore di results[1].score \\
\hline
TU-34 & \ok & dao.getIndexedDocumentsCount() deve essere 1 \\
\hline
TU-35 & \ok & found non deve essere null \\
\hline
TU-36 & \ok & found?.getFilename() deve essere "main.xml" \\
\hline
TU-37 & \ok & found?.getDocumentId() deve essere documentId \\
\hline
TU-38 & \ok & found?.getIsMain() deve essere true \\
\hline
TU-39 & \ok & rows deve avere lunghezza 1 \\
\hline
TU-40 & \ok & rows[0].id deve essere updated.getId() \\
\hline
TU-41 & \ok & rows[0].filename deve essere "main-v2.xml" \\
\hline
TU-42 & \ok & rows[0].hash deve essere "hash-new" \\
\hline
TU-43 & \ok & dao.getByDocumentId(documentId) deve avere lunghezza 2 \\
\hline
TU-44 & \ok & dao.getByStatus(IntegrityStatusEnum.VALID) deve avere lunghezza 1 \\
\hline
TU-45 & \ok & dao.getByStatus(IntegrityStatusEnum.INVALID) deve avere lunghezza 1 \\
\hline
TU-46 & \ok & dip.getId() deve essere 10 \\
\hline
TU-47 & \ok & dip.getUuid() deve essere "dip-uuid" \\
\hline
TU-48 & \ok & dip.getIntegrityStatus() deve essere IntegrityStatusEnum.VALID \\
\hline
TU-49 & \ok & dip.getIntegrityStatus() deve essere IntegrityStatusEnum.UNKNOWN \\
\hline
TU-50 & \ok & dto.id deve essere 42 \\
\hline
TU-51 & \ok & dto.uuid deve essere "dip-uuid" \\
\hline
TU-52 & \ok & dto.integrityStatus deve essere IntegrityStatusEnum.INVALID \\
\hline
TU-53 & \ok & () => DipMapper.toDTO(dip) deve lanciare "Cannot convert Dip to DTO: id is null" \\
\hline
TU-54 & \ok & model deve essere uguale a { uuid: "dip-uuid", integrityStatus: IntegrityStatusEnum.VALID, } \\
\hline
TU-55 & \ok & entity.getId() deve essere 7 \\
\hline
TU-56 & \ok & entity.getDipId() deve essere 3 \\
\hline
TU-57 & \ok & entity.getDipUuid() deve essere "dip-uuid" \\
\hline
TU-58 & \ok & entity.getUuid() deve essere "dc-uuid" \\
\hline
TU-59 & \ok & entity.getIntegrityStatus() deve essere IntegrityStatusEnum.VALID \\
\hline
TU-60 & \ok & entity.getIntegrityStatus() deve essere IntegrityStatusEnum.UNKNOWN \\
\hline
TU-61 & \ok & dto deve essere uguale a { id: 12, dipId: 5, uuid: "dc-uuid", name: "Classe", timestamp: "2026-01-01T00:00:00Z", integrityStatus: IntegrityStatusEnum.INVALID, } \\
\hline
TU-62 & \ok & dto.dipId deve essere -1 \\
\hline
TU-63 & \ok & () => DocumentClassMapper.toDTO(entity) deve lanciare "Cannot convert to DTO: DocumentClass entity is not yet persisted and has no ID." \\
\hline
TU-64 & \ok & result deve essere uguale a { id: 21, uuid: "dc-uuid", integrityStatus: IntegrityStatusEnum.VALID, name: "Classe Risorse Umane", type: "CLASSE", } \\
\hline
TU-65 & \ok & () => DocumentClassMapper.toSearchResult(entity) deve lanciare "Cannot convert to SearchResult: DocumentClass entity is not yet persisted and has no ID." \\
\hline
TU-66 & \ok & model deve essere uguale a { dipUuid: "dip-uuid", uuid: "dc-uuid", integrityStatus: IntegrityStatusEnum.VALID, name: "Classe", timestamp: "2026-01-01T00:00:00Z", } \\
\hline
TU-67 & \ok & entity.getId() deve essere 4 \\
\hline
TU-68 & \ok & entity.getUuid() deve essere "doc-uuid" \\
\hline
TU-69 & \ok & entity.getProcessId() deve essere 8 \\
\hline
TU-70 & \ok & entity.getProcessUuid() deve essere "proc-uuid" \\
\hline
TU-71 & \ok & entity.getIntegrityStatus() deve essere IntegrityStatusEnum.VALID \\
\hline
TU-72 & \ok & entity.getMetadata().getName() deve essere "Documento" \\
\hline
TU-73 & \ok & entity.getMetadata().getChildren() deve avere lunghezza 1 \\
\hline
TU-74 & \ok & DocumentMapper.fromPersistence( missingStatus, metadataRows, ).getIntegrityStatus() deve essere IntegrityStatusEnum.UNKNOWN \\
\hline
TU-75 & \ok & DocumentMapper.fromPersistence( invalidStatus, metadataRows, ).getIntegrityStatus() deve essere IntegrityStatusEnum.UNKNOWN \\
\hline
TU-76 & \ok & dto.id deve essere 10 \\
\hline
TU-77 & \ok & dto.processId deve essere 20 \\
\hline
TU-78 & \ok & dto.uuid deve essere "doc-uuid" \\
\hline
TU-79 & \ok & dto.integrityStatus deve essere IntegrityStatusEnum.INVALID \\
\hline
TU-80 & \ok & dto.metadata.name deve essere "Documento" \\
\hline
TU-81 & \ok & () => DocumentMapper.toDTO(entity) deve lanciare "Cannot convert to DTO: Document entity is not yet persisted and has no ID." \\
\hline
TU-82 & \ok & dto.processId deve essere -1 \\
\hline
TU-83 & \ok & model.uuid deve essere "doc-uuid" \\
\hline
TU-84 & \ok & model.integrityStatus deve essere IntegrityStatusEnum.VALID \\
\hline
TU-85 & \ok & model.processUuid deve essere "proc-uuid" \\
\hline
TU-86 & \ok & model.metadata.map((m) => m.getName()) deve essere uguale a ["Titolo", "Nome"] \\
\hline
TU-87 & \ok & Array.isArray(ruolo) deve essere true \\
\hline
TU-88 & \ok & ruolo deve avere lunghezza 2 \\
\hline
TU-89 & \ok & ruolo[0].Altro.PF.Cognome deve essere "Rossi" \\
\hline
TU-90 & \ok & ruolo[0].Altro.PF.Nome deve essere "Paolo" \\
\hline
TU-91 & \ok & ruolo[1].ResponsabileGestioneDocumentale.PF.Cognome deve essere "Colombo" \\
\hline
TU-92 & \ok & soggetti non deve essere null \\
\hline
TU-93 & \ok & ruoloNodes deve avere lunghezza 2 \\
\hline
TU-94 & \ok & root.findNodeByName("Altro")?.findNodeByName("Cognome")?.getStringValue() deve essere "Rossi" \\
\hline
TU-95 & \ok & root .findNodeByName("ResponsabileGestioneDocumentale") ?.findNodeByName("Cognome") ?.getStringValue() deve essere "Colombo" \\
\hline
TU-96 & \ok & entity.getId() deve essere 9 \\
\hline
TU-97 & \ok & entity.getUuid() deve essere "file-uuid" \\
\hline
TU-98 & \ok & entity.getFilename() deve essere "main.xml" \\
\hline
TU-99 & \ok & entity.getIsMain() deve essere true \\
\hline
TU-100 & \ok & entity.getDocumentId() deve essere 3 \\
\hline
TU-101 & \ok & entity.getDocumentUuid() deve essere "doc-uuid" \\
\hline
TU-102 & \ok & entity.getIntegrityStatus() deve essere IntegrityStatusEnum.VALID \\
\hline
TU-103 & \ok & entity.getIsMain() deve essere false \\
\hline
TU-104 & \ok & entity.getIntegrityStatus() deve essere IntegrityStatusEnum.UNKNOWN \\
\hline
TU-105 & \ok & dto deve essere uguale a { id: 20, documentId: 33, filename: "main.xml", path: "/pkg/main.xml", hash: "hash", integrityStatus: IntegrityStatusEnum.INVALID, isMain: true, } \\
\hline
TU-106 & \ok & dto.documentId deve essere -1 \\
\hline
TU-107 & \ok & () => FileMapper.toDTO(entity) deve lanciare "Cannot convert to DTO: File entity is not yet persisted and has no ID." \\
\hline
TU-108 & \ok & model deve essere uguale a { filename: "main.xml", path: "/pkg/main.xml", hash: "hash", integrityStatus: IntegrityStatusEnum.VALID, isMain: 0, uuid: "file-uuid", documentUuid: "doc-uuid", } \\
\hline
TU-109 & \ok & MetadataKeyMapper.toPascalCase( "DocumentoInformatico.Soggetti.Ruolo.Destinatario.PG.CodiceFiscale\_PartitaIva", ) deve essere "DocumentoInformatico.Soggetti.Ruolo.Destinatario.PG.CodiceFiscale\_PartitaIva" \\
\hline
TU-110 & \ok & MetadataKeyMapper.toPascalCase("soggetti.codice\_fiscale") deve essere "Soggetti.CodiceFiscale" \\
\hline
TU-111 & \ok & MetadataKeyMapper.toPascalCase("soggetti.codice-fiscale") deve essere "Soggetti.CodiceFiscale" \\
\hline
TU-112 & \ok & group.logicOperator deve essere "AND" \\
\hline
TU-113 & \ok & group.items deve essere uguale a [ { path: "Soggetto.Nome", operator: "EQ", value: "Mario", }, ] \\
\hline
TU-114 & \ok & group.items deve essere uguale a [ { path: "Flags.Attivo", operator: "EQ", value: false, }, { path: "Metrics.Count", operator: "EQ", value: 0, }, ] \\
\hline
TU-115 & \ok & nested.items[0].path deve essere "Destinatario.PG.CodiceFiscale\_PartitaIva" \\
\hline
TU-116 & \ok & () => MetadataMapper.fromFlatRows([]) deve lanciare "Cannot build metadata tree from empty rows" \\
\hline
TU-117 & \ok & root non deve essere null \\
\hline
TU-118 & \ok & root?.getName() deve essere "root" \\
\hline
TU-119 & \ok & root?.getType() deve essere MetadataType.COMPOSITE \\
\hline
TU-120 & \ok & root?.getChildren() deve avere lunghezza 2 \\
\hline
TU-121 & \ok & subject?.getName() deve essere "Soggetto" \\
\hline
TU-122 & \ok & subject?.getType() deve essere MetadataType.COMPOSITE \\
\hline
TU-123 & \ok & subject?.getChildren() deve avere lunghezza 1 \\
\hline
TU-124 & \ok & subject?.getChildren()[0].getName() deve essere "Nome" \\
\hline
TU-125 & \ok & subject?.getChildren()[0].getStringValue() deve essere "Mario" \\
\hline
TU-126 & \ok & root?.getType() deve essere MetadataType.COMPOSITE \\
\hline
TU-127 & \ok & root?.getChildren() deve avere lunghezza 1 \\
\hline
TU-128 & \ok & root?.getChildren()[0].getType() deve essere MetadataType.NUMBER \\
\hline
TU-129 & \ok & () => MetadataMapper.fromFlatRows(rows) deve lanciare "Invalid metadata tree: missing parent row for id 100 (parent\_id=999)" \\
\hline
TU-130 & \ok & () => MetadataMapper.fromFlatRows(rows) deve lanciare "Invalid metadata tree: duplicate row id 1" \\
\hline
TU-131 & \ok & () => MetadataMapper.fromFlatRows(rows) deve lanciare "Invalid metadata tree: no root nodes found" \\
\hline
TU-132 & \ok & () => MetadataMapper.fromFlatRows(rows) deve lanciare "Invalid metadata type: invalid" \\
\hline
TU-133 & \ok & flat deve avere lunghezza 1 \\
\hline
TU-134 & \ok & flat[0].getName() deve essere "Titolo" \\
\hline
TU-135 & \ok & flat[0].getStringValue() deve essere "Doc A" \\
\hline
TU-136 & \ok & flat deve avere lunghezza 2 \\
\hline
TU-137 & \ok & flat.map((node) => node.getName()) deve essere uguale a ["A", "C"] \\
\hline
TU-138 & \ok & entity.getId() deve essere 3 \\
\hline
TU-139 & \ok & entity.getDocumentClassId() deve essere 9 \\
\hline
TU-140 & \ok & entity.getDocumentClassUuid() deve essere "dc-uuid" \\
\hline
TU-141 & \ok & entity.getUuid() deve essere "proc-uuid" \\
\hline
TU-142 & \ok & entity.getIntegrityStatus() deve essere IntegrityStatusEnum.VALID \\
\hline
TU-143 & \ok & entity.getMetadata().getChildren() deve avere lunghezza 1 \\
\hline
TU-144 & \ok & ProcessMapper.fromPersistence(missingStatus, metadataRows).getIntegrityStatus() deve essere IntegrityStatusEnum.UNKNOWN \\
\hline
TU-145 & \ok & ProcessMapper.fromPersistence(invalidStatus, metadataRows).getIntegrityStatus() deve essere IntegrityStatusEnum.UNKNOWN \\
\hline
TU-146 & \ok & dto.id deve essere 8 \\
\hline
TU-147 & \ok & dto.documentClassId deve essere 11 \\
\hline
TU-148 & \ok & dto.uuid deve essere "proc-uuid" \\
\hline
TU-149 & \ok & dto.integrityStatus deve essere IntegrityStatusEnum.INVALID \\
\hline
TU-150 & \ok & dto.metadata.name deve essere "Process" \\
\hline
TU-151 & \ok & dto.documentClassId deve essere -1 \\
\hline
TU-152 & \ok & () => ProcessMapper.toDTO(entity) deve lanciare "Cannot convert to DTO: Process entity is not yet persisted and has no ID." \\
\hline
TU-153 & \ok & model.documentClassUuid deve essere "dc-uuid" \\
\hline
TU-154 & \ok & model.uuid deve essere "proc-uuid" \\
\hline
TU-155 & \ok & model.integrityStatus deve essere IntegrityStatusEnum.VALID \\
\hline
TU-156 & \ok & model.metadata.map((m) => m.getName()) deve essere uguale a ["Fase", "Step"] \\
\hline
TU-157 & \ok & rows deve avere lunghezza 4 \\
\hline
TU-158 & \ok & rows[0] deve essere uguale a { owner\_id: 1, parent\_id: null, name: "Root", value: "", type: "COMPOSITE", } \\
\hline
TU-159 & \ok & rows[1].name deve essere "Title" \\
\hline
TU-160 & \ok & rows[2].name deve essere "Soggetto" \\
\hline
TU-161 & \ok & rows[3].name deve essere "Nome" \\
\hline
TU-162 & \ok & loaded deve avere lunghezza 2 \\
\hline
TU-163 & \ok & loaded.map((r) => r.name) deve essere uguale a ["A", "B"] \\
\hline
TU-164 & \ok & count deve essere 2 \\
\hline
TU-165 & \ok & found non deve essere null \\
\hline
TU-166 & \ok & found?.getUuid() deve essere "proc-a" \\
\hline
TU-167 & \ok & found?.getDocumentClassId() deve essere documentClassId \\
\hline
TU-168 & \ok & found?.getMetadata().getName() deve essere "Process" \\
\hline
TU-169 & \ok & found?.getMetadata().findNodeByName("Versione")?.getStringValue() deve essere "v1" \\
\hline
TU-170 & \ok & dao.getByDocumentClassId(documentClassId) deve avere lunghezza 2 \\
\hline
TU-171 & \ok & dao.getByStatus(IntegrityStatusEnum.VALID) deve avere lunghezza 1 \\
\hline
TU-172 & \ok & dao.getByStatus(IntegrityStatusEnum.INVALID) deve avere lunghezza 1 \\
\hline
TU-173 & \ok & result.sql deve essere "" \\
\hline
TU-174 & \ok & result.params deve essere uguale a [] \\
\hline
TU-175 & \ok & result.sql deve essere "(json\_extract(document.metadata, '\$.Doc.Type') = ? COLLATE NOCASE)" \\
\hline
TU-176 & \ok & result.params deve essere uguale a ["Invoice"] \\
\hline
TU-177 & \ok & result.sql deve essere "(json\_extract(document.metadata, '\$.DocumentoInformatico.Riservato') = ? COLLATE NOCASE)" \\
\hline
TU-178 & \ok & result.params deve essere uguale a [1] \\
\hline
TU-179 & \ok & result.sql deve essere "(json\_extract(document.metadata, '\$.Amount') > ?)" \\
\hline
TU-180 & \ok & result.params deve essere uguale a [100] \\
\hline
TU-181 & \ok & result.sql deve essere "(json\_extract(document.metadata, '\$.Amount') < ?)" \\
\hline
TU-182 & \ok & result.params deve essere uguale a [50] \\
\hline
TU-183 & \ok & result.sql deve essere "(json\_extract(document.metadata, '\$.Name') LIKE ? COLLATE NOCASE)" \\
\hline
TU-184 & \ok & result.params deve essere uguale a ["\%test\%"] \\
\hline
TU-185 & \ok & result.sql deve essere "(json\_extract(document.metadata, '\$.Status') IN (?, ?, ?))" \\
\hline
TU-186 & \ok & result.params deve essere uguale a ["A", "B", "C"] \\
\hline
TU-187 & \ok & result.sql deve essere "" \\
\hline
TU-188 & \ok & result.params deve essere uguale a [] \\
\hline
TU-189 & \ok & result.sql deve essere "(json\_extract(document.metadata, '\$.Doc.Type') = ? COLLATE NOCASE AND json\_extract(document.metadata, '\$.Status') IN (?, ?))" \\
\hline
TU-190 & \ok & result.params deve essere uguale a ["Invoice", "Paid", "Sent"] \\
\hline
TU-191 & \ok & result.sql deve essere "(json\_extract(document.metadata, '\$.Type') = ? COLLATE NOCASE OR json\_extract(document.metadata, '\$.Type') = ? COLLATE NOCASE)" \\
\hline
TU-192 & \ok & result.params deve essere uguale a ["A", "B"] \\
\hline
TU-193 & \ok & result.sql deve essere "(json\_extract(document.metadata, '\$.Year') > ? AND (json\_extract(document.metadata, '\$.Dept') = ? COLLATE NOCASE OR json\_extract(document.metadata, '\$.Dept') = ? COLLATE NOCASE))" \\
\hline
TU-194 & \ok & result.params deve essere uguale a [2020, "IT", "HR"] \\
\hline
TU-195 & \ok & result.sql deve essere "(EXISTS (SELECT 1 FROM json\_each(document.metadata, '\$.Items') AS el WHERE (json\_extract(el.value, '\$.Code') = ? COLLATE NOCASE AND json\_extract(el.value, '\$.Price') > ?)))" \\
\hline
TU-196 & \ok & result.params deve essere uguale a ["XYZ", 50] \\
\hline
TU-197 & \ok & () => builder.build(query) deve lanciare /ELEM\_MATCH requires a nested SearchGroup value/ \\
\hline
TU-198 & \ok & result.sql deve essere "(json\_extract(document.metadata, '\$.Name') = ? COLLATE NOCASE)" \\
\hline
TU-199 & \ok & result.params deve essere uguale a ["Valid"] \\
\hline
TU-200 & \ok & result.sql deve essere "(json\_extract(document.metadata, '\$.Dept') = ? COLLATE NOCASE)" \\
\hline
TU-201 & \ok & result.params deve essere uguale a ["IT"] \\
\hline
TU-202 & \ok & result.sql deve essere "(json\_extract(document.metadata, '\$.Status') IN (?, ?))" \\
\hline
TU-203 & \ok & result.params deve essere uguale a ["Paid", "Sent"] \\
\hline
TU-204 & \ok & result.sql deve essere "" \\
\hline
TU-205 & \ok & result.params deve essere uguale a [] \\
\hline
TU-206 & \ok & result.sql deve essere "(json\_extract(document.metadata, '\$.My.Path') = ? COLLATE NOCASE)" \\
\hline
TU-207 & \ok & result.params deve essere uguale a ["X"] \\
\hline
TU-208 & \ok & () => builder.build(badQuery1) deve lanciare /SearchCondition path cannot be empty/ \\
\hline
TU-209 & \ok & () => builder.build(badQuery2) deve lanciare /Invalid SearchCondition path segment/ \\
\hline
TU-210 & \ok & retrievedVector non deve essere null \\
\hline
TU-211 & \ok & retrievedVector?.getDocumentId() deve essere 1 \\
\hline
TU-212 & \ok & retrievedVector?.getEmbedding() deve essere uguale a new Float32Array([0.1, 0.2, 0.3]) \\
\hline
TU-213 & \ok & result deve essere null \\
\hline
TU-214 & \ok & result deve essere uguale a [] \\
\hline
TU-215 & \ok & dip.getUuid() deve essere "abc-123" \\
\hline
TU-216 & \ok & dip.getIntegrityStatus() deve essere IntegrityStatusEnum.UNKNOWN \\
\hline
TU-217 & \ok & dip.getId() deve essere null \\
\hline
TU-218 & \ok & dip.getIntegrityStatus() deve essere IntegrityStatusEnum.VALID \\
\hline
TU-219 & \ok & doc.getUuid() deve essere "doc-uuid" \\
\hline
TU-220 & \ok & doc.getMetadata() deve essere meta \\
\hline
TU-221 & \ok & doc.getProcessUuid() deve essere "process-uuid" \\
\hline
TU-222 & \ok & doc.getIntegrityStatus() deve essere IntegrityStatusEnum.UNKNOWN \\
\hline
TU-223 & \ok & doc.getId() deve essere null \\
\hline
TU-224 & \ok & doc.getIntegrityStatus() deve essere IntegrityStatusEnum.INVALID \\
\hline
TU-225 & \ok & dc.getDipUuid() deve essere "dip-uuid" \\
\hline
TU-226 & \ok & dc.getUuid() deve essere "dc-uuid" \\
\hline
TU-227 & \ok & dc.getName() deve essere "Contratti" \\
\hline
TU-228 & \ok & dc.getTimestamp() deve essere "2024-01-01T00:00:00Z" \\
\hline
TU-229 & \ok & dc.getIntegrityStatus() deve essere IntegrityStatusEnum.UNKNOWN \\
\hline
TU-230 & \ok & dc.getId() deve essere null \\
\hline
TU-231 & \ok & dc.getIntegrityStatus() deve essere IntegrityStatusEnum.INVALID \\
\hline
TU-232 & \ok & file.getFilename() deve essere "documento.xml" \\
\hline
TU-233 & \ok & file.getPath() deve essere "/dip/subdir/documento.xml" \\
\hline
TU-234 & \ok & file.getIsMain() deve essere true \\
\hline
TU-235 & \ok & file.getDocumentUuid() deve essere "doc-uuid" \\
\hline
TU-236 & \ok & file.getUuid() deve essere "uuid" \\
\hline
TU-237 & \ok & file.getIntegrityStatus() deve essere IntegrityStatusEnum.UNKNOWN \\
\hline
TU-238 & \ok & file.getId() deve essere null \\
\hline
TU-239 & \ok & file.getIsMain() deve essere false \\
\hline
TU-240 & \ok & file.getIntegrityStatus() deve essere IntegrityStatusEnum.INVALID \\
\hline
TU-241 & \ok & proc.getDocumentClassUuid() deve essere "dc-uuid" \\
\hline
TU-242 & \ok & proc.getUuid() deve essere "proc-uuid" \\
\hline
TU-243 & \ok & proc.getMetadata() deve essere meta \\
\hline
TU-244 & \ok & proc.getIntegrityStatus() deve essere IntegrityStatusEnum.UNKNOWN \\
\hline
TU-245 & \ok & proc.getId() deve essere null \\
\hline
TU-246 & \ok & proc.getIntegrityStatus() deve essere IntegrityStatusEnum.INVALID \\
\hline
TU-247 & \ok & resolveMock deve essere stato chiamato con DocumentoUC.GET\_BY\_ID \\
\hline
TU-248 & \ok & resolveMock deve essere stato chiamato con DocumentoUC.GET\_BY\_PROCESS \\
\hline
TU-249 & \ok & resolveMock deve essere stato chiamato con DocumentoUC.GET\_BY\_STATUS \\
\hline
TU-250 & \ok & resolveMock deve essere stato chiamato con FileUC.GET\_BY\_ID \\
\hline
TU-251 & \ok & resolveMock deve essere stato chiamato con FileUC.GET\_BY\_DOCUMENT \\
\hline
TU-252 & \ok & resolveMock deve essere stato chiamato con FileUC.GET\_BY\_STATUS \\
\hline
TU-253 & \ok & resolveMock deve essere stato chiamato con ProcessUC.GET\_BY\_ID \\
\hline
TU-254 & \ok & resolveMock deve essere stato chiamato con ProcessUC.GET\_BY\_STATUS \\
\hline
TU-255 & \ok & resolveMock deve essere stato chiamato con ProcessUC.GET\_BY\_DOCUMENT\_CLASS \\
\hline
TU-256 & \ok & resolveMock deve essere stato chiamato con DocumentClassUC.GET\_BY\_DIP\_ID \\
\hline
TU-257 & \ok & resolveMock deve essere stato chiamato con DocumentClassUC.GET\_BY\_STATUS \\
\hline
TU-258 & \ok & resolveMock deve essere stato chiamato con DocumentClassUC.GET\_BY\_ID \\
\hline
TU-259 & \ok & resolveMock deve essere stato chiamato con DipUC.GET\_BY\_ID \\
\hline
TU-260 & \ok & resolveMock deve essere stato chiamato con DipUC.GET\_BY\_STATUS \\
\hline
TU-261 & \ok & handleMock deve essere stato chiamato 15 volte \\
\hline
TU-262 & \ok & handlers.has(IpcChannels.BROWSE\_GET\_DOCUMENT\_BY\_ID) deve essere true \\
\hline
TU-263 & \ok & handlers.has(IpcChannels.BROWSE\_GET\_DOCUMENTS\_BY\_PROCESS) deve essere true \\
\hline
TU-264 & \ok & handlers.has(IpcChannels.BROWSE\_GET\_DOCUMENTS\_BY\_STATUS) deve essere true \\
\hline
TU-265 & \ok & handlers.has(IpcChannels.BROWSE\_GET\_FILE\_BY\_ID) deve essere true \\
\hline
TU-266 & \ok & handlers.has(IpcChannels.BROWSE\_GET\_FILE\_BUFFER\_BY\_ID) deve essere true \\
\hline
TU-267 & \ok & handlers.has(IpcChannels.BROWSE\_GET\_FILE\_BY\_DOCUMENT) deve essere true \\
\hline
TU-268 & \ok & handlers.has(IpcChannels.BROWSE\_GET\_FILE\_BY\_STATUS) deve essere true \\
\hline
TU-269 & \ok & handlers.has(IpcChannels.BROWSE\_GET\_PROCESS\_BY\_ID) deve essere true \\
\hline
TU-270 & \ok & handlers.has(IpcChannels.BROWSE\_GET\_PROCESS\_BY\_STATUS) deve essere true \\
\hline
TU-271 & \ok & handlers.has(IpcChannels.BROWSE\_GET\_PROCESS\_BY\_DOCUMENT\_CLASS) deve essere true \\
\hline
TU-272 & \ok & handlers.has(IpcChannels.BROWSE\_GET\_DOCUMENT\_CLASS\_BY\_DIP\_ID) deve essere true \\
\hline
TU-273 & \ok & handlers.has(IpcChannels.BROWSE\_GET\_DOCUMENT\_CLASS\_BY\_STATUS) deve essere true \\
\hline
TU-274 & \ok & handlers.has(IpcChannels.BROWSE\_GET\_DOCUMENT\_CLASS\_BY\_ID) deve essere true \\
\hline
TU-275 & \ok & handlers.has(IpcChannels.BROWSE\_GET\_DIP\_BY\_ID) deve essere true \\
\hline
TU-276 & \ok & handlers.has(IpcChannels.BROWSE\_GET\_DIP\_BY\_STATUS) deve essere true \\
\hline
TU-277 & \ok & handlers.has(IpcChannels.BROWSE\_GET\_DIP\_BY\_DOCUMENT\_CLASS) deve essere false \\
\hline
TU-278 & \ok & getDocByIdUC.execute deve essere stato chiamato con 10 \\
\hline
TU-279 & \ok & getFileByIdUC.execute deve essere stato chiamato con 11 \\
\hline
TU-280 & \ok & getProcessByIdUC.execute deve essere stato chiamato con 12 \\
\hline
TU-281 & \ok & getDocClassByIdUC.execute deve essere stato chiamato con 13 \\
\hline
TU-282 & \ok & getDipByIdUC.execute deve essere stato chiamato con 14 \\
\hline
TU-283 & \ok & documentToDtoMock deve essere stato chiamato con docEntity \\
\hline
TU-284 & \ok & fileToDtoMock deve essere stato chiamato con fileEntity \\
\hline
TU-285 & \ok & processToDtoMock deve essere stato chiamato con processEntity \\
\hline
TU-286 & \ok & documentClassToDtoMock deve essere stato chiamato con docClassEntity \\
\hline
TU-287 & \ok & dipToDtoMock deve essere stato chiamato con dipEntity \\
\hline
TU-288 & \ok & docResult deve essere uguale a { id: 1, dto: "doc" } \\
\hline
TU-289 & \ok & fileResult deve essere uguale a { id: 2, dto: "file" } \\
\hline
TU-290 & \ok & processResult deve essere uguale a { id: 3, dto: "process" } \\
\hline
TU-291 & \ok & docClassResult deve essere uguale a { id: 4, dto: "doc-class" } \\
\hline
TU-292 & \ok & dipResult deve essere uguale a { id: 5, dto: "dip" } \\
\hline
TU-293 & \ok & docResult deve essere null \\
\hline
TU-294 & \ok & fileResult deve essere null \\
\hline
TU-295 & \ok & processResult deve essere null \\
\hline
TU-296 & \ok & docClassResult deve essere null \\
\hline
TU-297 & \ok & dipResult deve essere null \\
\hline
TU-298 & \ok & documentToDtoMock non deve essere stato chiamato \\
\hline
TU-299 & \ok & fileToDtoMock non deve essere stato chiamato \\
\hline
TU-300 & \ok & processToDtoMock non deve essere stato chiamato \\
\hline
TU-301 & \ok & documentClassToDtoMock non deve essere stato chiamato \\
\hline
TU-302 & \ok & dipToDtoMock non deve essere stato chiamato \\
\hline
TU-303 & \ok & getDocByProcessUC.execute deve essere stato chiamato con 200 \\
\hline
TU-304 & \ok & getDocByStatusUC.execute deve essere stato chiamato con IntegrityStatusEnum.VALID \\
\hline
TU-305 & \ok & documentToDtoMock deve essere stato chiamato con doc1 \\
\hline
TU-306 & \ok & documentToDtoMock deve essere stato chiamato con doc2 \\
\hline
TU-307 & \ok & byProcessResult deve essere uguale a [{ mapped: doc1 }, { mapped: doc2 }] \\
\hline
TU-308 & \ok & byStatusResult deve essere uguale a [{ mapped: doc1 }] \\
\hline
TU-309 & \ok & getFileByDocUC.execute deve essere stato chiamato con 300 \\
\hline
TU-310 & \ok & getFileByStatusUC.execute deve essere stato chiamato con IntegrityStatusEnum.UNKNOWN \\
\hline
TU-311 & \ok & fileToDtoMock deve essere stato chiamato con file1 \\
\hline
TU-312 & \ok & fileToDtoMock deve essere stato chiamato con file2 \\
\hline
TU-313 & \ok & byDocumentResult deve essere uguale a [{ mapped: file1 }, { mapped: file2 }] \\
\hline
TU-314 & \ok & byStatusResult deve essere uguale a [{ mapped: file2 }] \\
\hline
TU-315 & \ok & getProcessByStatusUC.execute deve essere stato chiamato con IntegrityStatusEnum.INVALID \\
\hline
TU-316 & \ok & getProcessByDocumentClassUC.execute deve essere stato chiamato con 400 \\
\hline
TU-317 & \ok & processToDtoMock deve essere stato chiamato con process1 \\
\hline
TU-318 & \ok & processToDtoMock deve essere stato chiamato con process2 \\
\hline
TU-319 & \ok & byStatusResult deve essere uguale a [ { mapped: process1 }, { mapped: process2 }, ] \\
\hline
TU-320 & \ok & byDocClassResult deve essere uguale a [{ mapped: process2 }] \\
\hline
TU-321 & \ok & getDocClassByDipIdUC.execute deve essere stato chiamato con 500 \\
\hline
TU-322 & \ok & getDocClassByStatusUC.execute deve essere stato chiamato con IntegrityStatusEnum.VALID \\
\hline
TU-323 & \ok & documentClassToDtoMock deve essere stato chiamato con docClass1 \\
\hline
TU-324 & \ok & documentClassToDtoMock deve essere stato chiamato con docClass2 \\
\hline
TU-325 & \ok & byDipResult deve essere uguale a [{ mapped: docClass1 }, { mapped: docClass2 }] \\
\hline
TU-326 & \ok & byStatusResult deve essere uguale a [{ mapped: docClass2 }] \\
\hline
TU-327 & \ok & getDipByStatusUC.execute deve essere stato chiamato con IntegrityStatusEnum.UNKNOWN \\
\hline
TU-328 & \ok & dipToDtoMock deve essere stato chiamato con dip1 \\
\hline
TU-329 & \ok & dipToDtoMock deve essere stato chiamato con dip2 \\
\hline
TU-330 & \ok & result deve essere uguale a [{ mapped: dip1 }, { mapped: dip2 }] \\
\hline
TU-331 & \ok & resolveMock deve essere stato chiamato con DocumentoUC.CHECK\_INTEGRITY\_STATUS \\
\hline
TU-332 & \ok & resolveMock deve essere stato chiamato con FileUC.CHECK\_INTEGRITY\_STATUS \\
\hline
TU-333 & \ok & resolveMock deve essere stato chiamato con ProcessUC.CHECK\_INTEGRITY\_STATUS \\
\hline
TU-334 & \ok & resolveMock deve essere stato chiamato con DocumentClassUC.CHECK\_INTEGRITY\_STATUS \\
\hline
TU-335 & \ok & resolveMock deve essere stato chiamato con DipUC.CHECK\_INTEGRITY\_STATUS \\
\hline
TU-336 & \ok & handleMock deve essere stato chiamato 5 volte \\
\hline
TU-337 & \ok & handleMock.mock.calls.some( (call) => call[0] === IpcChannels.CHECK\_DOCUMENT\_INTEGRITY\_STATUS, ) deve essere true \\
\hline
TU-338 & \ok & handleMock.mock.calls.some( (call) => call[0] === IpcChannels.CHECK\_FILE\_INTEGRITY\_STATUS, ) deve essere true \\
\hline
TU-339 & \ok & handleMock.mock.calls.some( (call) => call[0] === IpcChannels.CHECK\_PROCESS\_INTEGRITY\_STATUS, ) deve essere true \\
\hline
TU-340 & \ok & handleMock.mock.calls.some( (call) => call[0] === IpcChannels.CHECK\_DOCUMENT\_CLASS\_INTEGRITY\_STATUS, ) deve essere true \\
\hline
TU-341 & \ok & handleMock.mock.calls.some( (call) => call[0] === IpcChannels.CHECK\_DIP\_INTEGRITY\_STATUS, ) deve essere true \\
\hline
TU-342 & \ok & checkDocumentIntegrityUC.execute deve essere stato chiamato con 10 \\
\hline
TU-343 & \ok & checkProcessIntegrityUC.execute deve essere stato chiamato con 11 \\
\hline
TU-344 & \ok & checkDocumentClassIntegrityUC.execute deve essere stato chiamato con 12 \\
\hline
TU-345 & \ok & checkDipIntegrityUC.execute deve essere stato chiamato con 13 \\
\hline
TU-346 & \ok & documentResult deve essere IntegrityStatusEnum.VALID \\
\hline
TU-347 & \ok & processResult deve essere IntegrityStatusEnum.INVALID \\
\hline
TU-348 & \ok & documentClassResult deve essere IntegrityStatusEnum.UNKNOWN \\
\hline
TU-349 & \ok & dipResult deve essere IntegrityStatusEnum.VALID \\
\hline
TU-350 & \ok & checkFileIntegrityUC.execute deve essere stato chiamato con 22 \\
\hline
TU-351 & \ok & result deve essere IntegrityStatusEnum.UNKNOWN \\
\hline
TU-352 & \ok & getRegisteredHandler(handleMock, IpcChannels.CHECK\_FILE\_INTEGRITY\_STATUS)( {}, 99, ) deve lanciare "File with id 99 not found" \\
\hline
TU-353 & \ok & registeredChannels deve contenere IpcChannels.FILE\_DOWNLOAD \\
\hline
TU-354 & \ok & exportFileUC.execute deve essere stato chiamato con 1 \\
\hline
TU-355 & \ok & result.success deve essere true \\
\hline
TU-356 & \ok & result.success deve essere false \\
\hline
TU-357 & \ok & result.errorCode deve essere "NOT\_FOUND" \\
\hline
TU-358 & \ok & result.success deve essere false \\
\hline
TU-359 & \ok & result.errorCode deve essere "WRITE\_ERROR" \\
\hline
TU-360 & \ok & ipcMain.invoke(IpcChannels.FILE\_DOWNLOAD, 1) deve lanciare "export exploded" \\
\hline
TU-361 & \ok & registeredChannels deve contenere IpcChannels.SEARCH\_CLASSES \\
\hline
TU-362 & \ok & registeredChannels deve contenere IpcChannels.SEARCH\_PROCESSES \\
\hline
TU-363 & \ok & registeredChannels deve contenere IpcChannels.SEARCH\_DOCUMENTS \\
\hline
TU-364 & \ok & registeredChannels deve contenere IpcChannels.SEARCH\_SEMANTIC \\
\hline
TU-365 & \ok & registeredChannels deve contenere IpcChannels.SEARCH\_FULLTEXT \\
\hline
TU-366 & \ok & registeredChannels deve contenere IpcChannels.SEARCH\_GET\_AI\_STATE \\
\hline
TU-367 & \ok & registeredChannels deve contenere IpcChannels.SEARCH\_CUSTOM\_METADATA\_KEYS \\
\hline
TU-368 & \ok & searchClassesUC.execute deve essere stato chiamato con "Contratti" \\
\hline
TU-369 & \ok & result deve essere uguale a [DocumentClassMapper.toSearchResult(dc)] \\
\hline
TU-370 & \ok & searchClassesUC.execute deve essere stato chiamato con "" \\
\hline
TU-371 & \ok & result deve essere uguale a [] \\
\hline
TU-372 & \ok & ipcMain.invoke(IpcChannels.SEARCH\_CLASSES, "Contratti") deve lanciare "classes failed" \\
\hline
TU-373 & \ok & searchProcessiUC.execute deve essere stato chiamato con "proc-uuid-abc" \\
\hline
TU-374 & \ok & result deve essere uguale a [ProcessMapper.toSearchResult(proc)] \\
\hline
TU-375 & \ok & searchProcessiUC.execute deve essere stato chiamato con "" \\
\hline
TU-376 & \ok & result deve essere uguale a [] \\
\hline
TU-377 & \ok & ipcMain.invoke(IpcChannels.SEARCH\_PROCESSES, "uuid-1") deve lanciare "processes failed" \\
\hline
TU-378 & \ok & searchDocumentsUC.execute deve essere stato chiamato \\
\hline
TU-379 & \ok & result deve essere uguale a [DocumentMapper.toSearchResult(doc, null)] \\
\hline
TU-380 & \ok & result deve essere uguale a [] \\
\hline
TU-381 & \ok & ipcMain.invoke(IpcChannels.SEARCH\_DOCUMENTS, filters) deve lanciare "documents failed" \\
\hline
TU-382 & \ok & searchSemanticUC.execute deve essere stato chiamato con "ricerca semantica" \\
\hline
TU-383 & \ok & result deve essere uguale a [DocumentMapper.toSearchResult(doc, 0.91)] \\
\hline
TU-384 & \ok & result deve essere uguale a [] \\
\hline
TU-385 & \ok & searchSemanticUC.execute deve essere stato chiamato con "" \\
\hline
TU-386 & \ok & ipcMain.invoke(IpcChannels.SEARCH\_SEMANTIC, query) deve lanciare "semantic failed" \\
\hline
TU-387 & \ok & calledQuery.filter.items[0] deve essere uguale a { path: "NomeDelDocumento", operator: "LIKE", value: "\%fattura.pdf\%", } \\
\hline
TU-388 & \ok & result deve essere uguale a expectedResults.map((doc) => DocumentMapper.toSearchResult(doc, null)) \\
\hline
TU-389 & \ok & result deve essere uguale a [] \\
\hline
TU-390 & \ok & ipcMain.invoke(IpcChannels.SEARCH\_FULLTEXT, query) deve lanciare "fulltext failed" \\
\hline
TU-391 & \ok & result.status deve essere "IDLE" \\
\hline
TU-392 & \ok & result.progressPercentage deve essere 0 \\
\hline
TU-393 & \ok & result.lastIndexedAt deve essere null \\
\hline
TU-394 & \ok & result.indexedDocuments deve essere 0 \\
\hline
TU-395 & \ok & result.status deve essere "READY" \\
\hline
TU-396 & \ok & result.progressPercentage deve essere 100 \\
\hline
TU-397 & \ok & result.lastIndexedAt deve essere null \\
\hline
TU-398 & \ok & result.indexedDocuments deve essere 42 \\
\hline
TU-399 & \ok & result.indexedDocuments deve essere 150 \\
\hline
TU-400 & \ok & ipcMain.invoke(IpcChannels.SEARCH\_GET\_AI\_STATE) deve lanciare "ai state unavailable" \\
\hline
TU-401 & \ok & getCustomMetadataKeysUC.execute deve essere stato chiamato con 7 \\
\hline
TU-402 & \ok & result deve essere uguale a ["DataProtocollo", "NomeCliente", "NumeroPratica"] \\
\hline
TU-403 & \ok & getCustomMetadataKeysUC.execute deve essere stato chiamato con null \\
\hline
TU-404 & \ok & result deve essere uguale a ["NomeCliente"] \\
\hline
TU-405 & \ok & dao.save deve essere stato chiamato con input \\
\hline
TU-406 & \ok & result deve essere savedDip \\
\hline
TU-407 & \ok & dao.getById deve essere stato chiamato con 12 \\
\hline
TU-408 & \ok & dao.getByUuid deve essere stato chiamato con "dip-2" \\
\hline
TU-409 & \ok & byId?.getUuid() deve essere "dip-2" \\
\hline
TU-410 & \ok & byUuid?.getId() deve essere 12 \\
\hline
TU-411 & \ok & dao.updateIntegrityStatus deve essere stato chiamato con 13, IntegrityStatusEnum.VALID \\
\hline
TU-412 & \ok & dao.getByStatus deve essere stato chiamato con IntegrityStatusEnum.VALID \\
\hline
TU-413 & \ok & found deve avere lunghezza 1 \\
\hline
TU-414 & \ok & found[0].getUuid() deve essere "dip-3" \\
\hline
TU-415 & \ok & dao.getByStatus deve essere stato chiamato con IntegrityStatusEnum.VALID \\
\hline
TU-416 & \ok & found deve avere lunghezza 0 \\
\hline
TU-417 & \ok & dao.save deve essere stato chiamato con dip \\
\hline
TU-418 & \ok & saved.getId() deve essere 99 \\
\hline
TU-419 & \ok & dao.save deve essere stato chiamato con dip \\
\hline
TU-420 & \ok & saved.getId() deve essere 88 \\
\hline
TU-421 & \ok & saved.getUuid() deve essere "dip-update" \\
\hline
TU-422 & \ok & dao.save deve essere stato chiamato con input \\
\hline
TU-423 & \ok & dao.getById deve essere stato chiamato con 21 \\
\hline
TU-424 & \ok & found?.getDipId() deve essere 10 \\
\hline
TU-425 & \ok & found?.getUuid() deve essere "dc-1" \\
\hline
TU-426 & \ok & repo.getByDipId(20) deve avere lunghezza 1 \\
\hline
TU-427 & \ok & dao.updateIntegrityStatus deve essere stato chiamato con 31, IntegrityStatusEnum.VALID \\
\hline
TU-428 & \ok & byStatus deve avere lunghezza 1 \\
\hline
TU-429 & \ok & dao.search deve essere stato chiamato con "Verbali" \\
\hline
TU-430 & \ok & dao.search deve essere stato chiamato con "inesistente" \\
\hline
TU-431 & \ok & found non deve essere null \\
\hline
TU-432 & \ok & found?.[0].getName() deve contenere "Verbali" \\
\hline
TU-433 & \ok & notFound deve essere null \\
\hline
TU-434 & \ok & () => repo.save(input).toThrowError("Failed to save DocumentClass with uuid=dc-err") \\
\hline
TU-435 & \ok & dao.getById deve essere stato chiamato con 404 \\
\hline
TU-436 & \ok & result deve essere null \\
\hline
TU-437 & \ok & dao.save deve essere stato chiamato con input \\
\hline
TU-438 & \ok & result deve essere saved \\
\hline
TU-439 & \ok & result.getUuid() deve essere "doc-unique" \\
\hline
TU-440 & \ok & dao.save deve essere stato chiamato 1 volte \\
\hline
TU-441 & \ok & dao.getByProcessId deve essere stato chiamato con 1 \\
\hline
TU-442 & \ok & results deve avere lunghezza 1 \\
\hline
TU-443 & \ok & results[0].getUuid() deve essere "doc-p1" \\
\hline
TU-444 & \ok & results deve avere lunghezza 0 \\
\hline
TU-445 & \ok & dao.getByStatus deve essere stato chiamato con IntegrityStatusEnum.VALID \\
\hline
TU-446 & \ok & results deve avere lunghezza 1 \\
\hline
TU-447 & \ok & repo.getByStatus(IntegrityStatusEnum.VALID) deve avere lunghezza 0 \\
\hline
TU-448 & \ok & dao.updateIntegrityStatus deve essere stato chiamato con 7, IntegrityStatusEnum.VALID \\
\hline
TU-449 & \ok & result.getId() deve essere 33 \\
\hline
TU-450 & \ok & dao.getById deve essere stato chiamato con 999 \\
\hline
TU-451 & \ok & result deve essere null \\
\hline
TU-452 & \ok & dao.searchDocument deve essere stato chiamato con emptyFilters \\
\hline
TU-453 & \ok & results deve essere uguale a [] \\
\hline
TU-454 & \ok & dao.searchDocumentSemantic deve essere stato chiamato con queryVector \\
\hline
TU-455 & \ok & result deve essere uguale a semantic \\
\hline
TU-456 & \ok & dao.getIndexedDocumentsCount deve essere stato chiamato 1 volte \\
\hline
TU-457 & \ok & count deve essere 42 \\
\hline
TU-458 & \ok & dto.id deve essere 1 \\
\hline
TU-459 & \ok & dto.uuid deve essere "test-uuid" \\
\hline
TU-460 & \ok & dto.integrityStatus deve essere IntegrityStatusEnum.VALID \\
\hline
TU-461 & \ok & () => EntityToDtoConverter.dipToDto(dip) deve lanciare "Cannot convert Dip to DTO: id is null" \\
\hline
TU-462 & \ok & dto.integrityStatus deve essere IntegrityStatusEnum.UNKNOWN \\
\hline
TU-463 & \ok & dto.id deve essere 2 \\
\hline
TU-464 & \ok & dto.dipId deve essere 1 \\
\hline
TU-465 & \ok & dto.uuid deve essere "dc-uuid" \\
\hline
TU-466 & \ok & dto.name deve essere "Ricevute" \\
\hline
TU-467 & \ok & dto.timestamp deve essere "2025-03-01T08:00:00Z" \\
\hline
TU-468 & \ok & dto.integrityStatus deve essere IntegrityStatusEnum.VALID \\
\hline
TU-469 & \ok & () => EntityToDtoConverter.documentClassToDto(dc) deve lanciare "Cannot convert DocumentClass to DTO: id or dipId is null" \\
\hline
TU-470 & \ok & dto.integrityStatus deve essere IntegrityStatusEnum.INVALID \\
\hline
TU-471 & \ok & dto.id deve essere 5 \\
\hline
TU-472 & \ok & dto.processId deve essere 10 \\
\hline
TU-473 & \ok & dto.uuid deve essere "doc-uuid" \\
\hline
TU-474 & \ok & dto.integrityStatus deve essere IntegrityStatusEnum.VALID \\
\hline
TU-475 & \ok & dto.metadata deve essere definito \\
\hline
TU-476 & \ok & dto.metadata.name deve essere "DocumentoInformatico" \\
\hline
TU-477 & \ok & dto.metadata.value deve avere lunghezza 2 \\
\hline
TU-478 & \ok & () => EntityToDtoConverter.documentToDto(doc) deve lanciare "Cannot convert Document to DTO: id or processId is null" \\
\hline
TU-479 & \ok & dto.metadata.name deve essere "root" \\
\hline
TU-480 & \ok & Array.isArray(dto.metadata.value) deve essere true \\
\hline
TU-481 & \ok & (dto.metadata.value as MetadataDTO[]).length deve essere 0 \\
\hline
TU-482 & \ok & dto.id deve essere 7 \\
\hline
TU-483 & \ok & dto.documentId deve essere 5 \\
\hline
TU-484 & \ok & dto.filename deve essere "documento.pdf" \\
\hline
TU-485 & \ok & dto.path deve essere "/docs/documento.pdf" \\
\hline
TU-486 & \ok & dto.hash deve essere "abc123hash" \\
\hline
TU-487 & \ok & dto.integrityStatus deve essere IntegrityStatusEnum.VALID \\
\hline
TU-488 & \ok & dto.isMain deve essere true \\
\hline
TU-489 & \ok & () => EntityToDtoConverter.fileToDto(file) deve lanciare "Cannot convert File to DTO: id or documentId is null" \\
\hline
TU-490 & \ok & dto.isMain deve essere false \\
\hline
TU-491 & \ok & dto.integrityStatus deve essere IntegrityStatusEnum.INVALID \\
\hline
TU-492 & \ok & dto.id deve essere 10 \\
\hline
TU-493 & \ok & dto.documentClassId deve essere 2 \\
\hline
TU-494 & \ok & dto.uuid deve essere "proc-uuid" \\
\hline
TU-495 & \ok & dto.integrityStatus deve essere IntegrityStatusEnum.VALID \\
\hline
TU-496 & \ok & (dto.metadata.value as MetadataDTO[]).length deve essere 2 \\
\hline
TU-497 & \ok & (dto.metadata.value[0] as MetadataDTO).name deve essere "fase" \\
\hline
TU-498 & \ok & (dto.metadata.value[0] as MetadataDTO).value deve essere "Acquisizione" \\
\hline
TU-499 & \ok & (dto.metadata.value[1] as MetadataDTO).name deve essere "ordine" \\
\hline
TU-500 & \ok & (dto.metadata.value[1] as MetadataDTO).value deve essere "1" \\
\hline
TU-501 & \ok & () => EntityToDtoConverter.processToDto(proc) deve lanciare "Cannot convert Process to DTO: id or documentClassId is null" \\
\hline
TU-502 & \ok & dto.metadata.name deve essere "ProcessoVuoto" \\
\hline
TU-503 & \ok & Array.isArray(dto.metadata.value) deve essere true \\
\hline
TU-504 & \ok & (dto.metadata.value as MetadataDTO[]).length deve essere 0 \\
\hline
TU-505 & \ok & dto.name deve essere "autore" \\
\hline
TU-506 & \ok & dto.value deve essere "Mario Rossi" \\
\hline
TU-507 & \ok & dto.type deve essere MetadataType.STRING \\
\hline
TU-508 & \ok & dto.name deve essere "anno" \\
\hline
TU-509 & \ok & dto.value deve essere "2025" \\
\hline
TU-510 & \ok & dto.type deve essere MetadataType.NUMBER \\
\hline
TU-511 & \ok & dto.name deve essere "riservato" \\
\hline
TU-512 & \ok & dto.value deve essere "true" \\
\hline
TU-513 & \ok & dto.type deve essere MetadataType.BOOLEAN \\
\hline
TU-514 & \ok & dto.name deve essere "root" \\
\hline
TU-515 & \ok & dto.type deve essere MetadataType.COMPOSITE \\
\hline
TU-516 & \ok & Array.isArray(dto.value) deve essere true \\
\hline
TU-517 & \ok & (dto.value as any[]).length deve essere 2 \\
\hline
TU-518 & \ok & dto.name deve essere "emptyRoot" \\
\hline
TU-519 & \ok & dto.type deve essere MetadataType.COMPOSITE \\
\hline
TU-520 & \ok & Array.isArray(dto.value) deve essere true \\
\hline
TU-521 & \ok & (dto.value as MetadataDTO[]).length deve essere 0 \\
\hline
TU-522 & \ok & dto.name deve essere "level1" \\
\hline
TU-523 & \ok & level2.name deve essere "level2" \\
\hline
TU-524 & \ok & level3.name deve essere "level3" \\
\hline
TU-525 & \ok & deep.name deve essere "deep" \\
\hline
TU-526 & \ok & deep.value deve essere "value" \\
\hline
TU-527 & \ok & dao.save deve essere stato chiamato con input \\
\hline
TU-528 & \ok & dao.getById deve essere stato chiamato con 71 \\
\hline
TU-529 & \ok & found?.getFilename() deve essere "main.xml" \\
\hline
TU-530 & \ok & found?.getDocumentId() deve essere 3 \\
\hline
TU-531 & \ok & repo.getByDocumentId(4) deve avere lunghezza 1 \\
\hline
TU-532 & \ok & dao.updateIntegrityStatus deve essere stato chiamato con 72, IntegrityStatusEnum.INVALID \\
\hline
TU-533 & \ok & dao.getByStatus deve essere stato chiamato con IntegrityStatusEnum.INVALID \\
\hline
TU-534 & \ok & byStatus deve avere lunghezza 1 \\
\hline
TU-535 & \ok & dao.save deve essere stato chiamato con input \\
\hline
TU-536 & \ok & result.getId() deve essere 88 \\
\hline
TU-537 & \ok & dao.getById deve essere stato chiamato con 404 \\
\hline
TU-538 & \ok & result deve essere null \\
\hline
TU-539 & \ok & result.success deve essere true \\
\hline
TU-540 & \ok & fileContent deve essere data \\
\hline
TU-541 & \ok & result.success deve essere false \\
\hline
TU-542 & \ok & result.errorCode deve essere "WRITE\_ERROR" \\
\hline
TU-543 & \ok & adapter.setDipPath(dipPath) deve lanciare "DiP index file not found in '/mock/dip/path'. Expected format: DiPIndex.<uuid>.xml" \\
\hline
TU-544 & \ok & result deve essere un'istanza di Dip \\
\hline
TU-545 & \ok & result.getUuid() deve essere "dip-from-a" \\
\hline
TU-546 & \ok & results deve avere lunghezza 2 \\
\hline
TU-547 & \ok & results[0] deve essere un'istanza di DocumentClass \\
\hline
TU-548 & \ok & results[0].getDipUuid() deve essere "dip-uuid" \\
\hline
TU-549 & \ok & results[0].getUuid() deve essere "dc-1" \\
\hline
TU-550 & \ok & results[0].getName() deve essere "doc1" \\
\hline
TU-551 & \ok & results[1] deve essere un'istanza di DocumentClass \\
\hline
TU-552 & \ok & results[1].getDipUuid() deve essere "dip-uuid" \\
\hline
TU-553 & \ok & results[1].getUuid() deve essere "dc-2" \\
\hline
TU-554 & \ok & results[1].getName() deve essere "doc2" \\
\hline
TU-555 & \ok & results deve avere lunghezza 1 \\
\hline
TU-556 & \ok & results[0] deve essere un'istanza di Process \\
\hline
TU-557 & \ok & results[0].getDocumentClassUuid() deve essere "dc-uuid" \\
\hline
TU-558 & \ok & results[0].getUuid() deve essere "proc-meta" \\
\hline
TU-559 & \ok & results[0].getMetadata() deve essere uguale a new Metadata("meta", "value") \\
\hline
TU-560 & \ok & results deve avere lunghezza 1 \\
\hline
TU-561 & \ok & results[0] deve essere un'istanza di Process \\
\hline
TU-562 & \ok & results[0].getDocumentClassUuid() deve essere "dc-uuid" \\
\hline
TU-563 & \ok & results[0].getUuid() deve essere "proc-null-metadata" \\
\hline
TU-564 & \ok & results[0].getMetadata() deve essere uguale a new Metadata("meta", "value") \\
\hline
TU-565 & \ok & results deve avere lunghezza 1 \\
\hline
TU-566 & \ok & results[0] deve essere un'istanza di Document \\
\hline
TU-567 & \ok & results[0].getUuid() deve essere "doc-without-metadata" \\
\hline
TU-568 & \ok & results[0].getProcessUuid() deve essere "proc-uuid" \\
\hline
TU-569 & \ok & results[0].getMetadata() deve essere uguale a new Metadata("meta", "value") \\
\hline
TU-570 & \ok & results deve avere lunghezza 1 \\
\hline
TU-571 & \ok & results[0] deve essere un'istanza di File \\
\hline
TU-572 & \ok & results[0].getFilename() deve essere "some-file" \\
\hline
TU-573 & \ok & results[0].getPath() deve essere "path/without-meta" \\
\hline
TU-574 & \ok & results[0].getHash() deve essere "hash" \\
\hline
TU-575 & \ok & results[0].getIsMain() deve essere true \\
\hline
TU-576 & \ok & results[0].getDocumentUuid() deve essere "doc-uuid" \\
\hline
TU-577 & \ok & BrowserWindow deve essere stato chiamato con { show: false, webPreferences: { plugins: true }, } \\
\hline
TU-578 & \ok & mockLoadURL deve essere stato chiamato con "file:///tmp/doc.pdf" \\
\hline
TU-579 & \ok & mockPrint deve essere stato chiamato con opts, any Function \\
\hline
TU-580 & \ok & result deve essere uguale a { success: true } \\
\hline
TU-581 & \ok & mockDestroy.toHaveBeenCalledOnce() \\
\hline
TU-582 & \ok & result.success deve essere false \\
\hline
TU-583 & \ok & result.errorMessage deve essere "printer not found" \\
\hline
TU-584 & \ok & result.errorCode deve essere "PRINT\_ERROR" \\
\hline
TU-585 & \ok & mockDestroy.toHaveBeenCalledOnce() \\
\hline
TU-586 & \ok & result.success deve essere false \\
\hline
TU-587 & \ok & result.errorMessage deve essere "did-fail-load: ERR\_FILE\_NOT\_FOUND (codice: -6)" \\
\hline
TU-588 & \ok & result.errorCode deve essere "LOAD\_ERROR" \\
\hline
TU-589 & \ok & mockDestroy.toHaveBeenCalledOnce() \\
\hline
TU-590 & \ok & mockPrint non deve essere stato chiamato \\
\hline
TU-591 & \ok & dao.save deve essere stato chiamato con input \\
\hline
TU-592 & \ok & result.getId() deve essere 81 \\
\hline
TU-593 & \ok & result.getDocumentClassId() deve essere 11 \\
\hline
TU-594 & \ok & repo.getByDocumentClassId(12) deve avere lunghezza 1 \\
\hline
TU-595 & \ok & dao.updateIntegrityStatus deve essere stato chiamato con 82, IntegrityStatusEnum.VALID \\
\hline
TU-596 & \ok & dao.getByStatus deve essere stato chiamato con IntegrityStatusEnum.VALID \\
\hline
TU-597 & \ok & byStatus deve avere lunghezza 1 \\
\hline
TU-598 & \ok & byStatus[0].getUuid() deve essere "proc-2" \\
\hline
TU-599 & \ok & dao.searchDocumentSemantic deve essere stato chiamato con queryVector \\
\hline
TU-600 & \ok & dao.searchDocumentSemantic deve essere stato chiamato 1 volte \\
\hline
TU-601 & \ok & result deve essere semanticResults \\
\hline
TU-602 & \ok & result deve avere lunghezza 2 \\
\hline
TU-603 & \ok & result[0].score deve essere 0.91 \\
\hline
TU-604 & \ok & result[1].document.getUuid() deve essere "doc-2" \\
\hline
TU-605 & \ok & repo.searchDocumentSemantic(new Float32Array([0.7, 0.8, 0.9])) deve lanciare "semantic failed" \\
\hline
TU-606 & \ok & first deve avere lunghezza 1 \\
\hline
TU-607 & \ok & second deve avere lunghezza 0 \\
\hline
TU-608 & \ok & dao.searchDocumentSemantic deve essere stato chiamato 2 volte \\
\hline
TU-609 & \ok & results[0].score deve essere 1 \\
\hline
TU-610 & \ok & results[0].score deve essere 0 \\
\hline
TU-611 & \ok & scores[i] deve essere maggiore o uguale a scores[i + 1] \\
\hline
TU-612 & \ok & results[0].score deve essere vicino a 0.123456789, 9 \\
\hline
TU-613 & \ok & results[0].score deve essere maggiore di 0.8 \\
\hline
TU-614 & \ok & results[0].score deve essere minore di 0.5 \\
\hline
TU-615 & \ok & results.length deve essere minore o uguale a 10 \\
\hline
TU-616 & \ok & r.score deve essere maggiore o uguale a 0 \\
\hline
TU-617 & \ok & r.score deve essere minore o uguale a 1 \\
\hline
TU-618 & \ok & uc.execute("missing doc") deve lanciare "999" \\
\hline
TU-619 & \ok & aiAdapter.generateEmbedding deve essere stato chiamato con "test vettore" \\
\hline
TU-620 & \ok & vectorRepo.searchSimilarVectors deve essere stato chiamato con vector, 10 \\
\hline
TU-621 & \ok & result deve essere "ok" \\
\hline
TU-622 & \ok & db.exec.mock.calls deve essere uguale a [ ["BEGIN IMMEDIATE TRANSACTION"], ["COMMIT"], ] \\
\hline
TU-623 & \ok & manager.runInTransaction(async () => { throw error}) deve essere error \\
\hline
TU-625 & \ok & db.exec.mock.calls deve essere uguale a [ ["BEGIN IMMEDIATE TRANSACTION"], ["ROLLBACK"], ] \\
\hline
TU-626 & \ok & result deve essere payload \\
\hline
TU-627 & \ok & result.getUuid() deve essere "test-dip-uuid-1234" \\
\hline
TU-628 & \ok & result[0].getUuid() deve essere "class-1" \\
\hline
TU-629 & \ok & result[0].getName() deve essere "Class1" \\
\hline
TU-630 & \ok & result[0].metadataRelativePath deve essere "./class-1/aip-1/AiPInfo.aip-1.xml" \\
\hline
TU-631 & \ok & result[1].metadataRelativePath deve essere "./class-2/aip-2/AiPInfo.aip-2.xml" \\
\hline
TU-632 & \ok & proc.getUuid() deve essere "aip-1" \\
\hline
TU-633 & \ok & proc.getDocumentClassUuid() deve essere "class-1" \\
\hline
TU-634 & \ok & processMeta deve essere definito \\
\hline
TU-635 & \ok & startMeta deve essere definito \\
\hline
TU-636 & \ok & startMeta.getChildren() deve essere uguale a [ expect.objectContaining({ name: "Date", value: "2025-01-01T00:00:00Z", }), expect.objectContaining({ name: "UserUUID", value: "dummy-user-uuid" }), expect.objectContaining({ name: "Source", value: "DummySource" }), ] \\
\hline
TU-637 & \ok & endMeta deve essere definito \\
\hline
TU-638 & \ok & endMeta.getChildren() deve essere uguale a [ expect.objectContaining({ name: "Date", value: "2025-01-01T01:00:00Z", }), expect.objectContaining({ name: "UserUUID", value: "dummy-user-uuid" }), expect.objectContaining({ name: "Source", value: "DummySource" }), expect.objectContaining({ name: "Status", value: "Completed" }), ] \\
\hline
TU-639 & \ok & submissionMeta deve essere definito \\
\hline
TU-640 & \ok & submissionStart?.getChildren() deve essere uguale a [ expect.objectContaining({ name: "Date", value: "2025-01-01T00:00:00Z", }), expect.objectContaining({ name: "UserUUID", value: "dummy-user-uuid" }), expect.objectContaining({ name: "Source", value: "DummySource" }), ] \\
\hline
TU-641 & \ok & submissionEnd?.getChildren() deve essere uguale a [ expect.objectContaining({ name: "Date", value: "2025-01-01T00:30:00Z", }), expect.objectContaining({ name: "UserUUID", value: "dummy-user-uuid" }), expect.objectContaining({ name: "Source", value: "DummySource" }), expect.objectContaining({ name: "Status", value: "Completed" }), ] \\
\hline
TU-642 & \ok & preservationMeta deve essere definito \\
\hline
TU-643 & \ok & preservationStart?.getChildren() deve essere uguale a [ expect.objectContaining({ name: "Date", value: "2025-01-01T00:45:00Z", }), expect.objectContaining({ name: "UserUUID", value: "dummy-user-uuid" }), expect.objectContaining({ name: "Source", value: "DummySource" }), ] \\
\hline
TU-644 & \ok & preservationEnd.getChildren() deve essere definito \\
\hline
TU-645 & \ok & proc.getUuid() deve essere "aip-1" \\
\hline
TU-646 & \ok & proc.getDocumentClassUuid() deve essere "class-1" \\
\hline
TU-647 & \ok & proc.getMetadata() deve essere uguale a new Metadata("AiPInfo", [], MetadataType.COMPOSITE) \\
\hline
TU-648 & \ok & proc.getUuid() deve essere "aip-1" \\
\hline
TU-649 & \ok & proc.getDocumentClassUuid() deve essere "class-1" \\
\hline
TU-650 & \ok & proc.getMetadata() deve essere uguale a new Metadata("AiPInfo", [], MetadataType.COMPOSITE) \\
\hline
TU-651 & \ok & doc.getUuid() deve essere "doc-1" \\
\hline
TU-652 & \ok & doc.getProcessUuid() deve essere "aip-1" \\
\hline
TU-653 & \ok & doc.getIntegrityStatus() deve essere "UNKNOWN" \\
\hline
TU-654 & \ok & metadata.getName() deve essere "DocumentoInformatico" \\
\hline
TU-655 & \ok & metadata.getChildren() deve avere lunghezza 6 \\
\hline
TU-656 & \ok & metadata.getChildren()[0].getName() deve essere "IdDoc" \\
\hline
TU-657 & \ok & metadata.getChildren()[0].getChildren() deve avere lunghezza 2 \\
\hline
TU-658 & \ok & metadata.getChildren()[0].getChildren()[0].getName() deve essere "ImprontaCrittograficaDelDocumento" \\
\hline
TU-659 & \ok & metadata.getChildren()[0].getChildren()[0].getChildren() deve avere lunghezza 2 \\
\hline
TU-660 & \ok & metadata.getChildren()[0].getChildren()[0].getChildren()[0].getName() deve essere "Impronta" \\
\hline
TU-661 & \ok & metadata.getChildren()[0].getChildren()[0].getChildren()[1].getName() deve essere "Algoritmo" \\
\hline
TU-662 & \ok & metadata.getChildren()[0].getChildren()[1].getName() deve essere "Identificativo" \\
\hline
TU-663 & \ok & metadata.getChildren()[1].getName() deve essere "Riservato" \\
\hline
TU-664 & \ok & metadata.getChildren()[2].getName() deve essere "TempoDiConservazione" \\
\hline
TU-665 & \ok & metadata.getChildren()[3].getName() deve essere "Soggetti" \\
\hline
TU-666 & \ok & metadata.getChildren()[3].getChildren() deve avere lunghezza 2 \\
\hline
TU-667 & \ok & metadata.getChildren()[3].getChildren()[0].getName() deve essere "Ruolo" \\
\hline
TU-668 & \ok & metadata.getChildren()[3].getChildren()[0].getChildren()[0] deve essere uguale a new Metadata("Tipo", "Autore", MetadataType.STRING) \\
\hline
TU-669 & \ok & metadata.getChildren()[4].getName() deve essere "Titolo" \\
\hline
TU-670 & \ok & metadata.getChildren()[4] deve essere uguale a new Metadata("Titolo", "Documento di test", MetadataType.STRING) \\
\hline
TU-671 & \ok & result[0].metadataRelativePath deve essere "./class-1/aip-1/documents/doc-1/meta.xml" \\
\hline
TU-672 & \ok & result[1].metadataRelativePath deve essere "./class-2/aip-2/documents/doc-2/meta2.xml" \\
\hline
TU-673 & \ok & doc.getUuid() deve essere "doc-1" \\
\hline
TU-674 & \ok & doc.getProcessUuid() deve essere "aip-1" \\
\hline
TU-675 & \ok & doc.getIntegrityStatus() deve essere "UNKNOWN" \\
\hline
TU-676 & \ok & doc.getMetadata().getName() deve essere "Unknown" \\
\hline
TU-677 & \ok & doc.getMetadata().getChildren() deve avere lunghezza 0 \\
\hline
TU-678 & \ok & result deve avere lunghezza 5 \\
\hline
TU-679 & \ok & metadataPaths.filter( (p) => p === "./class-1/aip-1/documents/doc-1/meta.xml", ) deve avere lunghezza 1 \\
\hline
TU-680 & \ok & metadataPaths.filter( (p) => p === "./class-2/aip-2/documents/doc-2/meta2.xml", ) deve avere lunghezza 3 \\
\hline
TU-681 & \ok & metadataPaths.filter( (p) => p === "./class-2/aip-3/documents/doc-3/meta3.xml", ) deve avere lunghezza 1 \\
\hline
TU-682 & \ok & mappedFiles deve avere lunghezza 5 \\
\hline
TU-683 & \ok & mappedFiles[0].getUuid() deve essere "prim-1" \\
\hline
TU-684 & \ok & mappedFiles[0].getFilename() deve essere "./primary.pdf" \\
\hline
TU-685 & \ok & mappedFiles[0].getPath() deve essere "class-1/aip-1/documents/doc-1/primary.pdf" \\
\hline
TU-686 & \ok & mappedFiles[0].getHash() deve essere "hash-main-prim-1" \\
\hline
TU-687 & \ok & mappedFiles[0].getIsMain() deve essere true \\
\hline
TU-688 & \ok & mappedFiles[0].getDocumentId() deve essere null \\
\hline
TU-689 & \ok & mappedFiles[1].getUuid() deve essere "prim-2" \\
\hline
TU-690 & \ok & mappedFiles[1].getHash() deve essere "" \\
\hline
TU-691 & \ok & mappedFiles[1].getIsMain() deve essere true \\
\hline
TU-692 & \ok & mappedFiles[2].getUuid() deve essere "prim-3" \\
\hline
TU-693 & \ok & mappedFiles[2].getHash() deve essere "" \\
\hline
TU-694 & \ok & mappedFiles[2].getIsMain() deve essere true \\
\hline
TU-695 & \ok & mappedFiles[3].getUuid() deve essere "att-1" \\
\hline
TU-696 & \ok & mappedFiles[3].getFilename() deve essere "./att1.pdf" \\
\hline
TU-697 & \ok & mappedFiles[3].getPath() deve essere "class-2/aip-2/documents/doc-2/att1.pdf" \\
\hline
TU-698 & \ok & mappedFiles[3].getHash() deve essere "hash-att-1" \\
\hline
TU-699 & \ok & mappedFiles[3].getIsMain() deve essere false \\
\hline
TU-700 & \ok & mappedFiles[4].getUuid() deve essere "att-2" \\
\hline
TU-701 & \ok & mappedFiles[4].getFilename() deve essere "./att2.pdf" \\
\hline
TU-702 & \ok & mappedFiles[4].getPath() deve essere "class-2/aip-2/documents/doc-2/att2.pdf" \\
\hline
TU-703 & \ok & mappedFiles[4].getHash() deve essere "hash-att-2" \\
\hline
TU-704 & \ok & mappedFiles[4].getIsMain() deve essere false \\
\hline
TU-705 & \ok & ensureArray(undefined) deve essere uguale a [] \\
\hline
TU-706 & \ok & ensureArray(null) deve essere uguale a [] \\
\hline
TU-707 & \ok & ensureArray(input) deve essere input \\
\hline
TU-708 & \ok & ensureArray(1) deve essere uguale a [1] \\
\hline
TU-709 & \ok & ensureArray("test") deve essere uguale a ["test"] \\
\hline
TU-710 & \ok & ensureArray({ a: 1 }) deve essere uguale a [{ a: 1 }] \\
\hline
TU-711 & \ok & result deve essere un'istanza di Uint8Array \\
\hline
TU-712 & \ok & Buffer.from(result).toString("utf-8") deve essere "test content" \\
\hline
TU-713 & \ok & content deve essere "test content" \\
\hline
TU-714 & \ok & result deve essere "test content" \\
\hline
TU-715 & \ok & result deve essere true \\
\hline
TU-716 & \ok & result deve essere false \\
\hline
TU-717 & \ok & result.sort((a: string, b: string) => a.localeCompare(b)) deve essere uguale a [ "file1.txt", "file2.txt", ] \\
\hline
TU-718 & \ok & parsedObj.simpleNode deve essere "exampleText" \\
\hline
TU-719 & \ok & parsedObj.nestedNode.childNode deve essere uguale a ["childText", "child2Text"] \\
\hline
TU-720 & \ok & parsedObj deve essere uguale a {} \\
\hline
TU-721 & \ok & parsedObj.node["\#text"] deve essere "text" \\
\hline
TU-722 & \ok & mockDAO.save deve essere stato chiamato con vector \\
\hline
TU-723 & \ok & repository.saveVector(null) deve lanciare un errore \\
\hline
TU-724 & \ok & mockDAO.getByDocumentId deve essere stato chiamato con 1 \\
\hline
TU-725 & \ok & result deve essere uguale a new Float32Array([0.1, 0.2, 0.3]) \\
\hline
TU-726 & \ok & mockDAO.getByDocumentId deve essere stato chiamato con 1 \\
\hline
TU-727 & \ok & result deve essere null \\
\hline
TU-728 & \ok & mockDAO.getByDocumentId.toBeCalledWith(1) \\
\hline
TU-729 & \ok & result deve essere uguale a expectedVector \\
\hline
TU-730 & \ok & mockDAO.searchSimilar.toBeCalledWith(queryVector, topK) \\
\hline
TU-731 & \ok & result deve essere uguale a [ { documentId: 1, score: 0.9 }, { documentId: 2, score: 0.8 }, ] \\
\hline
TU-732 & \ok & fs.mkdirSync deve essere stato chiamato con string containing 'paraphrase-multilingual-MiniLM-L12-v2', { recursive: true } \\
\hline
TU-733 & \ok & fs.mkdirSync non deve essere stato chiamato \\
\hline
TU-734 & \ok & mkdirCall deve contenere 'Xenova' \\
\hline
TU-735 & \ok & mkdirCall deve contenere 'paraphrase-multilingual-MiniLM-L12-v2' \\
\hline
TU-736 & \ok & mkdirCall deve contenere 'onnx' \\
\hline
TU-737 & \ok & baseDir deve contenere 'Xenova' \\
\hline
TU-738 & \ok & fs.existsSync deve essere stato chiamato \\
\hline
TU-739 & \ok & https.get deve essere stato chiamato \\
\hline
TU-740 & \ok & http.get non deve essere stato chiamato \\
\hline
TU-741 & \ok & fs.createWriteStream deve essere stato chiamato con string containing 'model.onnx' \\
\hline
TU-742 & \ok & capturedRes.pipe deve essere stato chiamato con ws \\
\hline
TU-743 & \ok & ws.close deve essere stato chiamato \\
\hline
TU-744 & \ok & result deve contenere 'model.onnx' \\
\hline
TU-745 & \ok & result deve contenere '/tmp' \\
\hline
TU-746 & \ok & calls deve avere lunghezza 2 \\
\hline
TU-747 & \ok & calls[0] deve essere 'https://example.com/model.onnx' \\
\hline
TU-748 & \ok & calls[1] deve essere 'https://cdn.example.com/model.onnx' \\
\hline
TU-749 & \ok & ws.destroy deve essere stato chiamato \\
\hline
TU-750 & \ok & downloadFile('https://example.com/model.onnx', '/tmp', 'model.onnx') deve lanciare 'Errore HTTP 404' \\
\hline
TU-751 & \ok & downloadFile('https://example.com/model.onnx', '/tmp', 'model.onnx') deve lanciare 'Errore HTTP 500' \\
\hline
TU-752 & \ok & fs.unlink deve essere stato chiamato \\
\hline
TU-753 & \ok & ws.destroy deve essere stato chiamato \\
\hline
TU-754 & \ok & downloadFile('https://example.com/model.onnx', '/tmp', 'model.onnx') deve lanciare 'ECONNREFUSED' \\
\hline
TU-755 & \ok & fs.unlink deve essere stato chiamato \\
\hline
TU-756 & \ok & downloadFile('https://example.com/model.onnx', '/tmp', 'model.onnx') deve lanciare 'ENOSPC' \\
\hline
TU-757 & \ok & typeof text deve essere "string" \\
\hline
TU-758 & \ok & text.length deve essere maggiore di 0 \\
\hline
TU-759 & \ok & text.length deve essere minore di 10\_000\_000 \\
\hline
TU-760 & \ok & fileSystemProvider.openReadStream deve essere stato chiamato 1 volte \\
\hline
TU-761 & \ok & embedding deve essere un'istanza di Float32Array \\
\hline
TU-762 & \ok & embeddingAdapter.generateEmbedding deve essere stato chiamato \\
\hline
TU-763 & \ok & embedding deve essere un'istanza di Float32Array \\
\hline
TU-764 & \ok & embeddingAdapter.generateEmbedding deve essere stato chiamato 8 volte \\
\hline
TU-765 & \ok & result deve essere "VALID" \\
\hline
TU-766 & \ok & result deve essere "INVALID" \\
\hline
TU-767 & \ok & result deve essere "VALID" \\
\hline
TU-768 & \ok & h.service.checkFileIntegrityStatus(777) deve lanciare "File with id 777 not found" \\
\hline
TU-769 & \ok & h.transactionManager.runInTransaction deve essere stato chiamato 1 volte \\
\hline
TU-770 & \ok & h.service.checkDocumentIntegrityStatus(778) deve lanciare "Document with id 778 not found" \\
\hline
TU-771 & \ok & h.transactionManager.runInTransaction deve essere stato chiamato 1 volte \\
\hline
TU-772 & \ok & h.service.checkProcessIntegrityStatus(779) deve lanciare "Process with id 779 not found" \\
\hline
TU-773 & \ok & h.transactionManager.runInTransaction deve essere stato chiamato 1 volte \\
\hline
TU-774 & \ok & h.service.checkDocumentClassIntegrityStatus(780) deve lanciare "DocumentClass with id 780 not found" \\
\hline
TU-775 & \ok & h.transactionManager.runInTransaction deve essere stato chiamato 1 volte \\
\hline
TU-776 & \ok & h.service.checkDipIntegrityStatus(781) deve lanciare "Dip with id 781 not found" \\
\hline
TU-777 & \ok & h.transactionManager.runInTransaction deve essere stato chiamato 1 volte \\
\hline
TU-778 & \ok & status deve essere IntegrityStatusEnum.VALID \\
\hline
TU-779 & \ok & h.fileRepo.updateIntegrityStatus deve essere stato chiamato con 1, IntegrityStatusEnum.VALID \\
\hline
TU-780 & \ok & h.documentRepo.updateIntegrityStatus non deve essere stato chiamato \\
\hline
TU-781 & \ok & h.processRepo.updateIntegrityStatus non deve essere stato chiamato \\
\hline
TU-782 & \ok & h.documentClassRepo.updateIntegrityStatus non deve essere stato chiamato \\
\hline
TU-783 & \ok & h.dipRepo.updateIntegrityStatus non deve essere stato chiamato \\
\hline
TU-784 & \ok & status deve essere IntegrityStatusEnum.VALID \\
\hline
TU-785 & \ok & h.documentRepo.updateIntegrityStatus deve essere stato chiamato con 10, IntegrityStatusEnum.VALID \\
\hline
TU-786 & \ok & h.fileRepo.updateIntegrityStatus deve essere stato chiamato 2 volte \\
\hline
TU-787 & \ok & status deve essere IntegrityStatusEnum.VALID \\
\hline
TU-788 & \ok & h.processRepo.updateIntegrityStatus deve essere stato chiamato con 20, IntegrityStatusEnum.VALID \\
\hline
TU-789 & \ok & h.documentRepo.updateIntegrityStatus deve essere stato chiamato con 21, IntegrityStatusEnum.VALID \\
\hline
TU-790 & \ok & h.fileRepo.updateIntegrityStatus deve essere stato chiamato con 2101, IntegrityStatusEnum.VALID \\
\hline
TU-791 & \ok & status deve essere IntegrityStatusEnum.VALID \\
\hline
TU-792 & \ok & h.documentClassRepo.updateIntegrityStatus deve essere stato chiamato con 30, IntegrityStatusEnum.VALID \\
\hline
TU-793 & \ok & h.processRepo.updateIntegrityStatus deve essere stato chiamato con 31, IntegrityStatusEnum.VALID \\
\hline
TU-794 & \ok & h.documentRepo.updateIntegrityStatus deve essere stato chiamato con 32, IntegrityStatusEnum.VALID \\
\hline
TU-795 & \ok & h.fileRepo.updateIntegrityStatus deve essere stato chiamato con 3201, IntegrityStatusEnum.VALID \\
\hline
TU-796 & \ok & status deve essere IntegrityStatusEnum.VALID \\
\hline
TU-797 & \ok & h.transactionManager.runInTransaction deve essere stato chiamato 1 volte \\
\hline
TU-798 & \ok & fileValidUpdates deve avere lunghezza 16 \\
\hline
TU-799 & \ok & documentValidUpdates deve avere lunghezza 8 \\
\hline
TU-800 & \ok & processValidUpdates deve avere lunghezza 4 \\
\hline
TU-801 & \ok & documentClassValidUpdates deve avere lunghezza 2 \\
\hline
TU-802 & \ok & h.dipRepo.updateIntegrityStatus deve essere stato chiamato con 1, IntegrityStatusEnum.VALID \\
\hline
TU-803 & \ok & status deve essere IntegrityStatusEnum.INVALID \\
\hline
TU-804 & \ok & h.documentRepo.updateIntegrityStatus deve essere stato chiamato con hierarchy.invalidChain.documentId, IntegrityStatusEnum.INVALID \\
\hline
TU-805 & \ok & h.processRepo.updateIntegrityStatus deve essere stato chiamato con hierarchy.invalidChain.processId, IntegrityStatusEnum.INVALID \\
\hline
TU-806 & \ok & h.documentClassRepo.updateIntegrityStatus deve essere stato chiamato con hierarchy.invalidChain.documentClassId, IntegrityStatusEnum.INVALID \\
\hline
TU-807 & \ok & h.dipRepo.updateIntegrityStatus deve essere stato chiamato con 1, IntegrityStatusEnum.INVALID \\
\hline
TU-808 & \ok & status deve essere IntegrityStatusEnum.VALID \\
\hline
TU-809 & \ok & h.dipRepo.updateIntegrityStatus deve essere stato chiamato con 99, IntegrityStatusEnum.VALID \\
\hline
TU-810 & \ok & h.documentClassRepo.updateIntegrityStatus non deve essere stato chiamato \\
\hline
TU-811 & \ok & h.processRepo.updateIntegrityStatus non deve essere stato chiamato \\
\hline
TU-812 & \ok & h.documentRepo.updateIntegrityStatus non deve essere stato chiamato \\
\hline
TU-813 & \ok & h.fileRepo.updateIntegrityStatus non deve essere stato chiamato \\
\hline
TU-814 & \ok & result deve essere un'istanza di Float32Array \\
\hline
TU-815 & \ok & result.length deve essere 384 \\
\hline
TU-816 & \ok & norm deve essere vicino a 1, 2 \\
\hline
TU-817 & \ok & v deve essere maggiore o uguale a -1 \\
\hline
TU-818 & \ok & v deve essere minore o uguale a 1 \\
\hline
TU-819 & \ok & areDifferent deve essere true \\
\hline
TU-820 & \ok & v1[i] deve essere vicino a v2[i], 5 \\
\hline
TU-821 & \ok & cosineSimilarity(v1, v2) deve essere maggiore di 0.7 \\
\hline
TU-822 & \ok & cosineSimilarity(v1, v2) deve essere minore di 0.5 \\
\hline
TU-823 & \ok & cosineSimilarity(v, v) deve essere vicino a 1, 4 \\
\hline
TU-824 & \ok & cosineSimilarity(v1, v2) deve essere maggiore di 0.75 \\
\hline
TU-825 & \ok & cosineSimilarity(v1, v2) deve essere maggiore di 0.7 \\
\hline
TU-826 & \ok & cosineSimilarity(v1, v2) deve essere vicino a cosineSimilarity(v2, v1), 5 \\
\hline
TU-827 & \ok & cosineSimilarity(v1, v2) deve essere maggiore di 0.6 \\
\hline
TU-828 & \ok & score deve essere maggiore o uguale a 0 \\
\hline
TU-829 & \ok & score deve essere minore o uguale a 1 \\
\hline
TU-830 & \ok & result deve essere un'istanza di Float32Array \\
\hline
TU-831 & \ok & result.length deve essere 384 \\
\hline
TU-832 & \ok & result deve essere un'istanza di Float32Array \\
\hline
TU-833 & \ok & result.length deve essere 384 \\
\hline
TU-834 & \ok & result deve essere un'istanza di Float32Array \\
\hline
TU-835 & \ok & result.length deve essere 384 \\
\hline
TU-836 & \ok & result deve essere un'istanza di Float32Array \\
\hline
TU-837 & \ok & result.length deve essere 384 \\
\hline
TU-838 & \ok & vectorArg deve essere un'istanza di Float32Array \\
\hline
TU-839 & \ok & vectorArg.length deve essere 384 \\
\hline
TU-840 & \ok & v1 non deve essere uguale a v2 \\
\hline
TU-841 & \ok & cosineSimilarity(v1, v2) deve essere maggiore di 0.7 \\
\hline
TU-842 & \ok & cosineSimilarity(v1, v2) deve essere minore di 0.5 \\
\hline
TU-843 & \ok & results deve avere lunghezza 2 \\
\hline
TU-844 & \ok & results[0].score deve essere vicino a 0.9, 5 \\
\hline
TU-845 & \ok & results[1].score deve essere vicino a 0.6, 5 \\
\hline
TU-846 & \ok & results[0].document deve essere fakeDoc1 \\
\hline
TU-847 & \ok & results[1].document deve essere fakeDoc2 \\
\hline
TU-848 & \ok & Array.from(v1) deve essere uguale a Array.from(v2) \\
\hline
TU-849 & \ok & results deve avere lunghezza 1 \\
\hline
TU-850 & \ok & results[0].score deve essere vicino a 0.85, 5 \\
\hline
TU-851 & \ok & hasNonZero deve essere true \\
\hline
TU-852 & \ok & uc.execute("query senza documenti") deve lanciare "999" \\
\hline
TU-853 & \ok & adapter.isInitialized() deve essere false \\
\hline
TU-854 & \ok & adapter.isInitialized() deve essere true \\
\hline
TU-855 & \ok & pipeline deve essere stato chiamato con 'feature-extraction', 'paraphrase-multilingual-MiniLM-L12-v2', { quantized: true, local\_files\_only: true } \\
\hline
TU-856 & \ok & pipeline deve essere stato chiamato 1 volte \\
\hline
TU-857 & \ok & env.allowLocalModels deve essere true \\
\hline
TU-858 & \ok & env.allowRemoteModels deve essere false \\
\hline
TU-859 & \ok & env.useBrowserCache deve essere false \\
\hline
TU-860 & \ok & env.backends.onnx.executionProviders deve essere uguale a ['cpu'] \\
\hline
TU-861 & \ok & adapter.initialize() deve lanciare '[AI] Cartella modelli non trovata' \\
\hline
TU-862 & \ok & adapter.isInitialized() deve essere false \\
\hline
TU-863 & \ok & adapter.initialize() deve lanciare 'model load failed' \\
\hline
TU-864 & \ok & adapter.isInitialized() deve essere false \\
\hline
TU-865 & \ok & adapter.initialize() deve lanciare 'temporary load failure' \\
\hline
TU-866 & \ok & adapter.isInitialized() deve essere false \\
\hline
TU-867 & \ok & adapter.isInitialized() deve essere true \\
\hline
TU-868 & \ok & pipeline deve essere stato chiamato 2 volte \\
\hline
TU-869 & \ok & result deve essere un'istanza di Float32Array \\
\hline
TU-870 & \ok & result deve essere fakeVector \\
\hline
TU-871 & \ok & mockEmbedder deve essere stato chiamato con 'testo', { pooling: 'mean', normalize: true, } \\
\hline
TU-872 & \ok & adapter.isInitialized() deve essere false \\
\hline
TU-873 & \ok & adapter.isInitialized() deve essere true \\
\hline
TU-874 & \ok & pipeline deve essere stato chiamato 1 volte \\
\hline
TU-875 & \ok & mockEmbedder deve essere stato chiamato 3 volte \\
\hline
TU-876 & \ok & result.length deve essere 384 \\
\hline
TU-877 & \ok & result1[0] deve essere vicino a 0.1, 5 \\
\hline
TU-878 & \ok & result2[0] deve essere vicino a 0.9, 5 \\
\hline
TU-879 & \ok & result1 non deve essere result2 \\
\hline
TU-880 & \ok & adapter.generateEmbedding('testo') deve lanciare 'inference failed' \\
\hline
TU-881 & \ok & repo.getById deve essere stato chiamato con 5 \\
\hline
TU-882 & \ok & result deve essere entity \\
\hline
TU-883 & \ok & repo.getByDipId deve essere stato chiamato con 9 \\
\hline
TU-884 & \ok & result deve essere list \\
\hline
TU-885 & \ok & repo.getByStatus deve essere stato chiamato con IntegrityStatusEnum.INVALID \\
\hline
TU-886 & \ok & result deve essere list \\
\hline
TU-887 & \ok & integrityService.checkDocumentClassIntegrityStatus deve essere stato chiamato con 1 \\
\hline
TU-888 & \ok & result deve essere IntegrityStatusEnum.UNKNOWN \\
\hline
TU-889 & \ok & uc.execute(14) deve lanciare "DocumentClass with id 14 not found" \\
\hline
TU-890 & \ok & repo.getById deve essere stato chiamato con 10 \\
\hline
TU-891 & \ok & result deve essere dip \\
\hline
TU-892 & \ok & repo.getById deve essere stato chiamato con 999 \\
\hline
TU-893 & \ok & result deve essere null \\
\hline
TU-894 & \ok & repo.getByStatus deve essere stato chiamato con IntegrityStatusEnum.VALID \\
\hline
TU-895 & \ok & result deve essere list \\
\hline
TU-896 & \ok & repo.getByStatus deve essere stato chiamato con IntegrityStatusEnum.INVALID \\
\hline
TU-897 & \ok & result deve essere uguale a [] \\
\hline
TU-898 & \ok & integrityService.checkDipIntegrityStatus deve essere stato chiamato con 1 \\
\hline
TU-899 & \ok & result deve essere IntegrityStatusEnum.VALID \\
\hline
TU-900 & \ok & uc.execute(3) deve lanciare "Dip with id 3 not found" \\
\hline
TU-901 & \ok & repo.getById deve essere stato chiamato con 10 \\
\hline
TU-902 & \ok & result deve essere entity \\
\hline
TU-903 & \ok & repo.getByProcessId deve essere stato chiamato con 3 \\
\hline
TU-904 & \ok & result deve essere list \\
\hline
TU-905 & \ok & repo.getByStatus deve essere stato chiamato con IntegrityStatusEnum.VALID \\
\hline
TU-906 & \ok & result deve essere list \\
\hline
TU-907 & \ok & integrityService.checkDocumentIntegrityStatus deve essere stato chiamato con 1 \\
\hline
TU-908 & \ok & result deve essere IntegrityStatusEnum.INVALID \\
\hline
TU-909 & \ok & uc.execute(77) deve lanciare "Document with id 77 not found" \\
\hline
TU-910 & \ok & repo.getById deve essere stato chiamato con 11 \\
\hline
TU-911 & \ok & result deve essere entity \\
\hline
TU-912 & \ok & repo.getByDocumentId deve essere stato chiamato con 8 \\
\hline
TU-913 & \ok & result deve essere list \\
\hline
TU-914 & \ok & repo.getByStatus deve essere stato chiamato con IntegrityStatusEnum.UNKNOWN \\
\hline
TU-915 & \ok & result deve essere list \\
\hline
TU-916 & \ok & result deve essere IntegrityStatusEnum.VALID \\
\hline
TU-917 & \ok & integrityService.checkFileIntegrityStatus deve essere stato chiamato con 8 \\
\hline
TU-918 & \ok & result deve essere IntegrityStatusEnum.UNKNOWN \\
\hline
TU-919 & \ok & integrityService.checkFileIntegrityStatus deve essere stato chiamato con 8 \\
\hline
TU-920 & \ok & uc.execute(99) deve lanciare "File with id 99 not found" \\
\hline
TU-921 & \ok & result.success deve essere false \\
\hline
TU-922 & \ok & result.errorCode deve essere "NOT\_FOUND" \\
\hline
TU-923 & \ok & dialogPort.showSaveDialog non deve essere stato chiamato \\
\hline
TU-924 & \ok & exportPort.exportFile non deve essere stato chiamato \\
\hline
TU-925 & \ok & result.success deve essere false \\
\hline
TU-926 & \ok & result.canceled deve essere true \\
\hline
TU-927 & \ok & exportPort.exportFile non deve essere stato chiamato \\
\hline
TU-928 & \ok & openReadStreamArg deve contenere "base/dip" \\
\hline
TU-929 & \ok & openReadStreamArg deve contenere "docs/f.pdf" \\
\hline
TU-930 & \ok & exportPort.exportFile deve essere stato chiamato con stream, "/dest/f.pdf" \\
\hline
TU-931 & \ok & result deve essere exported \\
\hline
TU-932 & \ok & makeUC().execute(1) deve lanciare "disk full" \\
\hline
TU-933 & \ok & result.canceled deve essere true \\
\hline
TU-934 & \ok & result.results deve avere lunghezza 0 \\
\hline
TU-935 & \ok & onProgress non deve essere stato chiamato \\
\hline
TU-936 & \ok & result.canceled deve essere false \\
\hline
TU-937 & \ok & result.results[0] deve corrispondere a { fileId: 99, success: false } \\
\hline
TU-938 & \ok & onProgress deve essere stato chiamato con 1, 1 \\
\hline
TU-939 & \ok & streamArg deve essere stream \\
\hline
TU-940 & \ok & destPathArg deve contenere "/dest" \\
\hline
TU-941 & \ok & destPathArg deve contenere "f.pdf" \\
\hline
TU-942 & \ok & result.canceled deve essere false \\
\hline
TU-943 & \ok & result.results[0] deve corrispondere a { fileId: 1, success: true } \\
\hline
TU-944 & \ok & onProgress deve essere stato chiamato con 1, 1 \\
\hline
TU-945 & \ok & result.success deve essere false \\
\hline
TU-946 & \ok & result.errorCode deve essere "NOT\_FOUND" \\
\hline
TU-947 & \ok & printPort.printSingle non deve essere stato chiamato \\
\hline
TU-948 & \ok & absolutePathArg deve contenere "base/dip" \\
\hline
TU-949 & \ok & absolutePathArg deve contenere "docs/report.pdf" \\
\hline
TU-950 & \ok & optsArg deve essere uguale a { silent: false, printBackground: true } \\
\hline
TU-951 & \ok & result deve essere ok \\
\hline
TU-952 & \ok & makeUC().execute(1) deve lanciare "unexpected crash" \\
\hline
TU-953 & \ok & fileRepo.getById deve essere stato chiamato con 42 \\
\hline
TU-954 & \ok & result.canceled deve essere true \\
\hline
TU-955 & \ok & result.results deve avere lunghezza 0 \\
\hline
TU-956 & \ok & onProgress non deve essere stato chiamato \\
\hline
TU-957 & \ok & result.canceled deve essere false \\
\hline
TU-958 & \ok & result.results[0] deve corrispondere a { fileId: 99, success: false } \\
\hline
TU-959 & \ok & onProgress deve essere stato chiamato con 1, 1 \\
\hline
TU-960 & \ok & absolutePathArg deve contenere "base/dip" \\
\hline
TU-961 & \ok & absolutePathArg deve contenere "docs/f.pdf" \\
\hline
TU-962 & \ok & optsArg deve essere uguale a { silent: false, printBackground: true } \\
\hline
TU-963 & \ok & result.canceled deve essere false \\
\hline
TU-964 & \ok & result.results[0] deve corrispondere a { fileId: 1, success: true } \\
\hline
TU-965 & \ok & onProgress deve essere stato chiamato con 1, 1 \\
\hline
TU-966 & \ok & dialogPort.showConfirmPrint deve essere stato chiamato con 3 \\
\hline
TU-967 & \ok & repo.getById deve essere stato chiamato con 12 \\
\hline
TU-968 & \ok & result deve essere entity \\
\hline
TU-969 & \ok & repo.getByDocumentClassId deve essere stato chiamato con 3 \\
\hline
TU-970 & \ok & result deve essere list \\
\hline
TU-971 & \ok & repo.getByStatus deve essere stato chiamato con IntegrityStatusEnum.VALID \\
\hline
TU-972 & \ok & result deve essere list \\
\hline
TU-973 & \ok & integrityService.checkProcessIntegrityStatus deve essere stato chiamato con 1 \\
\hline
TU-974 & \ok & result deve essere IntegrityStatusEnum.VALID \\
\hline
TU-975 & \ok & uc.execute(55) deve lanciare "Process with id 55 not found" \\
\hline
TU-976 & \ok & results deve avere lunghezza 1 \\
\hline
TU-977 & \ok & results[0].getId() deve essere 1 \\
\hline
TU-978 & \ok & results[0].getName() deve essere "Contratti" \\
\hline
TU-979 & \ok & results[0].getUuid() deve essere "uuid-1" \\
\hline
TU-980 & \ok & repo.searchDocumentalClasses deve essere stato chiamato con "Fatture 2026" \\
\hline
TU-981 & \ok & results deve avere lunghezza 0 \\
\hline
TU-982 & \ok & repo.searchDocumentalClasses deve essere stato chiamato con "" \\
\hline
TU-983 & \ok & results deve avere lunghezza 3 \\
\hline
TU-984 & \ok & results.map((r) => r.getName()) deve essere uguale a ["AAA", "BBB", "CCC"] \\
\hline
TU-985 & \ok & results[0].getId() deve essere 1 \\
\hline
TU-986 & \ok & results[0].getName() deve essere "Originale" \\
\hline
TU-987 & \ok & results[0].getIntegrityStatus() deve essere "UNKNOWN" \\
\hline
TU-988 & \ok & () => uc.execute("Contratti") deve lanciare "document class search failed" \\
\hline
TU-989 & \ok & results deve avere lunghezza 1 \\
\hline
TU-990 & \ok & results[0].getId() deve essere 1 \\
\hline
TU-991 & \ok & results[0].getUuid() deve essere "proc-uuid-abc" \\
\hline
TU-992 & \ok & results[0].getMetadata().findNodeByName("name")?.getStringValue() deve essere "Proc A" \\
\hline
TU-993 & \ok & results[0].getMetadata().findNodeByName("type")?.getStringValue() deve essere "PROCESSO" \\
\hline
TU-994 & \ok & repo.searchProcesses deve essere stato chiamato con "uuid-specifico-123" \\
\hline
TU-995 & \ok & results deve avere lunghezza 0 \\
\hline
TU-996 & \ok & repo.searchProcesses deve essere stato chiamato con "" \\
\hline
TU-997 & \ok & results deve avere lunghezza 2 \\
\hline
TU-998 & \ok & results.map((r) => r.getId()) deve essere uguale a [1, 2, 3] \\
\hline
TU-999 & \ok & results[0].getMetadata().findNodeByName("name") deve essere null \\
\hline
TU-1000 & \ok & results[0].getMetadata().findNodeByName("type") deve essere null \\
\hline
TU-1001 & \ok & repo.searchProcesses deve essere stato chiamato con "proc-" \\
\hline
TU-1002 & \ok & () => uc.execute("proc-uuid") deve lanciare "process search failed" \\
\hline
TU-1003 & \ok & results deve avere lunghezza 1 \\
\hline
TU-1004 & \ok & results[0].getUuid() deve essere "uuid-1" \\
\hline
TU-1005 & \ok & results[0] .getMetadata() .findNodeByName("NomeDelDocumento") ?.getStringValue() deve essere "fattura.pdf" \\
\hline
TU-1006 & \ok & results[0] .getMetadata() .findNodeByName("tipoDocumento") ?.getStringValue() deve essere "DOCUMENTO INFORMATICO" \\
\hline
TU-1007 & \ok & results deve avere lunghezza 0 \\
\hline
TU-1008 & \ok & results[0].getMetadata().findNodeByName("nome") deve essere null \\
\hline
TU-1009 & \ok & results[0].getMetadata().findNodeByName("tipoDocumento") deve essere null \\
\hline
TU-1010 & \ok & repo.searchDocument deve essere stato chiamato con filters \\
\hline
TU-1011 & \ok & results deve avere lunghezza 3 \\
\hline
TU-1012 & \ok & results.map((r) => r.getUuid()) deve essere uguale a [ "uuid-a", "uuid-b", "uuid-c", ] \\
\hline
TU-1013 & \ok & results[0] .getMetadata() .findNodeByName("NomeDelDocumento") ?.getStringValue() deve essere "contratto.pdf" \\
\hline
TU-1014 & \ok & results[0] .getMetadata() .findNodeByName("tipoDocumento") ?.getStringValue() deve essere "DOCUMENTO AMMINISTRATIVO INFORMATICO" \\
\hline
TU-1015 & \ok & results[0] .getMetadata() .findNodeByName("NomeDelDocumento") ?.getStringValue() deve essere "primo-nome.pdf" \\
\hline
TU-1016 & \ok & results[0] .getMetadata() .findNodeByName("tipoDocumento") ?.getStringValue() deve essere "TIPO-1" \\
\hline
TU-1017 & \ok & () => uc.execute(emptyFilters) deve lanciare "db unavailable" \\
\hline
TU-1018 & \ok & results deve avere lunghezza 1 \\
\hline
TU-1019 & \ok & results[0].document.getUuid() deve essere "uuid-s1" \\
\hline
TU-1020 & \ok & results[0].score deve essere 0.92 \\
\hline
TU-1021 & \ok & typeof results[0].score deve essere "number" \\
\hline
TU-1022 & \ok & results[0].score non deve essere null \\
\hline
TU-1023 & \ok & results deve avere lunghezza 0 \\
\hline
TU-1024 & \ok & aiAdapter.generateEmbedding deve essere stato chiamato con "ricerca semantica avanzata" \\
\hline
TU-1025 & \ok & vectorRepo.searchSimilarVectors deve essere stato chiamato con embedding, 10 \\
\hline
TU-1026 & \ok & results[0].score deve essere 0.95 \\
\hline
TU-1027 & \ok & results[1].score deve essere 0.78 \\
\hline
TU-1028 & \ok & results[2].score deve essere 0.61 \\
\hline
TU-1029 & \ok & results[0].score deve essere 1 \\
\hline
TU-1030 & \ok & results[1].score deve essere 0 \\
\hline
TU-1031 & \ok & results[0].document.getUuid() deve essere "uuid-empty" \\
\hline
TU-1032 & \ok & results[0].score deve essere 0.88 \\
\hline
TU-1033 & \ok & results[0].document.getUuid() deve essere "uuid-dup" \\
\hline
TU-1034 & \ok & results[0].score deve essere 0.5 \\
\hline
TU-1035 & \ok & uc.execute("query") deve lanciare "semantic fail" \\
\hline
TU-1036 & \ok & true deve essere true \\
\hline
TU-1037 & \ok & true deve essere true \\
\hline
TU-1038 & \ok & result.success deve essere true \\
\hline
TU-1039 & \ok & durationsMs deve avere lunghezza runs \\
\hline
TU-1040 & \ok & result deve essere uguale a { success: true } \\
\hline
TU-1041 & \ok & dips deve avere lunghezza 1 \\
\hline
TU-1042 & \ok & (dips[0] as any).uuid deve essere "dip-uuid-1" \\
\hline
TU-1043 & \ok & docClasses deve avere lunghezza 1 \\
\hline
TU-1044 & \ok & (docClasses[0] as any).uuid deve essere "dc-1" \\
\hline
TU-1045 & \ok & processes deve avere lunghezza 1 \\
\hline
TU-1046 & \ok & (processes[0] as any).uuid deve essere "proc-1" \\
\hline
TU-1047 & \ok & documents deve avere lunghezza 1 \\
\hline
TU-1048 & \ok & (documents[0] as any).uuid deve essere "doc-1" \\
\hline
TU-1049 & \ok & files deve avere lunghezza 1 \\
\hline
TU-1050 & \ok & (files[0] as any).filename deve essere "./primary.pdf" \\
\hline
TU-1051 & \ok & embeddingService.generateDocumentEmbedding deve essere stato chiamato 1 volte \\
\hline
TU-1052 & \ok & calledFile.getPath() deve contenere "primary.pdf" \\
\hline
TU-1053 & \ok & vectorRepository.saveVector deve essere stato chiamato 1 volte \\
\hline
TU-1054 & \ok & vectorRepository.saveVector deve essere stato chiamato con any Vector \\
\hline
TU-1055 & \ok & useCase.execute(dipPath) deve lanciare un errore \\
\hline
TU-1056 & \ok & dips deve avere lunghezza 0 \\
\hline
TU-1057 & \ok & docClasses deve avere lunghezza 0 \\
\hline
TU-1058 & \ok & result deve essere uguale a { success: true } \\
\hline
TU-1059 & \ok & files deve avere lunghezza 1 \\
\hline
TU-1060 & \ok & embeddingService.generateDocumentEmbedding deve essere stato chiamato 1 volte \\
\hline
TU-1061 & \ok & vectorRepository.saveVector non deve essere stato chiamato \\
\hline
TU-1062 & \ok & cacheByToken deve essere cacheByClass \\
\hline
TU-1063 & \ok & location.path() deve essere '/browse' \\
\hline
TU-1064 & \ok & searchRoute deve essere truthy \\
\hline
TU-1065 & \ok & searchRoute?.loadChildren deve essere definito \\
\hline
TU-1066 & \ok & typeof searchRoute.loadChildren deve essere 'function' \\
\hline
TU-1067 & \ok & detailRoute deve essere truthy \\
\hline
TU-1068 & \ok & detailRoute?.loadChildren deve essere definito \\
\hline
TU-1069 & \ok & typeof detailRoute.loadChildren deve essere 'function' \\
\hline
TU-1070 & \ok & integrityRoute deve essere truthy \\
\hline
TU-1071 & \ok & integrityRoute?.loadChildren deve essere definito \\
\hline
TU-1072 & \ok & typeof integrityRoute.loadChildren deve essere 'function' \\
\hline
TU-1073 & \ok & location.path() deve essere '/browse' \\
\hline
TU-1074 & \ok & location.path() deve essere '/browse' \\
\hline
TU-1075 & \ok & app deve essere truthy \\
\hline
TU-1076 & \ok & compiled.querySelector('router-outlet') non deve essere null \\
\hline
TU-1077 & \ok & mockBridge.invoke deve essere stato chiamato con 'ipc:log:write', mockLogEntry \\
\hline
TU-1078 & \ok & errorResult deve essere mockError \\
\hline
TU-1079 & \ok & mockBridge.invoke deve essere stato chiamato 1 volte \\
\hline
TU-1080 & \ok & error.code deve essere ErrorCode.DOC\_BLOB\_LOAD\_ERROR \\
\hline
TU-1081 & \ok & error.message deve essere 'File non trovato' \\
\hline
TU-1082 & \ok & error.source deve essere 'DocFacade' \\
\hline
TU-1083 & \ok & error.category deve essere ErrorCategory.IO \\
\hline
TU-1084 & \ok & error.severity deve essere ErrorSeverity.ERROR \\
\hline
TU-1085 & \ok & error.recoverable deve essere true \\
\hline
TU-1086 & \ok & error.severity deve essere ErrorSeverity.FATAL \\
\hline
TU-1087 & \ok & error.recoverable deve essere false \\
\hline
TU-1088 & \ok & result.code deve essere ErrorCode.DOC\_FORMAT\_UNSUPPORTED \\
\hline
TU-1089 & \ok & result.message deve essere 'Formato .xyz non supportato' \\
\hline
TU-1090 & \ok & result.source deve essere 'ElectronRenderer' \\
\hline
TU-1091 & \ok & result.category deve essere ErrorCategory.IO \\
\hline
TU-1092 & \ok & result.code deve essere ErrorCode.SEARCH\_ENGINE\_ERROR \\
\hline
TU-1093 & \ok & result.message deve essere 'Errore misterioso nel motore di ricerca' \\
\hline
TU-1094 & \ok & result.detail deve essere rawError.stack \\
\hline
TU-1095 & \ok & result.code deve essere ErrorCode.IPC\_ERROR \\
\hline
TU-1096 & \ok & result.category deve essere ErrorCategory.IPC \\
\hline
TU-1097 & \ok & result.code deve essere ErrorCode.SEARCH\_ENGINE\_ERROR \\
\hline
TU-1098 & \ok & result.message deve essere 'Qualcosa è andato storto' \\
\hline
TU-1099 & \ok & result.source deve essere 'IPC Gateway' \\
\hline
TU-1100 & \ok & error.category deve essere ErrorCategory.IPC \\
\hline
TU-1101 & \ok & error.severity deve essere ErrorSeverity.ERROR \\
\hline
TU-1102 & \ok & error.recoverable deve essere false \\
\hline
TU-1103 & \ok & result.code deve essere ErrorCode.SEARCH\_ENGINE\_ERROR \\
\hline
TU-1104 & \ok & result deve essere uguale a { id: 1, name: 'Documento' } \\
\hline
TU-1105 & \ok & result deve essere null \\
\hline
TU-1106 & \ok & result deve essere null \\
\hline
TU-1107 & \ok & cacheService.get('manual-key') deve essere null \\
\hline
TU-1108 & \ok & cacheService.get('search/123') deve essere null \\
\hline
TU-1109 & \ok & cacheService.get('search/456') deve essere null \\
\hline
TU-1110 & \ok & cacheService.get('doc/789') deve essere 'dati documento' \\
\hline
TU-1111 & \ok & mockLoggingChannel.writeLog deve essere stato chiamato con { level: LogLevel.INFO, event: TelemetryEvent.SEARCH\_EXECUTED, timestamp: MOCK\_TIME, payload: null, } \\
\hline
TU-1112 & \ok & mockLoggingChannel.writeLog deve essere stato chiamato con { level: LogLevel.INFO, event: TelemetryMetric.SEARCH\_LATENCY\_MS, timestamp: MOCK\_TIME, payload: { durationMs: 350 }, } \\
\hline
TU-1113 & \ok & mockLoggingChannel.writeLog deve essere stato chiamato con { level: LogLevel.ERROR, event: 'APP\_ERROR', timestamp: MOCK\_TIME, payload: mockError, } \\
\hline
TU-1114 & \ok & () => service.trackEvent(TelemetryEvent.SEARCH\_EXECUTED) non deve lanciare un errore \\
\hline
TU-1115 & \ok & console.error deve essere stato chiamato con 'Fallimento scrittura telemetria via IPC:', logError \\
\hline
TU-1116 & \ok & result.idAgg!.tipoAggregazione deve essere 'Fascicolo' \\
\hline
TU-1117 & \ok & result.idAgg!.idAggregazione deve essere '' \\
\hline
TU-1118 & \ok & result.tipologiaFascicolo deve essere 'procedimento amministrativo' \\
\hline
TU-1119 & \ok & result.soggetti deve essere uguale a [{ tipoRuolo: 'Sconosciuto', denominazione: 'N/A' }] \\
\hline
TU-1120 & \ok & result.dataApertura deve essere 'N/D' \\
\hline
TU-1121 & \ok & result.progressivo deve essere 1 \\
\hline
TU-1122 & \ok & result.chiaveDescrittiva!.oggetto deve essere 'Fascicolo n. ' \\
\hline
TU-1123 & \ok & result.procedimentoAmministrativo!.materiaArgomentoStruttura deve essere 'Standard' \\
\hline
TU-1124 & \ok & result.procedimentoAmministrativo!.procedimento deve essere 'N/A' \\
\hline
TU-1125 & \ok & result.procedimentoAmministrativo!.fasi deve essere uguale a [] \\
\hline
TU-1126 & \ok & result.posizioneFisicaAggregazioneDocumentale deve essere 'Archivio' \\
\hline
TU-1127 & \ok & result.tempoDiConservazione deve essere 10 \\
\hline
TU-1128 & \ok & result.indiceDocumenti deve essere uguale a [] \\
\hline
TU-1129 & \ok & result.customMetadata deve essere uguale a [] \\
\hline
TU-1130 & \ok & result.idAgg.tipoAggregazione deve essere 'Serie' \\
\hline
TU-1131 & \ok & result.idAgg.idAggregazione deve essere 'AGG-001' \\
\hline
TU-1132 & \ok & result.tipologiaFascicolo deve essere 'affare' \\
\hline
TU-1133 & \ok & result.dataApertura deve essere '2023-01-01T00:00:00Z' \\
\hline
TU-1134 & \ok & result.dataChiusura deve essere '2024-01-01T00:00:00Z' \\
\hline
TU-1135 & \ok & result.progressivo deve essere 42 \\
\hline
TU-1136 & \ok & result.chiaveDescrittiva.oggetto deve essere 'Fascicolo di Test' \\
\hline
TU-1137 & \ok & result.chiaveDescrittiva.paroleChiave deve essere 'test, mock' \\
\hline
TU-1138 & \ok & result.posizioneFisicaAggregazioneDocumentale deve essere 'Scaffale A' \\
\hline
TU-1139 & \ok & result.tempoDiConservazione deve essere 99 \\
\hline
TU-1140 & \ok & result.note deve essere 'Nota di test' \\
\hline
TU-1141 & \ok & result.idAggPrimario deve essere 'PRIM-001' \\
\hline
TU-1142 & \ok & result.processSummary!.uuid deve essere 'abc-uuid' \\
\hline
TU-1143 & \ok & result.processSummary!.integrityStatus deve essere 'VALID' \\
\hline
TU-1144 & \ok & result.processSummary!.timestamp deve essere '2023-01-01T00:00:00Z' \\
\hline
TU-1145 & \ok & result.assegnazione.tipoAssegnazione deve essere 'Per conoscenza' \\
\hline
TU-1146 & \ok & result.assegnazione.soggettoAssegnatario.denominazione deve essere 'Mario Rossi' \\
\hline
TU-1147 & \ok & result.assegnazione.soggettoAssegnatario.codiceFiscale deve essere 'RSSMRA' \\
\hline
TU-1148 & \ok & result.assegnazione.dataInizioAssegnazione deve essere '2023-05-01' \\
\hline
TU-1149 & \ok & result.assegnazione.dataFineAssegnazione deve essere '2023-06-01' \\
\hline
TU-1150 & \ok & result.soggetti deve avere lunghezza 1 \\
\hline
TU-1151 & \ok & result.soggetti[0] deve essere uguale a { tipoRuolo: 'Mittente', denominazione: 'Giuseppe Verdi', codiceFiscale: 'VRDGPP', indirizzoDigitale: 'giuseppe@pec.it', } \\
\hline
TU-1152 & \ok & result.soggetti deve avere lunghezza 1 \\
\hline
TU-1153 & \ok & result.soggetti[0] deve essere uguale a { tipoRuolo: 'Destinatario', denominazione: 'Azienda Test SRL', codiceFiscale: '12345678901', indirizzoDigitale: 'azienda@pec.it', } \\
\hline
TU-1154 & \ok & result.soggetti deve avere lunghezza 1 \\
\hline
TU-1155 & \ok & result.soggetti[0].tipoRuolo deve essere 'Produttore' \\
\hline
TU-1156 & \ok & result.soggetti[0].denominazione deve essere 'IPA123' \\
\hline
TU-1157 & \ok & result.soggetti deve avere lunghezza 1 \\
\hline
TU-1158 & \ok & result.soggetti[0].tipoRuolo deve essere 'RUP' \\
\hline
TU-1159 & \ok & result.soggetti[0].denominazione deve essere 'Luca Bianchi' \\
\hline
TU-1160 & \ok & result.soggetti[0].denominazione deve essere 'Soggetto Generico' \\
\hline
TU-1161 & \ok & result.soggetti[0].tipoRuolo deve essere 'Sconosciuto' \\
\hline
TU-1162 & \ok & result.indiceDocumenti deve avere lunghezza 1 \\
\hline
TU-1163 & \ok & result.indiceDocumenti[0] deve essere uguale a { tipoDocumento: 'DocumentoAmministativoinformatico', identificativo: 'DOC-AMM-1', impronta: 'HASH123', } \\
\hline
TU-1164 & \ok & result.indiceDocumenti deve avere lunghezza 1 \\
\hline
TU-1165 & \ok & result.indiceDocumenti[0] deve essere uguale a { tipoDocumento: 'Documentoinformatico', identificativo: 'DOC-INF-2', impronta: 'HASH456', } \\
\hline
TU-1166 & \ok & result.indiceDocumenti deve avere lunghezza 1 \\
\hline
TU-1167 & \ok & result.indiceDocumenti[0] deve essere uguale a { tipoDocumento: 'Documento', identificativo: 'DOC-GEN-3', impronta: undefined, } \\
\hline
TU-1168 & \ok & result.procedimentoAmministrativo!.fasi deve avere lunghezza 2 \\
\hline
TU-1169 & \ok & result.procedimentoAmministrativo!.fasi[0] deve essere uguale a { tipoFase: 'Istruttoria', dataInizio: '2023-01-01', dataFine: '2023-01-31', } \\
\hline
TU-1170 & \ok & result.procedimentoAmministrativo!.fasi[1] deve essere uguale a { tipoFase: 'Decisionale', dataInizio: '2023-02-01', dataFine: '2023-02-15', } \\
\hline
TU-1171 & \ok & result.customMetadata deve essere uguale a array containing [ { nome: 'CampoExtra1', valore: 'Valore 1' }, { nome: 'CampoExtra2', valore: 'Valore 2' }, ] \\
\hline
TU-1172 & \ok & result.customMetadata deve essere uguale a array containing [{ nome: 'UnknownTag123', valore: 'Value123' }] \\
\hline
TU-1173 & \ok & mockCache.get deve essere stato chiamato con 'aggregate:123' \\
\hline
TU-1174 & \ok & mockGateway.invoke non deve essere stato chiamato \\
\hline
TU-1175 & \ok & facade.getState()().detail deve essere uguale a mockAggregateData \\
\hline
TU-1176 & \ok & mockTelemetry.trackTiming deve essere stato chiamato \\
\hline
TU-1177 & \ok & mockGateway.invoke deve essere stato chiamato con IpcChannels.BROWSE\_GET\_PROCESS\_BY\_ID, 123, null \\
\hline
TU-1178 & \ok & mockGateway.invoke deve essere stato chiamato con IpcChannels.BROWSE\_GET\_DOCUMENTS\_BY\_PROCESS, 123, null \\
\hline
TU-1179 & \ok & mockCache.set deve essere stato chiamato con 'aggregate:123', any Object, 300000 \\
\hline
TU-1180 & \ok & facade.getState()().detail deve essere truthy \\
\hline
TU-1181 & \ok & mockErrorHandler.handle deve essere stato chiamato con rawError \\
\hline
TU-1182 & \ok & facade.getState()().error deve essere uguale a mappedError \\
\hline
TU-1183 & \ok & mockTelemetry.trackError deve essere stato chiamato con mappedError \\
\hline
TU-1184 & \ok & service deve essere truthy \\
\hline
TU-1185 & \ok & mockElectronAPI.classeDocumentale.list deve essere stato chiamato \\
\hline
TU-1186 & \ok & result deve essere uguale a mockClasses \\
\hline
TU-1187 & \ok & mockElectronAPI.classeDocumentale.get deve essere stato chiamato con 1 \\
\hline
TU-1188 & \ok & result deve essere uguale a mockClasse \\
\hline
TU-1189 & \ok & mockElectronAPI.classeDocumentale.create deve essere stato chiamato con 'Nuova Classe' \\
\hline
TU-1190 & \ok & result deve essere uguale a mockClasse \\
\hline
TU-1191 & \ok & result.documentId deve essere '' \\
\hline
TU-1192 & \ok & result.fileName deve essere 'Documento' \\
\hline
TU-1193 & \ok & result.mimeType deve essere MimeType.UNSUPPORTED \\
\hline
TU-1194 & \ok & result.metadata!.identificativo deve essere 'N/A' \\
\hline
TU-1195 & \ok & result.metadata!.impronta deve essere 'N/A' \\
\hline
TU-1196 & \ok & result.integrityStatus deve essere undefined \\
\hline
TU-1197 & \ok & result.classification!.indice deve essere 'N/A' \\
\hline
TU-1198 & \ok & result.registration!.tipoRegistro deve essere 'N/A' \\
\hline
TU-1199 & \ok & result.format!.tipo deve essere 'N/A' \\
\hline
TU-1200 & \ok & result.verification!.firmaDigitale deve essere 'N/A' \\
\hline
TU-1201 & \ok & result.attachments!.numero deve essere 0 \\
\hline
TU-1202 & \ok & result.attachments!.allegati deve essere uguale a [] \\
\hline
TU-1203 & \ok & result.subjects deve essere uguale a [] \\
\hline
TU-1204 & \ok & result.customMetadata deve essere uguale a [] \\
\hline
TU-1205 & \ok & result.documentId deve essere 'DB-ID-123' \\
\hline
TU-1206 & \ok & result.aipInfo!.uuid deve essere 'UUID-456' \\
\hline
TU-1207 & \ok & result.fileName deve essere 'relazione\_tecnica.pdf' \\
\hline
TU-1208 & \ok & result.mimeType deve essere MimeType.PDF \\
\hline
TU-1209 & \ok & result.documentId deve essere 'META-ID-789' \\
\hline
TU-1210 & \ok & result.fileName deve essere 'image\_no\_ext' \\
\hline
TU-1211 & \ok & result.mimeType deve essere MimeType.IMAGE \\
\hline
TU-1212 & \ok & result.metadata!.identificativo deve essere 'DOC-123' \\
\hline
TU-1213 & \ok & result.metadata!.impronta deve essere 'HASH123' \\
\hline
TU-1214 & \ok & result.metadata!.algoritmoImpronta deve essere 'SHA-512' \\
\hline
TU-1215 & \ok & result.metadata!.oggetto deve essere 'Oggetto Test' \\
\hline
TU-1216 & \ok & result.metadata!.tipoDocumentale deve essere 'Contratto' \\
\hline
TU-1217 & \ok & result.metadata!.riservatezza deve essere 'true' \\
\hline
TU-1218 & \ok & result.metadata!.versione deve essere '2.0' \\
\hline
TU-1219 & \ok & result.metadata!.tempoDiConservazione deve essere '99' \\
\hline
TU-1220 & \ok & result.metadata!.paroleChiave deve essere uguale a ['test', 'mock'] \\
\hline
TU-1221 & \ok & result.registration!.tipoRegistro deve essere 'Protocollo Ufficiale' \\
\hline
TU-1222 & \ok & result.registration!.data deve essere '2023-01-01' \\
\hline
TU-1223 & \ok & result.registration!.numero deve essere 'PROT-001' \\
\hline
TU-1224 & \ok & result.registration!.tipoRegistro deve essere 'Registro Interno Nested' \\
\hline
TU-1225 & \ok & result.registration!.tipoRegistro deve essere 'Repertorio 123' \\
\hline
TU-1226 & \ok & result.subjects deve avere lunghezza 1 \\
\hline
TU-1227 & \ok & result.subjects![0].tipo deve essere SubjectType.PF \\
\hline
TU-1228 & \ok & result.subjects![0].ruolo deve essere 'Autore' \\
\hline
TU-1229 & \ok & result.subjects![0].campiSpecifici['Nome'] deve essere 'Mario' \\
\hline
TU-1230 & \ok & result.subjects![0].campiSpecifici['CodiceFiscale'] deve essere 'RSSMRA' \\
\hline
TU-1231 & \ok & result.subjects![0].campiSpecifici['IndirizziDigitaliDiRiferimento'] deve essere 'mario.rossi@pec.it' \\
\hline
TU-1232 & \ok & result.subjects![0].tipo deve essere SubjectType.PG \\
\hline
TU-1233 & \ok & result.subjects![0].campiSpecifici['DenominazioneOrganizzazione'] deve essere 'Azienda SRL' \\
\hline
TU-1234 & \ok & result.subjects![0].campiSpecifici['CodiceFiscale\_PartitaIva'] deve essere 'IT123456789' \\
\hline
TU-1235 & \ok & result.subjects![0].tipo deve essere SubjectType.PAI \\
\hline
TU-1236 & \ok & result.subjects![0].campiSpecifici['IPAAmm'] deve essere 'IPA-111' \\
\hline
TU-1237 & \ok & result.subjects![0].tipo deve essere SubjectType.PF \\
\hline
TU-1238 & \ok & result.subjects![0].campiSpecifici['Nome'] deve essere 'Giulia' \\
\hline
TU-1239 & \ok & result.attachments!.numero deve essere 1 \\
\hline
TU-1240 & \ok & result.attachments!.allegati deve avere lunghezza 1 \\
\hline
TU-1241 & \ok & result.attachments!.allegati![0].id deve essere 'All-01' \\
\hline
TU-1242 & \ok & result.attachments!.allegati![0].descrizione deve essere 'Planimetria' \\
\hline
TU-1243 & \ok & result.changeTracking!.tipo deve essere 'Aggiornamento' \\
\hline
TU-1244 & \ok & result.changeTracking!.soggetto deve essere 'Anna Verdi' \\
\hline
TU-1245 & \ok & result.changeTracking!.data deve essere '2023-10-10' \\
\hline
TU-1246 & \ok & result.changeTracking!.idVersionePrecedente deve essere 'v1.0' \\
\hline
TU-1247 & \ok & result.customMetadata deve essere uguale a array containing [ { nome: 'SistemaSorgente', valore: 'ExtApp' }, { nome: 'FieldNotPresentInXSD', valore: '1234' }, ] \\
\hline
TU-1248 & \ok & mockGateway.invoke deve essere stato chiamato con IpcChannels.BROWSE\_GET\_DOCUMENT\_BY\_ID, 456, null \\
\hline
TU-1249 & \ok & mockCache.set deve essere stato chiamato con 'document:456', any Object, 300000 \\
\hline
TU-1250 & \ok & facade.getState()().detail deve essere truthy \\
\hline
TU-1251 & \ok & mockGateway.invoke deve essere stato chiamato alla posizione 1 con IpcChannels.BROWSE\_GET\_FILE\_BY\_DOCUMENT, 456, null \\
\hline
TU-1252 & \ok & mockGateway.invoke deve essere stato chiamato alla posizione 2 con IpcChannels.BROWSE\_GET\_FILE\_BUFFER\_BY\_ID, 999, null \\
\hline
TU-1253 & \ok & mockGateway.invoke non deve essere stato chiamato con IpcChannels.BROWSE\_GET\_FILE\_BY\_ID, 456, null \\
\hline
TU-1254 & \ok & mockCache.set deve essere stato chiamato con 'blob:456', mockBlob, 60000 \\
\hline
TU-1255 & \ok & result deve essere uguale a mockBlob \\
\hline
TU-1256 & \ok & facade.getFileBlob('456') deve essere uguale a mappedError \\
\hline
TU-1257 & \ok & mockTelemetry.trackError deve essere stato chiamato con mappedError \\
\hline
TU-1258 & \ok & mockGateway.invoke non deve essere stato chiamato \\
\hline
TU-1259 & \ok & facade.getState()().detail deve essere uguale a mockDocumentData \\
\hline
TU-1260 & \ok & facade.getState()().error deve essere uguale a mappedError \\
\hline
TU-1261 & \ok & facade.getState()().loading deve essere false \\
\hline
TU-1262 & \ok & mockTelemetry.trackError deve essere stato chiamato con mappedError \\
\hline
TU-1263 & \ok & mockGateway.invoke deve essere stato chiamato 2 volte \\
\hline
TU-1264 & \ok & facade.getState()().detail?.mimeType deve essere MimeType.PDF \\
\hline
TU-1265 & \ok & mockGateway.invoke non deve essere stato chiamato \\
\hline
TU-1266 & \ok & result deve essere uguale a mockBlob \\
\hline
TU-1267 & \ok & facade.getFileBlob('abc') deve essere uguale a mappedError \\
\hline
TU-1268 & \ok & mockErrorHandler.handle deve essere stato chiamato \\
\hline
TU-1269 & \ok & facade.getFileBlob('456') deve essere uguale a mappedError \\
\hline
TU-1270 & \ok & state.phase() deve essere ExportPhase.IDLE \\
\hline
TU-1271 & \ok & state.loading() deve essere false \\
\hline
TU-1272 & \ok & state.progress() deve essere 0 \\
\hline
TU-1273 & \ok & state.phase() deve essere ExportPhase.PROCESSING \\
\hline
TU-1274 & \ok & state.loading() deve essere true \\
\hline
TU-1275 & \ok & state.outputContext() deve essere OutputContext.SINGLE\_EXPORT \\
\hline
TU-1276 & \ok & state.progress() deve essere 0 \\
\hline
TU-1277 & \ok & state.progress() deve essere 46 \\
\hline
TU-1278 & \ok & state.progress() deve essere 100 \\
\hline
TU-1279 & \ok & state.phase() deve essere ExportPhase.SUCCESS \\
\hline
TU-1280 & \ok & state.result() deve essere uguale a mockResult \\
\hline
TU-1281 & \ok & state.loading() deve essere false \\
\hline
TU-1282 & \ok & state.progress() deve essere 100 \\
\hline
TU-1283 & \ok & state.phase() deve essere ExportPhase.ERROR \\
\hline
TU-1284 & \ok & state.error() deve essere uguale a mockError \\
\hline
TU-1285 & \ok & state.loading() deve essere false \\
\hline
TU-1286 & \ok & state.phase() deve essere ExportPhase.IDLE \\
\hline
TU-1287 & \ok & state.loading() deve essere false \\
\hline
TU-1288 & \ok & state.outputContext() deve essere null \\
\hline
TU-1289 & \ok & result.success deve essere false \\
\hline
TU-1290 & \ok & result.errorCode deve essere 'BRIDGE\_UNAVAILABLE' \\
\hline
TU-1291 & \ok & result deve essere uguale a { canceled: false, results: [] } \\
\hline
TU-1292 & \ok & result deve essere null \\
\hline
TU-1293 & \ok & result deve essere uguale a [] \\
\hline
TU-1294 & \ok & result.success deve essere false \\
\hline
TU-1295 & \ok & result deve essere uguale a { canceled: false, results: [] } \\
\hline
TU-1296 & \ok & () => unsub() non deve lanciare un errore \\
\hline
TU-1297 & \ok & () => unsub() non deve lanciare un errore \\
\hline
TU-1298 & \ok & ipc.invoke deve essere stato chiamato con IpcChannels.FILE\_DOWNLOAD, 42 \\
\hline
TU-1299 & \ok & result deve essere uguale a ipcResult \\
\hline
TU-1300 & \ok & gateway.exportFile(1) deve lanciare 'IPC crash' \\
\hline
TU-1301 & \ok & ipc.invoke deve essere stato chiamato con IpcChannels.FILE\_DOWNLOAD\_MANY, [1, 2, 3] \\
\hline
TU-1302 & \ok & result deve essere uguale a ipcResult \\
\hline
TU-1303 & \ok & result.canceled deve essere true \\
\hline
TU-1304 & \ok & gateway.exportFiles([1]) deve lanciare 'IPC crash' \\
\hline
TU-1305 & \ok & ipc.invoke deve essere stato chiamato con 'browse:get-file-by-id', 7 \\
\hline
TU-1306 & \ok & result deve essere uguale a dto \\
\hline
TU-1307 & \ok & result deve essere null \\
\hline
TU-1308 & \ok & ipc.invoke deve essere stato chiamato con 'browse:get-file-by-document', 10 \\
\hline
TU-1309 & \ok & result deve essere uguale a dtos \\
\hline
TU-1310 & \ok & result deve avere lunghezza 2 \\
\hline
TU-1311 & \ok & result deve essere uguale a [] \\
\hline
TU-1312 & \ok & ipc.invoke deve essere stato chiamato con IpcChannels.FILE\_PRINT, 5 \\
\hline
TU-1313 & \ok & result deve essere uguale a { success: true } \\
\hline
TU-1314 & \ok & result.success deve essere false \\
\hline
TU-1315 & \ok & result.error deve essere 'Stampante non trovata' \\
\hline
TU-1316 & \ok & gateway.printFile(1) deve lanciare 'IPC crash' \\
\hline
TU-1317 & \ok & ipc.invoke deve essere stato chiamato con IpcChannels.FILE\_PRINT\_MANY, [10, 11] \\
\hline
TU-1318 & \ok & result deve essere uguale a ipcResult \\
\hline
TU-1319 & \ok & result.canceled deve essere true \\
\hline
TU-1320 & \ok & gateway.printFiles([1]) deve lanciare 'IPC crash' \\
\hline
TU-1321 & \ok & ipc.on deve essere stato chiamato con IpcChannels.FILE\_DOWNLOAD\_PROGRESS, any Function \\
\hline
TU-1322 & \ok & result deve essere unsub \\
\hline
TU-1323 & \ok & ipc.on deve essere stato chiamato con IpcChannels.FILE\_PRINT\_PROGRESS, any Function \\
\hline
TU-1324 & \ok & result deve essere unsub \\
\hline
TU-1325 & \ok & facade.phase deve essere state.phase \\
\hline
TU-1326 & \ok & facade.outputContext deve essere state.outputContext \\
\hline
TU-1327 & \ok & facade.result deve essere state.result \\
\hline
TU-1328 & \ok & facade.progress deve essere state.progress \\
\hline
TU-1329 & \ok & facade.error deve essere state.error \\
\hline
TU-1330 & \ok & facade.loading deve essere state.loading \\
\hline
TU-1331 & \ok & facade.queue deve essere state.queue \\
\hline
TU-1332 & \ok & state.reset deve essere stato chiamato 1 volte \\
\hline
TU-1333 & \ok & state.setProcessing deve essere stato chiamato con OutputContext.SINGLE\_EXPORT \\
\hline
TU-1334 & \ok & state.reset deve essere stato chiamato 1 volte \\
\hline
TU-1335 & \ok & state.setSuccess non deve essere stato chiamato \\
\hline
TU-1336 & \ok & state.setSuccess deve essere stato chiamato con object containing { outputContext: OutputContext.SINGLE\_EXPORT, successCount: 1, failedCount: 0, } \\
\hline
TU-1337 & \ok & state.setError deve essere stato chiamato 1 volte \\
\hline
TU-1338 & \ok & error.message deve essere 'ERR\_IO' \\
\hline
TU-1339 & \ok & error.recoverable deve essere true \\
\hline
TU-1340 & \ok & error.message deve essere 'IPC crash' \\
\hline
TU-1341 & \ok & error.message deve essere 'Errore sconosciuto' \\
\hline
TU-1342 & \ok & state.setProcessing deve essere stato chiamato con OutputContext.MULTI\_EXPORT \\
\hline
TU-1343 & \ok & state.setError deve essere stato chiamato 1 volte \\
\hline
TU-1344 & \ok & queue deve avere lunghezza 2 \\
\hline
TU-1345 & \ok & queue.map((q: any) => q.filename) deve essere uguale a ['a.pdf', 'c.pdf'] \\
\hline
TU-1346 & \ok & state.reset deve essere stato chiamato 1 volte \\
\hline
TU-1347 & \ok & state.setSuccess non deve essere stato chiamato \\
\hline
TU-1348 & \ok & calls deve essere uguale a [ expect.closeTo(33.33, 1), expect.closeTo(66.67, 1), 100, ] \\
\hline
TU-1349 & \ok & result.successCount deve essere 2 \\
\hline
TU-1350 & \ok & result.failedCount deve essere 1 \\
\hline
TU-1351 & \ok & state.updateQueueItem deve essere stato chiamato con 7, { status: 'done', error: undefined } \\
\hline
TU-1352 & \ok & state.updateQueueItem deve essere stato chiamato con 5, object containing { status: 'error', error: 'IO fail' } \\
\hline
TU-1353 & \ok & result.outputContext deve essere OutputContext.MULTI\_EXPORT \\
\hline
TU-1354 & \ok & state.setError deve essere stato chiamato 1 volte \\
\hline
TU-1355 & \ok & state.setProcessing deve essere stato chiamato con OutputContext.SINGLE\_PRINT \\
\hline
TU-1356 & \ok & gateway.printFile deve essere stato chiamato con 1 \\
\hline
TU-1357 & \ok & result.outputContext deve essere OutputContext.SINGLE\_PRINT \\
\hline
TU-1358 & \ok & result.successCount deve essere 1 \\
\hline
TU-1359 & \ok & state.setUnavailable deve essere stato chiamato 1 volte \\
\hline
TU-1360 & \ok & gateway.printFile non deve essere stato chiamato \\
\hline
TU-1361 & \ok & error.message deve contenere 'report.xlsx' \\
\hline
TU-1362 & \ok & error.recoverable deve essere false \\
\hline
TU-1363 & \ok & state.setError deve essere stato chiamato 1 volte \\
\hline
TU-1364 & \ok & state.setError deve essere stato chiamato 1 volte \\
\hline
TU-1365 & \ok & state.setSuccess deve essere stato chiamato \\
\hline
TU-1366 & \ok & state.setUnavailable non deve essere stato chiamato \\
\hline
TU-1367 & \ok & state.setSuccess deve essere stato chiamato \\
\hline
TU-1368 & \ok & gateway.printFiles deve essere stato chiamato 1 volte \\
\hline
TU-1369 & \ok & gateway.printFiles deve essere stato chiamato con [1, 2] \\
\hline
TU-1370 & \ok & gateway.printFiles deve essere stato chiamato con [3] \\
\hline
TU-1371 & \ok & state.setUnavailable deve essere stato chiamato 1 volte \\
\hline
TU-1372 & \ok & gateway.printFiles non deve essere stato chiamato \\
\hline
TU-1373 & \ok & state.setSuccess deve essere stato chiamato 1 volte \\
\hline
TU-1374 & \ok & state.reset deve essere stato chiamato 1 volte \\
\hline
TU-1375 & \ok & state.setSuccess non deve essere stato chiamato \\
\hline
TU-1376 & \ok & state.setError deve essere stato chiamato 1 volte \\
\hline
TU-1377 & \ok & pctElement.textContent deve contenere '45\%' \\
\hline
TU-1378 & \ok & labelElement.textContent deve essere 'Generazione PDF…' \\
\hline
TU-1379 & \ok & labelElement.textContent deve essere 'Stampa in corso…' \\
\hline
TU-1380 & \ok & fillElement.style.width deve essere '75\%' \\
\hline
TU-1381 & \ok & fixture.nativeElement.querySelector('.result-box') deve essere null \\
\hline
TU-1382 & \ok & component.phase() deve essere ExportPhase.IDLE \\
\hline
TU-1383 & \ok & component.result() deve essere null \\
\hline
TU-1384 & \ok & component.error() deve essere null \\
\hline
TU-1385 & \ok & fixture.nativeElement.querySelector('.result-success') non deve essere null \\
\hline
TU-1386 & \ok & fixture.nativeElement.querySelector('.result-error') deve essere null \\
\hline
TU-1387 & \ok & fixture.nativeElement.querySelector('.result-warning') deve essere null \\
\hline
TU-1388 & \ok & box.getAttribute('role') deve essere 'status' \\
\hline
TU-1389 & \ok & box.getAttribute('aria-live') deve essere 'polite' \\
\hline
TU-1390 & \ok & fixture.nativeElement.querySelector('.result-sub') deve essere null \\
\hline
TU-1391 & \ok & sub non deve essere null \\
\hline
TU-1392 & \ok & sub.textContent deve contenere '3 elementi non salvati' \\
\hline
TU-1393 & \ok & fixture.nativeElement.querySelector('.result-success') deve essere null \\
\hline
TU-1394 & \ok & component.successMessage() deve essere '' \\
\hline
TU-1395 & \ok & component.successMessage() deve essere msg \\
\hline
TU-1396 & \ok & component.successMessage() deve essere 'Operazione completata' \\
\hline
TU-1397 & \ok & box.getAttribute('aria-label') deve essere 'Documento salvato con successo' \\
\hline
TU-1398 & \ok & fixture.nativeElement.querySelector('.result-error') non deve essere null \\
\hline
TU-1399 & \ok & fixture.nativeElement.querySelector('.result-success') deve essere null \\
\hline
TU-1400 & \ok & fixture.nativeElement.querySelector('.result-warning') deve essere null \\
\hline
TU-1401 & \ok & box.getAttribute('role') deve essere 'alert' \\
\hline
TU-1402 & \ok & box.getAttribute('aria-live') deve essere 'assertive' \\
\hline
TU-1403 & \ok & title.textContent.trim() deve essere 'Qualcosa è andato storto' \\
\hline
TU-1404 & \ok & box.getAttribute('aria-label') deve essere 'Errore: Qualcosa è andato storto' \\
\hline
TU-1405 & \ok & fixture.nativeElement.querySelector('.retry-btn') deve essere null \\
\hline
TU-1406 & \ok & btn non deve essere null \\
\hline
TU-1407 & \ok & btn.getAttribute('aria-label') deve essere "Riprova l'operazione" \\
\hline
TU-1408 & \ok & fixture.nativeElement.querySelector('.result-error') deve essere null \\
\hline
TU-1409 & \ok & spy deve essere stato chiamato 1 volte \\
\hline
TU-1410 & \ok & emitted[0] deve essere undefined \\
\hline
TU-1411 & \ok & fixture.nativeElement.querySelector('.result-warning') non deve essere null \\
\hline
TU-1412 & \ok & fixture.nativeElement.querySelector('.result-success') deve essere null \\
\hline
TU-1413 & \ok & fixture.nativeElement.querySelector('.result-error') deve essere null \\
\hline
TU-1414 & \ok & box.getAttribute('role') deve essere 'status' \\
\hline
TU-1415 & \ok & box.getAttribute('aria-live') deve essere 'polite' \\
\hline
TU-1416 & \ok & title.textContent.trim() deve essere 'Funzione non disponibile' \\
\hline
TU-1417 & \ok & box.getAttribute('aria-label') deve essere 'Operazione non disponibile: Funzione non disponibile' \\
\hline
TU-1418 & \ok & fixture.nativeElement.querySelector('.retry-btn') deve essere null \\
\hline
TU-1419 & \ok & fixture.nativeElement.querySelector('.result-warning') deve essere null \\
\hline
TU-1420 & \ok & sub.getAttribute('role') deve essere 'alert' \\
\hline
TU-1421 & \ok & sub.getAttribute('aria-live') deve essere 'assertive' \\
\hline
TU-1422 & \ok & box.getAttribute('aria-atomic') deve essere 'true' \\
\hline
TU-1423 & \ok & mockGateway.getFilesByDocumentId deve essere stato chiamato con 42 \\
\hline
TU-1424 & \ok & mockGateway.getFilesByDocumentId non deve essere stato chiamato \\
\hline
TU-1425 & \ok & mockFacade.reset deve essere stato chiamato \\
\hline
TU-1426 & \ok & fixture.debugElement.query(By.css('app-export-progress')) non deve essere null \\
\hline
TU-1427 & \ok & fixture.debugElement.query(By.css('app-export-progress')) deve essere null \\
\hline
TU-1428 & \ok & fixture.debugElement.query(By.css('.queue-container')) non deve essere null \\
\hline
TU-1429 & \ok & fixture.debugElement.query(By.css('.queue-container')) deve essere null \\
\hline
TU-1430 & \ok & fixture.debugElement.query(By.css('app-export-result')) non deve essere null \\
\hline
TU-1431 & \ok & fixture.debugElement.query(By.css('app-export-result')) deve essere null \\
\hline
TU-1432 & \ok & fixture.debugElement.query(By.css('app-export-result')) deve essere null \\
\hline
TU-1433 & \ok & mockFacade.reset deve essere stato chiamato \\
\hline
TU-1434 & \ok & fixture.debugElement.query(By.css('app-export-result')) deve essere null \\
\hline
TU-1435 & \ok & mockFacade.exportFile deve essere stato chiamato con 7 \\
\hline
TU-1436 & \ok & component.showDownloadSelector() deve essere false \\
\hline
TU-1437 & \ok & component.showDownloadSelector() deve essere true \\
\hline
TU-1438 & \ok & component.selectedDownloadIds().has(1) deve essere true \\
\hline
TU-1439 & \ok & component.selectedDownloadIds().has(2) deve essere true \\
\hline
TU-1440 & \ok & component.selectedDownloadIds().has(1) deve essere false \\
\hline
TU-1441 & \ok & component.selectedDownloadIds().has(1) deve essere true \\
\hline
TU-1442 & \ok & mockFacade.exportFiles deve essere stato chiamato con [1, 2] \\
\hline
TU-1443 & \ok & component.showDownloadSelector() deve essere false \\
\hline
TU-1444 & \ok & mockFacade.exportFile deve essere stato chiamato con 1 \\
\hline
TU-1445 & \ok & mockFacade.exportFiles non deve essere stato chiamato \\
\hline
TU-1446 & \ok & mockFacade.exportFile non deve essere stato chiamato \\
\hline
TU-1447 & \ok & mockFacade.exportFiles non deve essere stato chiamato \\
\hline
TU-1448 & \ok & component.showDownloadSelector() deve essere false \\
\hline
TU-1449 & \ok & component.downloadableFiles() deve avere lunghezza 0 \\
\hline
TU-1450 & \ok & component.selectedDownloadIds().size deve essere 0 \\
\hline
TU-1451 & \ok & component.showDownloadSelector() deve essere false \\
\hline
TU-1452 & \ok & mockFacade.printDocument deve essere stato chiamato con 3 \\
\hline
TU-1453 & \ok & component.showPrintSelector() deve essere false \\
\hline
TU-1454 & \ok & mockFacade.printDocument non deve essere stato chiamato \\
\hline
TU-1455 & \ok & component.showPrintSelector() deve essere false \\
\hline
TU-1456 & \ok & component.showPrintSelector() deve essere true \\
\hline
TU-1457 & \ok & component.selectedPrintIds().has(10) deve essere true \\
\hline
TU-1458 & \ok & component.selectedPrintIds().has(11) deve essere true \\
\hline
TU-1459 & \ok & component.selectedPrintIds().has(10) deve essere false \\
\hline
TU-1460 & \ok & component.selectedPrintIds().has(10) deve essere true \\
\hline
TU-1461 & \ok & mockFacade.printDocuments deve essere stato chiamato con [10, 11] \\
\hline
TU-1462 & \ok & component.showPrintSelector() deve essere false \\
\hline
TU-1463 & \ok & mockFacade.printDocuments non deve essere stato chiamato \\
\hline
TU-1464 & \ok & component.showPrintSelector() deve essere false \\
\hline
TU-1465 & \ok & component.printableFiles() deve avere lunghezza 0 \\
\hline
TU-1466 & \ok & component.selectedPrintIds().size deve essere 0 \\
\hline
TU-1467 & \ok & component.selectedDownloadCount() deve essere 2 \\
\hline
TU-1468 & \ok & component.selectedDownloadCount() deve essere 1 \\
\hline
TU-1469 & \ok & component.selectedPrintCount() deve essere 2 \\
\hline
TU-1470 & \ok & component.selectedPrintCount() deve essere 1 \\
\hline
TU-1471 & \ok & mockCache.get deve essere stato chiamato con 'node-fallback:DIP:1' \\
\hline
TU-1472 & \ok & mockGateway.invoke non deve essere stato chiamato \\
\hline
TU-1473 & \ok & facade.getState()().detail deve essere uguale a cached \\
\hline
TU-1474 & \ok & mockGateway.invoke deve essere stato chiamato con IpcChannels.BROWSE\_GET\_DOCUMENT\_CLASS\_BY\_ID, 22, null \\
\hline
TU-1475 & \ok & mockGateway.invoke deve essere stato chiamato con IpcChannels.BROWSE\_GET\_PROCESS\_BY\_DOCUMENT\_CLASS, 22, null \\
\hline
TU-1476 & \ok & facade.getState()().detail?.title deve essere 'Classe Contratti' \\
\hline
TU-1477 & \ok & facade.getState()().detail?.relatedSection?.items deve essere uguale a [ { itemType: 'PROCESS', itemId: '31', label: 'Processo Contratti', description: 'UUID: PROC-31 - Stato: UNKNOWN', }, ] \\
\hline
TU-1478 & \ok & facade .getState()() .detail?.fields.some((field) => field.label.toLowerCase().includes('id interno')) deve essere false \\
\hline
TU-1479 & \ok & mockCache.set deve essere stato chiamato con 'node-fallback:DOCUMENT\_CLASS:22', any Object, 300000 \\
\hline
TU-1480 & \ok & facade.getState()().detail?.title deve essere 'Allegato\_1.pdf' \\
\hline
TU-1481 & \ok & mockErrorHandler.handle deve essere stato chiamato con rawError \\
\hline
TU-1482 & \ok & facade.getState()().error deve essere uguale a mappedError \\
\hline
TU-1483 & \ok & mockTelemetry.trackError deve essere stato chiamato con mappedError \\
\hline
TU-1484 & \ok & mockGateway.invoke non deve essere stato chiamato \\
\hline
TU-1485 & \ok & mockErrorHandler.handle deve essere stato chiamato \\
\hline
TU-1486 & \ok & facade.getState()().error deve essere uguale a mockError \\
\hline
TU-1487 & \ok & mockGateway.invoke deve essere stato chiamato con IpcChannels.BROWSE\_GET\_DIP\_BY\_ID, 5, null \\
\hline
TU-1488 & \ok & facade.getState()().detail?.type deve essere 'DIP' \\
\hline
TU-1489 & \ok & facade.getState()().detail?.fields deve contenere un elemento uguale a object containing { label: 'UUID', value: 'DIP-UUID-123' } \\
\hline
TU-1490 & \ok & mockErrorHandler.handle deve essere stato chiamato \\
\hline
TU-1491 & \ok & facade.getState()().error deve essere uguale a mockError \\
\hline
TU-1492 & \ok & detail?.title deve essere 'Classe Documentale' \\
\hline
TU-1493 & \ok & items.length deve essere 3 \\
\hline
TU-1494 & \ok & items[0].label deve essere 'Proc Alpha' \\
\hline
TU-1495 & \ok & items[1].label deve essere 'Agg Beta' \\
\hline
TU-1496 & \ok & items[2].label deve essere 'PROC-UUID-X' \\
\hline
TU-1497 & \ok & detail?.relatedSection?.items deve essere uguale a [] \\
\hline
TU-1498 & \ok & detail?.hint deve essere 'I metadati di processo verranno mostrati qui quando disponibili.' \\
\hline
TU-1499 & \ok & component deve essere truthy \\
\hline
TU-1500 & \ok & htmlElement.querySelector('.title')?.textContent deve contenere 'Documento di Test' \\
\hline
TU-1501 & \ok & htmlElement.querySelector('.badge')?.textContent deve contenere 'PDF' \\
\hline
TU-1502 & \ok & component deve essere truthy \\
\hline
TU-1503 & \ok & fixture.nativeElement.textContent deve contenere 'Documento' \\
\hline
TU-1504 & \ok & fixture.nativeElement.textContent deve contenere 'Processo' \\
\hline
TU-1505 & \ok & fixture.nativeElement.textContent deve contenere 'Classe Documentale' \\
\hline
TU-1506 & \ok & fixture.nativeElement.textContent deve contenere 'Fascicolo' \\
\hline
TU-1507 & \ok & mockIntegrityFacade.verifyItem deve essere stato chiamato con '123', 'DOCUMENT' \\
\hline
TU-1508 & \ok & component.manualStatus() deve essere 'VALID' \\
\hline
TU-1509 & \ok & spyEmit deve essere stato chiamato \\
\hline
TU-1510 & \ok & mockIntegrityFacade.verifyItem deve essere stato chiamato con '456', 'PROCESS' \\
\hline
TU-1511 & \ok & component.manualStatus() deve essere 'UNKNOWN' \\
\hline
TU-1512 & \ok & spyEmit deve essere stato chiamato \\
\hline
TU-1513 & \ok & component.verificationStatus() deve essere 'INVALID' \\
\hline
TU-1514 & \ok & component.verificationStatus() deve essere 'VALID' \\
\hline
TU-1515 & \ok & mockOutputFacade.exportPdf deve essere stato chiamato con { documentId: '789', tipo: 'PROCESS' } \\
\hline
TU-1516 & \ok & mockOutputFacade.print deve essere stato chiamato con { documentId: '789', tipo: 'PROCESS' } \\
\hline
TU-1517 & \ok & component.manualStatus() deve essere 'VALID' \\
\hline
TU-1518 & \ok & component.manualStatus() deve essere null \\
\hline
TU-1519 & \ok & mockAggregateFacade.loadAggregate deve essere stato chiamato con '123' \\
\hline
TU-1520 & \ok & mockDocumentFacade.loadDocument non deve essere stato chiamato \\
\hline
TU-1521 & \ok & mockFallbackFacade.loadNode non deve essere stato chiamato \\
\hline
TU-1522 & \ok & component.pageTitle() deve essere 'Fascicolo affare' \\
\hline
TU-1523 & \ok & component.isLoading() deve essere false \\
\hline
TU-1524 & \ok & mockDocumentFacade.loadDocument deve essere stato chiamato con '456' \\
\hline
TU-1525 & \ok & mockAggregateFacade.loadAggregate non deve essere stato chiamato \\
\hline
TU-1526 & \ok & mockProcessFacade.loadProcess non deve essere stato chiamato \\
\hline
TU-1527 & \ok & mockFallbackFacade.loadNode non deve essere stato chiamato \\
\hline
TU-1528 & \ok & component.pageTitle() deve essere 'delibera.pdf' \\
\hline
TU-1529 & \ok & mockProcessFacade.loadProcess deve essere stato chiamato con '31' \\
\hline
TU-1530 & \ok & mockAggregateFacade.loadAggregate non deve essere stato chiamato \\
\hline
TU-1531 & \ok & mockDocumentFacade.loadDocument non deve essere stato chiamato \\
\hline
TU-1532 & \ok & mockFallbackFacade.loadNode non deve essere stato chiamato \\
\hline
TU-1533 & \ok & component.pageTitle() deve essere 'PROC-31' \\
\hline
TU-1534 & \ok & compiled.querySelector('.process-sidebar') deve essere truthy \\
\hline
TU-1535 & \ok & compiled.querySelector('.process-main-content') deve essere truthy \\
\hline
TU-1536 & \ok & mockFallbackFacade.loadNode deve essere stato chiamato con 'DIP', '1' \\
\hline
TU-1537 & \ok & component.pageTitle() deve essere 'DIP' \\
\hline
TU-1538 & \ok & fixture.nativeElement.querySelector('app-node-fallback-panel') deve essere truthy \\
\hline
TU-1539 & \ok & routerMock.navigate deve essere stato chiamato con ['/detail', 'PROCESS', '31'] \\
\hline
TU-1540 & \ok & component.isLoading() deve essere true \\
\hline
TU-1541 & \ok & compiled.querySelector('.spinner-overlay') deve essere truthy \\
\hline
TU-1542 & \ok & component.currentError() deve essere uguale a fintoErrore \\
\hline
TU-1543 & \ok & compiled.querySelector('app-error-dialog') deve essere truthy \\
\hline
TU-1544 & \ok & mockAggregateFacade.loadAggregate deve essere stato chiamato con 'ERR-1' \\
\hline
TU-1545 & \ok & component.pageTitle() deve essere '' \\
\hline
TU-1546 & \ok & component deve essere truthy \\
\hline
TU-1547 & \ok & overlay deve essere truthy \\
\hline
TU-1548 & \ok & service deve essere truthy \\
\hline
TU-1549 & \ok & mapDipNodeTypeToDetailItemType('dip') deve essere 'DIP' \\
\hline
TU-1550 & \ok & mapDipNodeTypeToDetailItemType('documentClass') deve essere 'DOCUMENT\_CLASS' \\
\hline
TU-1551 & \ok & mapDipNodeTypeToDetailItemType('process') deve essere 'PROCESS' \\
\hline
TU-1552 & \ok & mapDipNodeTypeToDetailItemType('document') deve essere 'DOCUMENT' \\
\hline
TU-1553 & \ok & mapDipNodeTypeToDetailItemType('file') deve essere 'FILE' \\
\hline
TU-1554 & \ok & mapSearchResultTypeToDetailItemType('AGGREGAZIONE DOCUMENTALE') deve essere 'AGGREGATE' \\
\hline
TU-1555 & \ok & mapSearchResultTypeToDetailItemType('documento\_informatico') deve essere 'DOCUMENT' \\
\hline
TU-1556 & \ok & mapSearchResultTypeToDetailItemType('documento amministrativo informatico') deve essere 'DOCUMENT' \\
\hline
TU-1557 & \ok & mapSearchResultTypeToDetailItemType('processo') deve essere 'PROCESS' \\
\hline
TU-1558 & \ok & mapSearchResultTypeToDetailItemType('classe') deve essere 'DOCUMENT\_CLASS' \\
\hline
TU-1559 & \ok & mapSearchResultTypeToDetailItemType('tipo sconosciuto') deve essere null \\
\hline
TU-1560 & \ok & isRichDetailRouteItemType('AGGREGATE') deve essere true \\
\hline
TU-1561 & \ok & isRichDetailRouteItemType('DOCUMENT') deve essere true \\
\hline
TU-1562 & \ok & isRichDetailRouteItemType('PROCESS') deve essere true \\
\hline
TU-1563 & \ok & isRichDetailRouteItemType('DIP') deve essere false \\
\hline
TU-1564 & \ok & buildDetailRoute('DOCUMENT\_CLASS', 77) deve essere uguale a ['/detail', 'DOCUMENT\_CLASS', '77'] \\
\hline
TU-1565 & \ok & facade.getState()().phase deve essere 'idle' \\
\hline
TU-1566 & \ok & facade.getState()().rootNodes deve avere lunghezza 0 \\
\hline
TU-1567 & \ok & facade.getState()().nodeCache.size deve essere 0 \\
\hline
TU-1568 & \ok & facade.getState()().loadingNodeIds.size deve essere 0 \\
\hline
TU-1569 & \ok & facade.getState()().nodeChildrenErrors.size deve essere 0 \\
\hline
TU-1570 & \ok & phaseWhileLoading deve essere 'loading' \\
\hline
TU-1571 & \ok & facade.getState()().phase deve essere 'ready' \\
\hline
TU-1572 & \ok & facade.getState()().rootNodes deve essere uguale a [rootNode] \\
\hline
TU-1573 & \ok & gatewayMock.getRootDip deve essere stato chiamato con 42 \\
\hline
TU-1574 & \ok & facade.getState()().rootError deve essere definito \\
\hline
TU-1575 & \ok & facade.getState()().rootError deve essere undefined \\
\hline
TU-1576 & \ok & facade.getState()().phase deve essere 'idle' \\
\hline
TU-1577 & \ok & facade.getState()().rootError deve essere uguale a error \\
\hline
TU-1578 & \ok & facade.getState()().rootNodes deve avere lunghezza 0 \\
\hline
TU-1579 & \ok & gatewayMock.getChildren non deve essere stato chiamato \\
\hline
TU-1580 & \ok & facade.getState()().phase deve essere stateBefore.phase \\
\hline
TU-1581 & \ok & loadingDuringCall deve essere true \\
\hline
TU-1582 & \ok & facade.getState()().loadingNodeIds.has(keyFor(parent)) deve essere false \\
\hline
TU-1583 & \ok & facade.getState()().nodeCache.get(keyFor(parent)) deve essere uguale a children \\
\hline
TU-1584 & \ok & facade.getState()().nodeChildrenErrors.has(keyFor(parent)) deve essere true \\
\hline
TU-1585 & \ok & facade.getState()().nodeChildrenErrors.has(keyFor(parent)) deve essere false \\
\hline
TU-1586 & \ok & facade.getState()().nodeCache.get(keyFor(child)) deve essere uguale a [grandchild] \\
\hline
TU-1587 & \ok & facade.getState()().loadingNodeIds.has(keyFor(parent)) deve essere false \\
\hline
TU-1588 & \ok & facade.getState()().nodeChildrenErrors.get(keyFor(parent)) deve essere uguale a error \\
\hline
TU-1589 & \ok & facade.getState()().phase deve essere 'ready' \\
\hline
TU-1590 & \ok & facade.getState()().nodeChildrenErrors.has(keyFor(child2)) deve essere true \\
\hline
TU-1591 & \ok & facade.getState()().nodeChildrenErrors.has(keyFor(root)) deve essere false \\
\hline
TU-1592 & \ok & mockBridge.invoke deve essere stato chiamato con IpcChannels.BROWSE\_GET\_DOCUMENT\_CLASS\_BY\_DIP\_ID, 1 \\
\hline
TU-1593 & \ok & result deve essere uguale a [] \\
\hline
TU-1594 & \ok & result deve essere uguale a [ { id: 10, name: 'Classe Contratti', type: 'documentClass', hasChildren: true, } ] \\
\hline
TU-1595 & \ok & result deve essere uguale a [ { id: 1234567, name: 'Classe Documentale', type: 'documentClass', hasChildren: true, } ] \\
\hline
TU-1596 & \ok & result.map((node) => node.name) deve essere uguale a [ 'Processo Contratti', 'Processo', ] \\
\hline
TU-1597 & \ok & result[0].name deve essere 'Allegato\_1.pdf' \\
\hline
TU-1598 & \ok & result[0].type deve essere 'file' \\
\hline
TU-1599 & \ok & result[0].hasChildren deve essere false \\
\hline
TU-1600 & \ok & result deve essere cached \\
\hline
TU-1601 & \ok & mockBridge.invoke non deve essere stato chiamato \\
\hline
TU-1602 & \ok & mockCache.set deve essere stato chiamato \\
\hline
TU-1603 & \ok & gateway.getChildren(parent) deve lanciare un errore \\
\hline
TU-1604 & \ok & mockErrorHandler.handle deve essere stato chiamato con error \\
\hline
TU-1605 & \ok & component deve essere truthy \\
\hline
TU-1606 & \ok & compiled.textContent deve contenere mockNode.name \\
\hline
TU-1607 & \ok & component.toggle.emit deve essere stato chiamato con mockNode \\
\hline
TU-1608 & \ok & component.nodeSelected.emit deve essere stato chiamato con mockNode \\
\hline
TU-1609 & \ok & spy deve essere stato chiamato \\
\hline
TU-1610 & \ok & component.nodeSelected.emit deve essere stato chiamato con mockNode \\
\hline
TU-1611 & \ok & component.toggle.emit deve essere stato chiamato con mockNode \\
\hline
TU-1612 & \ok & component.nodeSelected.emit deve essere stato chiamato con mockNode \\
\hline
TU-1613 & \ok & component.toggle.emit non deve essere stato chiamato \\
\hline
TU-1614 & \ok & component.nodeSelected.emit deve essere stato chiamato con mockNode \\
\hline
TU-1615 & \ok & component.toggle.emit non deve essere stato chiamato \\
\hline
TU-1616 & \ok & button deve essere truthy \\
\hline
TU-1617 & \ok & button deve essere falsy \\
\hline
TU-1618 & \ok & spinner deve essere truthy \\
\hline
TU-1619 & \ok & spinner deve essere falsy \\
\hline
TU-1620 & \ok & error deve essere falsy \\
\hline
TU-1621 & \ok & dot deve essere truthy \\
\hline
TU-1622 & \ok & dot.classList.contains('node-dot-closed') deve essere true \\
\hline
TU-1623 & \ok & dot deve essere truthy \\
\hline
TU-1624 & \ok & dot.classList.contains('node-dot-open') deve essere true \\
\hline
TU-1625 & \ok & component deve essere truthy \\
\hline
TU-1626 & \ok & el.textContent deve contenere 'Errore test' \\
\hline
TU-1627 & \ok & button deve essere truthy \\
\hline
TU-1628 & \ok & button.nativeElement.textContent deve contenere 'Riprova' \\
\hline
TU-1629 & \ok & button deve essere null \\
\hline
TU-1630 & \ok & spy deve essere stato chiamato \\
\hline
TU-1631 & \ok & component deve essere truthy \\
\hline
TU-1632 & \ok & fixture.nativeElement.textContent deve contenere 'Caricamento' \\
\hline
TU-1633 & \ok & routerMock.navigate deve essere stato chiamato con ['/detail', 'DOCUMENT', '1'] \\
\hline
TU-1634 & \ok & routerMock.navigate deve essere stato chiamato con ['/detail', 'PROCESS', '42'] \\
\hline
TU-1635 & \ok & routerMock.navigate deve essere stato chiamato con ['/detail', 'DIP', '7'] \\
\hline
TU-1636 & \ok & routerMock.navigate deve essere stato chiamato con ['/detail', 'DOCUMENT\_CLASS', '9'] \\
\hline
TU-1637 & \ok & routerMock.navigate deve essere stato chiamato con ['/detail', 'FILE', '11'] \\
\hline
TU-1638 & \ok & component deve essere truthy \\
\hline
TU-1639 & \ok & component.rootNodes deve essere uguale a nodes \\
\hline
TU-1640 & \ok & component.rootNodes deve avere lunghezza 0 \\
\hline
TU-1641 & \ok & component.rootNodes deve avere lunghezza 0 \\
\hline
TU-1642 & \ok & component.rootNodes deve essere uguale a nodes \\
\hline
TU-1643 & \ok & component.state() deve corrispondere a { phase: 'ready' } \\
\hline
TU-1644 & \ok & component.state().phase deve essere 'loading' \\
\hline
TU-1645 & \ok & component.state().phase deve essere 'idle' \\
\hline
TU-1646 & \ok & component.state().phase deve essere 'ready' \\
\hline
TU-1647 & \ok & routerMock.navigate deve essere stato chiamato con ['/detail', 'DIP', '3'] \\
\hline
TU-1648 & \ok & routerMock.navigate deve essere stato chiamato con ['/detail', 'DOCUMENT\_CLASS', '12'] \\
\hline
TU-1649 & \ok & routerMock.navigate deve essere stato chiamato con ['/detail', 'FILE', '44'] \\
\hline
TU-1650 & \ok & component deve essere truthy \\
\hline
TU-1651 & \ok & component.flatNodes()[0].node deve essere node \\
\hline
TU-1652 & \ok & ids deve contenere 2 \\
\hline
TU-1653 & \ok & ids non deve contenere 1 \\
\hline
TU-1654 & \ok & component.flatNodes() deve avere lunghezza 0 \\
\hline
TU-1655 & \ok & component.flatNodes()[0].depth deve essere 0 \\
\hline
TU-1656 & \ok & component.flatNodes()[0].isLoading deve essere true \\
\hline
TU-1657 & \ok & component.flatNodes()[0].isLoading deve essere false \\
\hline
TU-1658 & \ok & component.flatNodes()[0].childrenError deve essere error \\
\hline
TU-1659 & \ok & component.flatNodes()[0].childrenError deve essere undefined \\
\hline
TU-1660 & \ok & component.flatNodes() deve avere lunghezza 1 \\
\hline
TU-1661 & \ok & flat deve avere lunghezza 2 \\
\hline
TU-1662 & \ok & flat[1].node.id deve essere 2 \\
\hline
TU-1663 & \ok & flat[1].depth deve essere 1 \\
\hline
TU-1664 & \ok & component.flatNodes() deve avere lunghezza 2 \\
\hline
TU-1665 & \ok & component.flatNodes() deve avere lunghezza 1 \\
\hline
TU-1666 & \ok & component.flatNodes()[0].isExpanded deve essere true \\
\hline
TU-1667 & \ok & component.flatNodes()[0].isExpanded deve essere false \\
\hline
TU-1668 & \ok & component.flatNodes() deve avere lunghezza 2 \\
\hline
TU-1669 & \ok & component.flatNodes()[0].isExpanded deve essere true \\
\hline
TU-1670 & \ok & component.flatNodes()[1].isExpanded deve essere false \\
\hline
TU-1671 & \ok & facadeMock.loadChildren deve essere stato chiamato con child \\
\hline
TU-1672 & \ok & facadeMock.loadChildren deve essere stato chiamato con parent \\
\hline
TU-1673 & \ok & facadeMock.loadChildren non deve essere stato chiamato \\
\hline
TU-1674 & \ok & facadeMock.loadChildren non deve essere stato chiamato \\
\hline
TU-1675 & \ok & component.flatNodes() deve avere lunghezza 1 \\
\hline
TU-1676 & \ok & component.flatNodes() deve avere lunghezza 2 \\
\hline
TU-1677 & \ok & component.flatNodes()[0].isLoading deve essere true \\
\hline
TU-1678 & \ok & spy deve essere stato chiamato con node \\
\hline
TU-1679 & \ok & spy deve essere stato chiamato con nodes[1] \\
\hline
TU-1680 & \ok & spy non deve essere stato chiamato con nodes[0] \\
\hline
TU-1681 & \ok & el.querySelector('[data-testid="process-metadata-card-anagrafica"]') deve essere truthy \\
\hline
TU-1682 & \ok & el.querySelector('[data-testid="process-metadata-card-overview"]') deve essere truthy \\
\hline
TU-1683 & \ok & el.querySelector('[data-testid="process-metadata-card-submission"]') deve essere truthy \\
\hline
TU-1684 & \ok & el.querySelector('[data-testid="process-metadata-card-conservation"]') deve essere truthy \\
\hline
TU-1685 & \ok & el.querySelector('[data-testid="process-metadata-card-custom"]') deve essere null \\
\hline
TU-1686 & \ok & el.querySelector('[data-testid="process-metadata-heading-overview"]')?.textContent?.trim() deve essere 'Contesto del Processo' \\
\hline
TU-1687 & \ok & el.querySelector('[data-testid="process-metadata-row-process-id"]')?.textContent deve contenere '31' \\
\hline
TU-1688 & \ok & el.querySelector('[data-testid="process-metadata-row-process-status"]')?.textContent deve contenere 'TERMINATO' \\
\hline
TU-1689 & \ok & el.querySelector('[data-testid="process-metadata-row-document-class-timestamp"]') ?.textContent deve contenere formatReadableDate('2026-04-08T10:30:00Z') \\
\hline
TU-1690 & \ok & el.querySelector('[data-testid="process-metadata-row-submission-data-inizio"]')?.textContent deve contenere formatReadableDate('2026-04-08') \\
\hline
TU-1691 & \ok & el.querySelector('[data-testid="process-metadata-row-conservation-data-inizio"]') ?.textContent deve contenere formatReadableDate('2026-04-08') \\
\hline
TU-1692 & \ok & el.querySelector('[data-testid="process-metadata-row-overview-oggetto"]')?.textContent deve contenere 'Processo Contratti' \\
\hline
TU-1693 & \ok & el.querySelector('[data-testid="process-metadata-row-conservation-data-fine"]') deve essere truthy \\
\hline
TU-1694 & \ok & el.querySelector('[data-testid="process-metadata-row-conservation-uuid-attivatore"]') deve essere truthy \\
\hline
TU-1695 & \ok & el.querySelector('[data-testid="process-metadata-row-conservation-uuid-terminatore"]') deve essere truthy \\
\hline
TU-1696 & \ok & el.querySelector('[data-testid="process-metadata-row-conservation-canale-attivazione"]') deve essere truthy \\
\hline
TU-1697 & \ok & el.querySelector('[data-testid="process-metadata-row-conservation-canale-terminazione"]') deve essere truthy \\
\hline
TU-1698 & \ok & el.querySelector('[data-testid="process-metadata-row-conservation-stato"]') deve essere truthy \\
\hline
TU-1699 & \ok & el.querySelector('[data-testid="process-metadata-row-submission-data-fine"]') deve essere truthy \\
\hline
TU-1700 & \ok & el.querySelector('[data-testid="process-metadata-row-submission-uuid-attivatore"]') deve essere truthy \\
\hline
TU-1701 & \ok & el.querySelector('[data-testid="process-metadata-row-submission-uuid-terminatore"]') deve essere truthy \\
\hline
TU-1702 & \ok & el.querySelector('[data-testid="process-metadata-row-submission-canale-attivazione"]') deve essere truthy \\
\hline
TU-1703 & \ok & el.querySelector('[data-testid="process-metadata-row-submission-canale-terminazione"]') deve essere truthy \\
\hline
TU-1704 & \ok & el.querySelector('[data-testid="process-metadata-row-submission-stato"]') deve essere truthy \\
\hline
TU-1705 & \ok & el.querySelector('[data-testid="process-metadata-card-overview"]') deve essere null \\
\hline
TU-1706 & \ok & el.querySelector('[data-testid="process-metadata-overview-empty"]') deve essere truthy \\
\hline
TU-1707 & \ok & el.querySelector('[data-testid="optional-field-absent-message"]')?.textContent deve contenere 'Nessun contesto del processo disponibile nei metadati' \\
\hline
TU-1708 & \ok & el.querySelector('[data-testid="process-metadata-row-conservation-data-fine"]') deve essere null \\
\hline
TU-1709 & \ok & el.querySelector('[data-testid="process-metadata-row-conservation-uuid-attivatore"]') deve essere null \\
\hline
TU-1710 & \ok & el.querySelector('[data-testid="process-metadata-row-conservation-uuid-terminatore"]') deve essere null \\
\hline
TU-1711 & \ok & el.querySelector('[data-testid="process-metadata-row-conservation-canale-attivazione"]') deve essere null \\
\hline
TU-1712 & \ok & el.querySelector('[data-testid="process-metadata-row-conservation-canale-terminazione"]') deve essere null \\
\hline
TU-1713 & \ok & el.querySelector('[data-testid="process-metadata-row-conservation-stato"]') deve essere null \\
\hline
TU-1714 & \ok & el.querySelector('[data-testid="process-metadata-row-submission-data-fine"]') deve essere null \\
\hline
TU-1715 & \ok & el.querySelector('[data-testid="process-metadata-row-submission-uuid-attivatore"]') deve essere null \\
\hline
TU-1716 & \ok & el.querySelector('[data-testid="process-metadata-row-submission-uuid-terminatore"]') deve essere null \\
\hline
TU-1717 & \ok & el.querySelector('[data-testid="process-metadata-row-submission-canale-attivazione"]') deve essere null \\
\hline
TU-1718 & \ok & el.querySelector('[data-testid="process-metadata-row-submission-canale-terminazione"]') deve essere null \\
\hline
TU-1719 & \ok & el.querySelector('[data-testid="process-metadata-row-submission-stato"]') deve essere null \\
\hline
TU-1720 & \ok & el.querySelector('[data-testid="process-metadata-row-document-class-name"]')?.textContent deve contenere 'N/A' \\
\hline
TU-1721 & \ok & el.querySelector('[data-testid="process-metadata-row-document-class-uuid"]')?.textContent deve contenere 'N/A' \\
\hline
TU-1722 & \ok & el.querySelector('[data-testid="process-metadata-row-document-class-timestamp"]') ?.textContent deve contenere 'N/A' \\
\hline
TU-1723 & \ok & el.querySelector('[data-testid="process-metadata-card-custom"]') deve essere null \\
\hline
TU-1724 & \ok & el.querySelector('[data-testid="process-metadata-custom-empty"]') deve essere null \\
\hline
TU-1725 & \ok & detail.processId deve essere '31' \\
\hline
TU-1726 & \ok & detail.processUuid deve essere 'PROC-31' \\
\hline
TU-1727 & \ok & detail.integrityStatus deve essere 'VALID' \\
\hline
TU-1728 & \ok & detail.metadata.processId deve essere '31' \\
\hline
TU-1729 & \ok & detail.metadata.processUuid deve essere 'PROC-31' \\
\hline
TU-1730 & \ok & detail.submission.processo deve essere 'PROC-XML-UUID' \\
\hline
TU-1731 & \ok & detail.submission.sessione deve essere 'SUB-SESSION-01' \\
\hline
TU-1732 & \ok & detail.submission.dataInizio deve essere '2025-12-11T16:48:48.087Z' \\
\hline
TU-1733 & \ok & detail.submission.dataFine deve essere '2025-12-11T16:50:39.671Z' \\
\hline
TU-1734 & \ok & detail.submission.uuidAttivatore deve essere 'USR-001' \\
\hline
TU-1735 & \ok & detail.submission.uuidTerminatore deve essere 'USR-001' \\
\hline
TU-1736 & \ok & detail.submission.canaleAttivazione deve essere 'WebGui' \\
\hline
TU-1737 & \ok & detail.submission.canaleTerminazione deve essere 'WebGui' \\
\hline
TU-1738 & \ok & detail.submission.stato deve essere 'TERMINATA' \\
\hline
TU-1739 & \ok & detail.conservation.processo deve essere 'PROC-XML-UUID' \\
\hline
TU-1740 & \ok & detail.conservation.sessione deve essere 'PRES-SESSION-01' \\
\hline
TU-1741 & \ok & detail.conservation.dataInizio deve essere '2025-12-11T16:50:39.671Z' \\
\hline
TU-1742 & \ok & detail.conservation.dataFine deve essere undefined \\
\hline
TU-1743 & \ok & detail.conservation.uuidAttivatore deve essere 'USR-001' \\
\hline
TU-1744 & \ok & detail.conservation.canaleAttivazione deve essere 'Other' \\
\hline
TU-1745 & \ok & detail.conservation.stato deve essere undefined \\
\hline
TU-1746 & \ok & detail.metadata.processStatus deve essere 'CHIUSO' \\
\hline
TU-1747 & \ok & detail.documentClass.id deve essere 22 \\
\hline
TU-1748 & \ok & detail.metadata.documentClassName deve essere 'N/A' \\
\hline
TU-1749 & \ok & detail.submission.processo deve essere 'PRES-31' \\
\hline
TU-1750 & \ok & detail.submission.sessione deve essere 'SESSION-99' \\
\hline
TU-1751 & \ok & detail.submission.dataInizio deve essere '2026-04-01' \\
\hline
TU-1752 & \ok & detail.submission.dataFine deve essere undefined \\
\hline
TU-1753 & \ok & detail.submission.uuidAttivatore deve essere 'USR-ATT-VERS-001' \\
\hline
TU-1754 & \ok & detail.submission.canaleAttivazione deve essere 'PortaleWeb' \\
\hline
TU-1755 & \ok & detail.submission.stato deve essere 'TERMINATA' \\
\hline
TU-1756 & \ok & detail.conservation.processo deve essere 'PRES-31' \\
\hline
TU-1757 & \ok & detail.conservation.sessione deve essere 'SESSION-99' \\
\hline
TU-1758 & \ok & detail.conservation.dataInizio deve essere '2026-04-08' \\
\hline
TU-1759 & \ok & detail.conservation.dataFine deve essere '2026-04-09' \\
\hline
TU-1760 & \ok & detail.conservation.uuidAttivatore deve essere 'USR-ATT-CONS-001' \\
\hline
TU-1761 & \ok & detail.conservation.canaleAttivazione deve essere 'ConservazioneBatch' \\
\hline
TU-1762 & \ok & detail.conservation.stato deve essere 'TERMINATA' \\
\hline
TU-1763 & \ok & detail.metadata.processStatus deve essere 'COMPLETATO' \\
\hline
TU-1764 & \ok & detail.processId deve essere '7' \\
\hline
TU-1765 & \ok & detail.processUuid deve essere 'N/A' \\
\hline
TU-1766 & \ok & detail.integrityStatus deve essere 'UNKNOWN' \\
\hline
TU-1767 & \ok & detail.overview.oggetto deve essere 'N/D Test' \\
\hline
TU-1768 & \ok & detail.conservation.processo deve essere 'N/A' \\
\hline
TU-1769 & \ok & detail.submission.processo deve essere 'N/A' \\
\hline
TU-1770 & \ok & detail.submission.dataInizio deve essere 'N/A' \\
\hline
TU-1771 & \ok & detail.documentClass.id deve essere null \\
\hline
TU-1772 & \ok & detail.documentClass.name deve essere 'N/A' \\
\hline
TU-1773 & \ok & detail.metadata.documentClassUuid deve essere 'N/A' \\
\hline
TU-1774 & \ok & mockCache.get deve essere stato chiamato con 'process:31' \\
\hline
TU-1775 & \ok & mockGateway.invoke non deve essere stato chiamato \\
\hline
TU-1776 & \ok & facade.getState()().detail deve essere uguale a cachedDetail \\
\hline
TU-1777 & \ok & mockGateway.invoke deve essere stato chiamato con IpcChannels.BROWSE\_GET\_PROCESS\_BY\_ID, 31, null \\
\hline
TU-1778 & \ok & mockGateway.invoke deve essere stato chiamato con IpcChannels.BROWSE\_GET\_DOCUMENTS\_BY\_PROCESS, 31, null \\
\hline
TU-1779 & \ok & mockGateway.invoke deve essere stato chiamato con IpcChannels.BROWSE\_GET\_DOCUMENT\_CLASS\_BY\_ID, 22, null \\
\hline
TU-1780 & \ok & facade.getState()().detail?.processUuid deve essere 'PROC-31' \\
\hline
TU-1781 & \ok & facade.getState()().detail?.documentClass.name deve essere 'Classe Contratti' \\
\hline
TU-1782 & \ok & facade.getState()().detail?.metadata.documentClassName deve essere 'Classe Contratti' \\
\hline
TU-1783 & \ok & facade.getState()().detail?.indiceDocumenti.length deve essere 1 \\
\hline
TU-1784 & \ok & mockCache.set deve essere stato chiamato con 'process:31', any Object, 300000 \\
\hline
TU-1785 & \ok & mockErrorHandler.handle deve essere stato chiamato con rawError \\
\hline
TU-1786 & \ok & facade.getState()().error deve essere uguale a mappedError \\
\hline
TU-1787 & \ok & mockTelemetry.trackError deve essere stato chiamato con mappedError \\
\hline
TU-1788 & \ok & result deve essere true \\
\hline
TU-1789 & \ok & mockGateway.invoke deve essere stato chiamato con IpcChannels.BROWSE\_GET\_PROCESS\_BY\_ID, 31, null \\
\hline
TU-1790 & \ok & result deve essere false \\
\hline
TU-1791 & \ok & mockGateway.invoke non deve essere stato chiamato \\
\hline
TU-1792 & \ok & mockElectronBridge.invoke deve essere stato chiamato con 'ipc:search:text', object containing { text: 'Ricerca Valida' }, any AbortSignal \\
\hline
TU-1793 & \ok & mockElectronBridge.invoke deve essere stato chiamato con 'ipc:search:advanced', object containing { filter: expect.objectContaining({ logicOperator: 'AND', }), }, any AbortSignal \\
\hline
TU-1794 & \ok & mockElectronBridge.invoke deve essere stato chiamato con 'ipc:search:semantic', object containing { text: 'Concetto chiave', useSemanticSearch: true }, any AbortSignal \\
\hline
TU-1795 & \ok & errorBox deve essere truthy \\
\hline
TU-1796 & \ok & errorBox.nativeElement.textContent deve contenere 'IPC\_TIMEOUT' \\
\hline
TU-1797 & \ok & mockElectronBridge.invoke non deve essere stato chiamato \\
\hline
TU-1798 & \ok & mockCache.get deve essere stato chiamato \\
\hline
TU-1799 & \ok & mockBridge.invoke non deve essere stato chiamato \\
\hline
TU-1800 & \ok & result deve essere uguale a mockResults \\
\hline
TU-1801 & \ok & mockBridge.invoke deve essere stato chiamato con 'ipc:search:text', mockQuery, abortController.signal \\
\hline
TU-1802 & \ok & mockCache.set deve essere stato chiamato \\
\hline
TU-1803 & \ok & result deve essere uguale a mockResults \\
\hline
TU-1804 & \ok & mockCache.invalidatePrefix deve essere stato chiamato con 'search:text' \\
\hline
TU-1805 & \ok & mockErrorHandler.handle non deve essere stato chiamato \\
\hline
TU-1806 & \ok & mockErrorHandler.handle deve essere stato chiamato con mockError \\
\hline
TU-1807 & \ok & (errorResult as any).message deve essere 'IPC timeout' \\
\hline
TU-1808 & \ok & mockBridge.invoke deve essere stato chiamato con 'ipc:search:advanced', mapped, abortController.signal \\
\hline
TU-1809 & \ok & result deve essere uguale a mockResults \\
\hline
TU-1810 & \ok & mockCache.get deve essere stato chiamato con `search:advanced:\${JSON.stringify(mapped)}` \\
\hline
TU-1811 & \ok & cachedResult deve essere uguale a mockResults \\
\hline
TU-1812 & \ok & mockBridge.invoke non deve essere stato chiamato con 'ipc:search:advanced', expect.anything(), expect.anything() \\
\hline
TU-1813 & \ok & result deve essere uguale a [] \\
\hline
TU-1814 & \ok & mockBridge.invoke deve essere stato chiamato con 'ipc:search:semantic', mockQuery, abortController.signal \\
\hline
TU-1815 & \ok & result deve essere uguale a mockResults \\
\hline
TU-1816 & \ok & mockCache.get deve essere stato chiamato con `search:semantic:\${JSON.stringify(mockQuery)}` \\
\hline
TU-1817 & \ok & cachedResult deve essere uguale a mockResults \\
\hline
TU-1818 & \ok & mockBridge.invoke deve essere stato chiamato con 'search:get-custom-metadata-keys', 1, abortController.signal \\
\hline
TU-1819 & \ok & mockCache.set deve essere stato chiamato \\
\hline
TU-1820 & \ok & result deve essere uguale a ['NomeCliente', 'NumeroPratica'] \\
\hline
TU-1821 & \ok & mockCache.get deve essere stato chiamato con 'search:custom-metadata-keys:1' \\
\hline
TU-1822 & \ok & mockBridge.invoke non deve essere stato chiamato con 'search:get-custom-metadata-keys', expect.anything(), expect.anything() \\
\hline
TU-1823 & \ok & result deve essere uguale a cachedKeys \\
\hline
TU-1824 & \ok & mockBridge.invoke non deve essere stato chiamato \\
\hline
TU-1825 & \ok & mockCache.invalidatePrefix deve essere stato chiamato con 'search:text' \\
\hline
TU-1826 & \ok & errorResult deve essere undefined \\
\hline
TU-1827 & \ok & nextCalled deve essere false \\
\hline
TU-1828 & \ok & errorCount deve essere 0 \\
\hline
TU-1829 & \ok & hasMeaningfulAdvancedFilters(filters) deve essere true \\
\hline
TU-1830 & \ok & dto non deve essere null \\
\hline
TU-1831 & \ok & dto?.filter.items deve contenere un elemento uguale a object containing { logicOperator: 'AND', items: expect.arrayContaining([ expect.objectContaining({ path: 'DocumentoInformatico.CustomMetadata.NomeCliente', operator: 'LIKE', value: '\%Alfa Solutions S.p.A.\%', }), ]), } \\
\hline
TU-1832 & \ok & dto non deve essere null \\
\hline
TU-1833 & \ok & dto?.filter.items deve contenere un elemento uguale a object containing { logicOperator: 'AND', items: expect.arrayContaining([ expect.objectContaining({ path: 'DocumentoAmministrativoInformatico.CustomMetadata.NomeCliente', operator: 'LIKE', value: '\%Alfa Solutions S.p.A.\%', }), ]), } \\
\hline
TU-1834 & \ok & dto non deve essere null \\
\hline
TU-1835 & \ok & dto?.filter.items deve contenere un elemento uguale a object containing { logicOperator: 'AND', items: expect.arrayContaining([ expect.objectContaining({ path: 'DocumentoInformatico.CustomMetadata.NomeCliente', operator: 'LIKE', value: '\%Alfa Solutions S.p.A.\%', }), ]), } \\
\hline
TU-1836 & \ok & dto non deve essere null \\
\hline
TU-1837 & \ok & dto?.filter.items deve contenere un elemento uguale a object containing { logicOperator: 'AND', items: expect.arrayContaining([ expect.objectContaining({ path: 'AggregazioneDocumentaliInformatiche.CustomMetadata.NomeCliente', operator: 'LIKE', value: '\%Alfa Solutions S.p.A.\%', }), ]), } \\
\hline
TU-1838 & \ok & conditions deve contenere un elemento uguale a object containing { path: 'DocumentoInformatico.CustomMetadata.NomeCliente', operator: 'LIKE', value: '\%Alfa\%', } \\
\hline
TU-1839 & \ok & conditions deve contenere un elemento uguale a object containing { path: 'DocumentoInformatico.CustomMetadata.Pagato', operator: 'LIKE', value: '\%true\%', } \\
\hline
TU-1840 & \ok & dto non deve essere null \\
\hline
TU-1841 & \ok & dto?.filter.items deve contenere un elemento uguale a object containing { logicOperator: 'AND', items: expect.arrayContaining([ expect.objectContaining({ logicOperator: 'OR', items: expect.arrayContaining([ expect.objectContaining({ path: 'DocumentoInformatico.CustomMetadata.NomeCliente', operator: 'LIKE', value: '\%Alfa\%', }), expect.objectContaining({ path: 'DocumentoAmministrativoInformatico.CustomMetadata.NomeCliente', operator: 'LIKE', value: '\%Alfa\%', }), expect.objectContaining({ path: 'AggregazioneDocumentaliInformatiche.CustomMetadata.NomeCliente', operator: 'LIKE', value: '\%Alfa\%', }), ]), }), ]), } \\
\hline
TU-1842 & \ok & dto deve essere null \\
\hline
TU-1843 & \ok & hasMeaningfulAdvancedFilters({ common: {}, diDai: {}, aggregate: {}, customMeta: null, subject: [], } as any) deve essere false \\
\hline
TU-1844 & \ok & dto non deve essere null \\
\hline
TU-1845 & \ok & dto?.filter.items deve contenere un elemento uguale a object containing { logicOperator: 'AND', items: expect.arrayContaining([ expect.objectContaining({ path: 'DocumentoInformatico.Note', operator: 'LIKE', value: '\%protocollo\%', }), ]), } \\
\hline
TU-1846 & \ok & dto?.filter.items deve contenere un elemento uguale a object containing { logicOperator: 'AND', items: expect.arrayContaining([ expect.objectContaining({ path: 'DocumentoInformatico.CustomMetadata.NomeCliente', operator: 'LIKE', value: '\%Acme\%', }), ]), } \\
\hline
TU-1847 & \ok & dto?.filter.items deve contenere un elemento uguale a object containing { logicOperator: 'AND', items: expect.arrayContaining([ expect.objectContaining({ path: 'AggregazioneDocumentale.TipoAgg.TipoAggregazione', operator: 'EQ', value: 'Fascicolo', }), expect.objectContaining({ path: 'AggregazioneDocumentale.TipoAgg.IdAggregazione', operator: 'EQ', value: 'A-123', }), ]), } \\
\hline
TU-1848 & \ok & subjectGroup deve essere definito \\
\hline
TU-1849 & \ok & subjectGroup.items[0] deve essere uguale a object containing { path: 'DocumentoInformatico.Soggetti.Ruolo', operator: 'ELEM\_MATCH', } \\
\hline
TU-1850 & \ok & nestedItems deve contenere un elemento uguale a object containing { path: 'Altro.TipoRuolo', operator: 'EQ', value: 'Altro', } \\
\hline
TU-1851 & \ok & nestedItems non deve contenere un elemento uguale a object containing { path: 'Altro.PF', } \\
\hline
TU-1852 & \ok & nestedItems deve contenere un elemento uguale a object containing { path: 'Altro.PF.Cognome', operator: 'LIKE', value: '\%Rossi\%', } \\
\hline
TU-1853 & \ok & nestedItems deve contenere un elemento uguale a object containing { path: 'Altro.PF.Nome', operator: 'LIKE', value: '\%Paolo\%', } \\
\hline
TU-1854 & \ok & nestedItems deve contenere un elemento uguale a object containing { path: 'ResponsabileGestioneDocumentale.TipoRuolo', operator: 'EQ', value: 'Responsabile della Gestione Documentale', } \\
\hline
TU-1855 & \ok & nestedItems deve contenere un elemento uguale a object containing { path: 'ResponsabileGestioneDocumentale.PF.Cognome', operator: 'LIKE', value: '\%Colombo\%', } \\
\hline
TU-1856 & \ok & nestedItems deve contenere un elemento uguale a object containing { path: 'ResponsabileGestioneDocumentale.PF.Nome', operator: 'LIKE', value: '\%Marco\%', } \\
\hline
TU-1857 & \ok & subjectGroup deve essere definito \\
\hline
TU-1858 & \ok & subjectGroup.items deve avere lunghezza 1 \\
\hline
TU-1859 & \ok & subjectGroup.items[0] deve essere uguale a object containing { path: 'DocumentoAmministrativoInformatico.Soggetti.Ruolo', operator: 'ELEM\_MATCH', } \\
\hline
TU-1860 & \ok & subjectGroup deve essere definito \\
\hline
TU-1861 & \ok & paths deve essere uguale a [ 'AggregazioneDocumentale.Soggetti.Ruolo', 'DocumentoAmministrativoInformatico.Soggetti.Ruolo', 'DocumentoInformatico.Soggetti.Ruolo', ] \\
\hline
TU-1862 & \ok & nestedItems deve contenere un elemento uguale a object containing { path: 'Destinatario.PG.CodiceFiscale\_PartitaIva', operator: 'EQ', value: '31140103768', } \\
\hline
TU-1863 & \ok & items deve contenere un elemento uguale a object containing { path: 'DocumentoInformatico.DatiDiRegistrazione.TipoRegistro.Repertorio\_Registro.TipoRegistro', operator: 'EQ', value: String.raw`Repertorio-Registro`, } \\
\hline
TU-1864 & \ok & items deve contenere un elemento uguale a object containing { path: 'DocumentoInformatico.ModalitaDiFormazione', operator: 'EQ', value: 'creazione tramite utilizzo di strumenti software che assicurino la produzione di documenti nei formati previsti in allegato 2', } \\
\hline
TU-1865 & \ok & items deve contenere un elemento uguale a object containing { path: 'AggregazioneDocumentale.TipoAgg.TipoAggregazione', operator: 'EQ', value: 'Serie Di Fascicoli', } \\
\hline
TU-1866 & \ok & items deve contenere un elemento uguale a object containing { path: 'DocumentoInformatico.Riservato', operator: 'EQ', value: true, } \\
\hline
TU-1867 & \ok & state.error deve essere uguale a appError \\
\hline
TU-1868 & \ok & state.loading deve essere false \\
\hline
TU-1869 & \ok & state.isSearching deve essere false \\
\hline
TU-1870 & \ok & mockErrorHandler.handle deve essere stato chiamato con rawError \\
\hline
TU-1871 & \ok & mockTelemetry.trackError deve essere stato chiamato con appError \\
\hline
TU-1872 & \ok & mockSearchChannel.search non deve essere stato chiamato \\
\hline
TU-1873 & \ok & state.validationErrors.get('dataDa') deve essere uguale a validationError \\
\hline
TU-1874 & \ok & mockSearchChannel.searchAdvanced non deve essere stato chiamato \\
\hline
TU-1875 & \ok & state.loading deve essere false \\
\hline
TU-1876 & \ok & state.isSearching deve essere false \\
\hline
TU-1877 & \ok & abortSpy deve essere stato chiamato \\
\hline
TU-1878 & \ok & mockSearchChannel.searchSemantic non deve essere stato chiamato \\
\hline
TU-1879 & \ok & mockSearchChannel.searchSemantic deve essere stato chiamato \\
\hline
TU-1880 & \ok & facade.getState()().results deve essere uguale a mockResults \\
\hline
TU-1881 & \ok & state.results deve essere uguale a mockResults \\
\hline
TU-1882 & \ok & state.loading deve essere false \\
\hline
TU-1883 & \ok & state.isSearching deve essere false \\
\hline
TU-1884 & \ok & mockTelemetry.trackEvent deve essere stato chiamato con any String \\
\hline
TU-1885 & \ok & mockLiveAnnouncer.announce deve essere stato chiamato con 'Trovati 2 risultati', 'polite' \\
\hline
TU-1886 & \ok & state.results deve essere uguale a mockResults \\
\hline
TU-1887 & \ok & state.validationErrors.size deve essere 0 \\
\hline
TU-1888 & \ok & mockSearchChannel.searchAdvanced deve essere stato chiamato con dummyFilters, any AbortSignal \\
\hline
TU-1889 & \ok & mockSearchChannel.searchAdvanced non deve essere stato chiamato \\
\hline
TU-1890 & \ok & facade.getState()().results deve essere uguale a [] \\
\hline
TU-1891 & \ok & facade.getState()().loading deve essere false \\
\hline
TU-1892 & \ok & facade.getState()().isSearching deve essere false \\
\hline
TU-1893 & \ok & mockErrorHandler.handle non deve essere stato chiamato \\
\hline
TU-1894 & \ok & mockTelemetry.trackError non deve essere stato chiamato \\
\hline
TU-1895 & \ok & facade.getState()().loading deve essere false \\
\hline
TU-1896 & \ok & facade.getState()().query deve essere uguale a newQuery \\
\hline
TU-1897 & \ok & facade.getState()().filters deve essere uguale a newFilters \\
\hline
TU-1898 & \ok & mockSearchChannel.search deve essere stato chiamato \\
\hline
TU-1899 & \ok & facade.getState()().results deve essere uguale a mockResults \\
\hline
TU-1900 & \ok & mockSearchChannel.search non deve essere stato chiamato \\
\hline
TU-1901 & \ok & mockSearchChannel.search non deve essere stato chiamato \\
\hline
TU-1902 & \ok & mockSearchChannel.search deve essere stato chiamato 1 volte \\
\hline
TU-1903 & \ok & facade.getState()().validationErrors.size deve essere 0 \\
\hline
TU-1904 & \ok & mockSearchChannel.search deve essere stato chiamato 1 volte \\
\hline
TU-1905 & \ok & facade.getState()().isSearching deve essere true \\
\hline
TU-1906 & \ok & mockSearchChannel.search deve essere stato chiamato 1 volte \\
\hline
TU-1907 & \ok & facade.getState()().isSearching deve essere true \\
\hline
TU-1908 & \ok & mockSearchChannel.searchAdvanced deve essere stato chiamato 1 volte \\
\hline
TU-1909 & \ok & facade.getState()().isSearching deve essere true \\
\hline
TU-1910 & \ok & mockSearchChannel.searchSemantic deve essere stato chiamato 1 volte \\
\hline
TU-1911 & \ok & mockSearchChannel.search non deve essere stato chiamato \\
\hline
TU-1912 & \ok & () => facade.cancelSearch() non deve lanciare un errore \\
\hline
TU-1913 & \ok & facade.getState()().isSearching deve essere false \\
\hline
TU-1914 & \ok & abortSpy deve essere stato chiamato \\
\hline
TU-1915 & \ok & mockErrorHandler.handle deve essere stato chiamato con rawError \\
\hline
TU-1916 & \ok & facade.getState()().error deve essere uguale a { code: 'ERR\_ADV', message: 'Errore' } \\
\hline
TU-1917 & \ok & mockErrorHandler.handle deve essere stato chiamato con rawError \\
\hline
TU-1918 & \ok & facade.getState()().error deve essere uguale a { code: 'ERR\_SEM', message: 'Errore' } \\
\hline
TU-1919 & \ok & mockErrorHandler.handle deve essere stato chiamato con null \\
\hline
TU-1920 & \ok & componentClass deve essere ProcessResultCardComponent \\
\hline
TU-1921 & \ok & componentClass deve essere ClassResultCardComponent \\
\hline
TU-1922 & \ok & componentClass deve essere DocumentResultCardComponent \\
\hline
TU-1923 & \ok & componentClass deve essere DocumentResultCardComponent \\
\hline
TU-1924 & \ok & componentClass deve essere AggregateResultCardComponent \\
\hline
TU-1925 & \ok & componentClass deve essere ProcessResultCardComponent \\
\hline
TU-1926 & \ok & componentClass deve essere DocumentResultCardComponent \\
\hline
TU-1927 & \ok & componentClass deve essere DocumentResultCardComponent \\
\hline
TU-1928 & \ok & currentState.status deve essere IndexingStatus.IDLE \\
\hline
TU-1929 & \ok & currentState.progressPercentage deve essere 0 \\
\hline
TU-1930 & \ok & mockIpcGateway.getIndexingStatus deve essere stato chiamato 1 volte \\
\hline
TU-1931 & \ok & currentState.status deve essere IndexingStatus.INDEXING \\
\hline
TU-1932 & \ok & currentState.progressPercentage deve essere 45 \\
\hline
TU-1933 & \ok & facade.getStatus()().status deve essere IndexingStatus.TIMEOUT \\
\hline
TU-1934 & \ok & facade.getStatus()().status deve essere IndexingStatus.INDEXING \\
\hline
TU-1935 & \ok & facade.getStatus()().status deve essere IndexingStatus.TIMEOUT \\
\hline
TU-1936 & \ok & facade.getStatus()().status deve essere IndexingStatus.READY \\
\hline
TU-1937 & \ok & facade.getStatus()().status deve essere IndexingStatus.ERROR \\
\hline
TU-1938 & \ok & facade.getStatus()().status deve essere IndexingStatus.ERROR \\
\hline
TU-1939 & \ok & roles.length deve essere maggiore di 0 \\
\hline
TU-1940 & \ok & titolareRole deve essere definito \\
\hline
TU-1941 & \ok & titolareRole?.label deve essere SubjectRoleType.AMMINISTRAZIONE\_TITOLARE \\
\hline
TU-1942 & \ok & types deve essere uguale a [SubjectType.PG, SubjectType.PAI, SubjectType.PAE] \\
\hline
TU-1943 & \ok & types deve essere uguale a [] \\
\hline
TU-1944 & \ok & roles.length deve essere 16 \\
\hline
TU-1945 & \ok & autoreRole?.label deve essere 'Autore' \\
\hline
TU-1946 & \ok & regRole?.label deve essere 'Sogg. Registrazione (DI)' \\
\hline
TU-1947 & \ok & types deve essere uguale a [SubjectType.PAI, SubjectType.PAE] \\
\hline
TU-1948 & \ok & types deve essere uguale a [] \\
\hline
TU-1949 & \ok & roles.length deve essere maggiore di 0 \\
\hline
TU-1950 & \ok & regRole deve essere definito \\
\hline
TU-1951 & \ok & regRole?.label deve essere SubjectRoleType.AMMINISTRAZIONE\_REGISTRAZIONE \\
\hline
TU-1952 & \ok & types deve essere uguale a [SubjectType.PAI] \\
\hline
TU-1953 & \ok & types deve essere uguale a [] \\
\hline
TU-1954 & \ok & roles.length deve essere maggiore di 0 \\
\hline
TU-1955 & \ok & autoreRole deve essere definito \\
\hline
TU-1956 & \ok & autoreRole?.label deve essere SubjectRoleType.AUTORE \\
\hline
TU-1957 & \ok & types deve essere uguale a [ SubjectType.PF, SubjectType.PG, SubjectType.PAI, SubjectType.PAE ] \\
\hline
TU-1958 & \ok & types deve essere uguale a [] \\
\hline
TU-1959 & \ok & component deve essere truthy \\
\hline
TU-1960 & \ok & component.form.contains('procedimento') deve essere true \\
\hline
TU-1961 & \ok & component.form.contains('assegnazione') deve essere true \\
\hline
TU-1962 & \ok & component.fasiFormArray.length deve essere 0 \\
\hline
TU-1963 & \ok & emitSpy deve essere stato chiamato \\
\hline
TU-1964 & \ok & emitted?.idAggregazione deve essere 'AGG-001' \\
\hline
TU-1965 & \ok & component.fasiFormArray.length deve essere 1 \\
\hline
TU-1966 & \ok & component.fasiFormArray.at(0).get('tipoFase') deve essere definito \\
\hline
TU-1967 & \ok & component.fasiFormArray.length deve essere 0 \\
\hline
TU-1968 & \ok & addBtn deve essere truthy \\
\hline
TU-1969 & \ok & component.fasiFormArray.length deve essere 1 \\
\hline
TU-1970 & \ok & removeBtn deve essere truthy \\
\hline
TU-1971 & \ok & component.fasiFormArray.length deve essere 0 \\
\hline
TU-1972 & \ok & component.form.get('tipoAggregazione')?.value deve essere AggregationType.FASCICOLO \\
\hline
TU-1973 & \ok & component.form.get('procedimento.materia')?.value deve essere 'Ambiente' \\
\hline
TU-1974 & \ok & component.fasiFormArray.length deve essere 1 \\
\hline
TU-1975 & \ok & patchSpy non deve essere stato chiamato \\
\hline
TU-1976 & \ok & error?.message deve essere 'URI non valido' \\
\hline
TU-1977 & \ok & () => fixture.detectChanges() non deve lanciare un errore \\
\hline
TU-1978 & \ok & errorEl?.nativeElement.textContent deve contenere 'La data di inizio non può essere successiva alla data di fine.' \\
\hline
TU-1979 & \ok & errors.some((el) => el.nativeElement.textContent.includes( 'La data di inizio non può essere successiva alla data di fine.', ), ) deve essere true \\
\hline
TU-1980 & \ok & errors.some((el) => el.nativeElement.textContent.includes( 'La data di inizio non può essere successiva alla data di fine.', ), ) deve essere true \\
\hline
TU-1981 & \ok & nextSpy deve essere stato chiamato \\
\hline
TU-1982 & \ok & completeSpy deve essere stato chiamato \\
\hline
TU-1983 & \ok & projected deve essere truthy \\
\hline
TU-1984 & \ok & projected.nativeElement.textContent deve contenere 'Contenuto Reale' \\
\hline
TU-1985 & \ok & fixture.debugElement.query(By.css('.async-loading-state')) deve essere null \\
\hline
TU-1986 & \ok & fixture.debugElement.query(By.css('.async-error-state')) deve essere null \\
\hline
TU-1987 & \ok & loader deve essere truthy \\
\hline
TU-1988 & \ok & loader.nativeElement.textContent deve contenere 'Caricamento in corso' \\
\hline
TU-1989 & \ok & fixture.debugElement.query(By.css('.projected-content')) deve essere null \\
\hline
TU-1990 & \ok & errorState deve essere truthy \\
\hline
TU-1991 & \ok & errorState.nativeElement.textContent deve contenere 'Il server non risponde' \\
\hline
TU-1992 & \ok & hostComponent.onRetry deve essere stato chiamato \\
\hline
TU-1993 & \ok & errorState.nativeElement.textContent deve contenere 'Si è verificato un errore imprevisto.' \\
\hline
TU-1994 & \ok & component deve essere truthy \\
\hline
TU-1995 & \ok & component.form.contains('chiaveDescrittiva') deve essere true \\
\hline
TU-1996 & \ok & component.form.contains('classificazione') deve essere true \\
\hline
TU-1997 & \ok & component.form.contains('conservazione') deve essere true \\
\hline
TU-1998 & \ok & component.form.contains('note') deve essere true \\
\hline
TU-1999 & \ok & component.form.contains('tipoDocumento') deve essere true \\
\hline
TU-2000 & \ok & emitSpy deve essere stato chiamato \\
\hline
TU-2001 & \ok & emittedValue?.note deve essere 'Test di nota' \\
\hline
TU-2002 & \ok & component.form.get('tipoDocumento')?.value deve essere DocumentType.DOCUMENTO\_INFORMATICO \\
\hline
TU-2003 & \ok & component.form.get('note')?.value deve essere 'Da verificare' \\
\hline
TU-2004 & \ok & component.form.get('classificazione.codice')?.value deve essere 'TIT.1' \\
\hline
TU-2005 & \ok & () => component.ngOnChanges({}) non deve lanciare un errore \\
\hline
TU-2006 & \ok & error deve essere truthy \\
\hline
TU-2007 & \ok & error?.message deve essere 'Oggetto obbligatorio' \\
\hline
TU-2008 & \ok & errorDiv deve essere truthy \\
\hline
TU-2009 & \ok & errorDiv.nativeElement.textContent deve contenere 'Oggetto obbligatorio' \\
\hline
TU-2010 & \ok & component.getError('note') deve essere undefined \\
\hline
TU-2011 & \ok & component.getError('classificazione.codice') deve essere undefined \\
\hline
TU-2012 & \ok & component.entries.at(0) deve essere original \\
\hline
TU-2013 & \ok & component.entries.at(0) non deve essere original \\
\hline
TU-2014 & \ok & component.entries.at(0).value deve essere uguale a { field: 'MetaKey2', value: 'MetaValue2' } \\
\hline
TU-2015 & \ok & component.entries.length deve essere 2 \\
\hline
TU-2016 & \ok & component.entries.length deve essere 2 \\
\hline
TU-2017 & \ok & component.entries.at(1) deve essere draftControl \\
\hline
TU-2018 & \ok & emitSpy deve essere stato chiamato con [{ field: 'Cliente', value: '' }] \\
\hline
TU-2019 & \ok & emitSpy deve essere stato chiamato con null \\
\hline
TU-2020 & \ok & component.entries.length deve essere 0 \\
\hline
TU-2021 & \ok & component.entries.length deve essere 1 \\
\hline
TU-2022 & \ok & component.entries.at(0).value deve essere uguale a { field: 'Chiave1', value: 'Valore1' } \\
\hline
TU-2023 & \ok & emitSpy non deve essere stato chiamato \\
\hline
TU-2024 & \ok & component.entries.at(0) deve essere firstControl \\
\hline
TU-2025 & \ok & component.entries.length deve essere 1 \\
\hline
TU-2026 & \ok & component.entries.length deve essere 0 \\
\hline
TU-2027 & \ok & emitSpy deve essere stato chiamato con [{ field: 'Chiave', value: '' }] \\
\hline
TU-2028 & \ok & emitSpy deve essere stato chiamato con null \\
\hline
TU-2029 & \ok & component.getError(0, 'field')?.message deve essere 'Chiave non valida' \\
\hline
TU-2030 & \ok & component.hasAnyError() deve essere true \\
\hline
TU-2031 & \ok & component.hasAnyError() deve essere false \\
\hline
TU-2032 & \ok & component.getSelectableKeys(0) deve essere uguale a ['Zeta', 'A', 'B'] \\
\hline
TU-2033 & \ok & component.getSelectableKeyLabel('Process.PreservationSession.DocumentsStats.DocumentsFilesCount') deve essere 'Documents Files Count' \\
\hline
TU-2034 & \ok & component deve essere truthy \\
\hline
TU-2035 & \ok & component.form.contains('registrazione') deve essere true \\
\hline
TU-2036 & \ok & component.form.contains('identificativoFormato') deve essere true \\
\hline
TU-2037 & \ok & component.form.contains('verifica') deve essere true \\
\hline
TU-2038 & \ok & component.tracciatureFormArray.length deve essere 0 \\
\hline
TU-2039 & \ok & emitSpy deve essere stato chiamato \\
\hline
TU-2040 & \ok & emittedValue?.tipologia deve essere 'Test Tipo' \\
\hline
TU-2041 & \ok & component.tracciatureFormArray.length deve essere 0 \\
\hline
TU-2042 & \ok & component.tracciatureFormArray.length deve essere 1 \\
\hline
TU-2043 & \ok & component.tracciatureFormArray.at(0).value deve essere uguale a { tipoModifica: null, dataModifica: null, oraModifica: null, idVersionePrec: null, } \\
\hline
TU-2044 & \ok & component.tracciatureFormArray.length deve essere 2 \\
\hline
TU-2045 & \ok & component.tracciatureFormArray.length deve essere 1 \\
\hline
TU-2046 & \ok & component.form.get('tipologia')?.value deve essere 'Bando' \\
\hline
TU-2047 & \ok & component.form.get('registrazione.tipologiaRegistro')?.value deve essere RegisterType.PROTOCOLLO \\
\hline
TU-2048 & \ok & component.tracciatureFormArray.length deve essere 2 \\
\hline
TU-2049 & \ok & component.tracciatureFormArray.at(0).get('tipoModifica')?.value deve essere ModificationType.INTEGRAZIONE \\
\hline
TU-2050 & \ok & component.tracciatureFormArray.at(1).get('tipoModifica')?.value deve essere ModificationType.ANNULLAMENTO \\
\hline
TU-2051 & \ok & () => component.ngOnChanges({}) non deve lanciare un errore \\
\hline
TU-2052 & \ok & component.tracciatureFormArray.length deve essere 1 \\
\hline
TU-2053 & \ok & component.tracciatureFormArray.length deve essere 0 \\
\hline
TU-2054 & \ok & error deve essere truthy \\
\hline
TU-2055 & \ok & error?.message deve essere 'Numero invalido' \\
\hline
TU-2056 & \ok & component.getError('campoInesistente') deve essere undefined \\
\hline
TU-2057 & \ok & component.getError('registrazione.tipologiaFlusso') deve essere undefined \\
\hline
TU-2058 & \ok & error?.code deve essere 'ERR\_CONF\_DT\_002' \\
\hline
TU-2059 & \ok & addBtn deve essere truthy \\
\hline
TU-2060 & \ok & component.tracciatureFormArray.length deve essere 1 \\
\hline
TU-2061 & \ok & removeBtn deve essere truthy \\
\hline
TU-2062 & \ok & component.tracciatureFormArray.length deve essere 0 \\
\hline
TU-2063 & \ok & component deve essere truthy \\
\hline
TU-2064 & \ok & spy deve essere stato chiamato con 'ciao' \\
\hline
TU-2065 & \ok & input deve essere truthy \\
\hline
TU-2066 & \ok & errorCmp.nativeElement.textContent deve contenere 'Data obbligatoria' \\
\hline
TU-2067 & \ok & select deve essere truthy \\
\hline
TU-2068 & \ok & options.length deve essere 2 \\
\hline
TU-2069 & \ok & spy deve essere stato chiamato \\
\hline
TU-2070 & \ok & input deve essere truthy \\
\hline
TU-2071 & \ok & spy deve essere stato chiamato \\
\hline
TU-2072 & \ok & select deve essere truthy \\
\hline
TU-2073 & \ok & spy deve essere stato chiamato \\
\hline
TU-2074 & \ok & component deve essere truthy \\
\hline
TU-2075 & \ok & component.message deve essere 'Si è verificato un errore.' \\
\hline
TU-2076 & \ok & wrapperEl deve essere truthy \\
\hline
TU-2077 & \ok & iconEl.nativeElement.className deve contenere 'bi-exclamation-triangle-fill' \\
\hline
TU-2078 & \ok & msgEl.nativeElement.textContent deve contenere 'Si è verificato un errore.' \\
\hline
TU-2079 & \ok & msgEl.nativeElement.textContent deve contenere 'Credenziali non valide.' \\
\hline
TU-2080 & \ok & nameEl.textContent deve contenere 'Fascicolo Personale Mario Rossi' \\
\hline
TU-2081 & \ok & idEl.textContent deve contenere 'ID Fascicolo: FAS-456' \\
\hline
TU-2082 & \ok & badgeEl.textContent deve contenere 'AGGREGAZIONE DOCUMENTALE' \\
\hline
TU-2083 & \ok & scoreBadge.textContent deve contenere '0.72' \\
\hline
TU-2084 & \ok & mockAction deve essere stato chiamato con mockAggregate \\
\hline
TU-2085 & \ok & nameEl.textContent deve contenere 'Risorse Umane' \\
\hline
TU-2086 & \ok & idEl.textContent deve contenere 'ID Classe: CLS-789' \\
\hline
TU-2087 & \ok & badgeEl.textContent deve contenere 'CLASSE' \\
\hline
TU-2088 & \ok & mockAction deve essere stato chiamato con mockClass \\
\hline
TU-2089 & \ok & nameEl.textContent deve contenere 'Relazione Tecnica' \\
\hline
TU-2090 & \ok & idEl.textContent deve contenere 'DOC-999' \\
\hline
TU-2091 & \ok & scoreBadge deve essere truthy \\
\hline
TU-2092 & \ok & scoreBadge deve essere falsy \\
\hline
TU-2093 & \ok & mockAction deve essere stato chiamato con mockDocument \\
\hline
TU-2094 & \ok & mockAction deve essere stato chiamato con mockDocument \\
\hline
TU-2095 & \ok & nameEl.textContent deve contenere 'Approvazione Bilancio' \\
\hline
TU-2096 & \ok & idEl.textContent deve contenere 'ID Processo: PRC-123' \\
\hline
TU-2097 & \ok & badgeEl.textContent deve contenere 'PROCESS' \\
\hline
TU-2098 & \ok & mockAction deve essere stato chiamato con mockProcess \\
\hline
TU-2099 & \ok & mockAction deve essere stato chiamato con mockProcess \\
\hline
TU-2100 & \ok & result deve essere uguale a { isSemanticForced: true, activeType: SearchQueryType.FREE, } \\
\hline
TU-2101 & \ok & result deve essere uguale a { isSemanticForced: false, activeType: requestedType, } \\
\hline
TU-2102 & \ok & component deve essere truthy \\
\hline
TU-2103 & \ok & component.form.value.text deve essere '' \\
\hline
TU-2104 & \ok & component.form.value.type deve essere MockQueryType.FREE \\
\hline
TU-2105 & \ok & component.form.value.useSemanticSearch deve essere false \\
\hline
TU-2106 & \ok & component.form.getRawValue().text deve essere 'Fascicolo 123' \\
\hline
TU-2107 & \ok & component.form.getRawValue().type deve essere MockQueryType.PROCESS\_ID \\
\hline
TU-2108 & \ok & component.form.getRawValue().useSemanticSearch deve essere true \\
\hline
TU-2109 & \ok & emitSpy non deve essere stato chiamato \\
\hline
TU-2110 & \ok & component.form.disabled deve essere true \\
\hline
TU-2111 & \ok & component.form.enabled deve essere true \\
\hline
TU-2112 & \ok & spinner deve essere truthy \\
\hline
TU-2113 & \ok & spinner.nativeElement.textContent deve contenere 'Ricerca in corso' \\
\hline
TU-2114 & \ok & nextSpy deve essere stato chiamato \\
\hline
TU-2115 & \ok & completeSpy deve essere stato chiamato \\
\hline
TU-2116 & \ok & component deve essere truthy \\
\hline
TU-2117 & \ok & component.results.length deve essere 0 \\
\hline
TU-2118 & \ok & emptyState deve essere truthy \\
\hline
TU-2119 & \ok & emptyState.componentInstance.message deve essere 'Nessun dato' \\
\hline
TU-2120 & \ok & emptyState.nativeElement.textContent deve contenere 'Nessun dato' \\
\hline
TU-2121 & \ok & emptyState deve essere null \\
\hline
TU-2122 & \ok & mockFactory.getComponentForType deve essere stato chiamato con ElementType.DOCUMENTO\_INFORMATICO \\
\hline
TU-2123 & \ok & mockFactory.getComponentForType deve essere stato chiamato con ElementType.DOCUMENTO\_AMMINISTRATIVO\_INFORMATICO \\
\hline
TU-2124 & \ok & stubs.length deve essere 2 \\
\hline
TU-2125 & \ok & inputs['result'] deve essere result \\
\hline
TU-2126 & \ok & inputs['isSemanticSearch'] deve essere true \\
\hline
TU-2127 & \ok & typeof inputs['onSelectAction'] deve essere 'function' \\
\hline
TU-2128 & \ok & emitSpy deve essere stato chiamato 1 volte \\
\hline
TU-2129 & \ok & emitSpy deve essere stato chiamato con result \\
\hline
TU-2130 & \ok & inputs['isSemanticSearch'] deve essere false \\
\hline
TU-2131 & \ok & component.statusClass deve essere 'status-unknown' \\
\hline
TU-2132 & \ok & component.statusLabel deve essere 'Sconosciuto' \\
\hline
TU-2133 & \ok & textEl.nativeElement.textContent deve contenere 'Sconosciuto' \\
\hline
TU-2134 & \ok & component.statusClass deve essere 'status-ready' \\
\hline
TU-2135 & \ok & component.statusLabel deve essere 'Pronto' \\
\hline
TU-2136 & \ok & component.statusClass deve essere 'status-error' \\
\hline
TU-2137 & \ok & component.statusLabel deve essere 'Errore Indice' \\
\hline
TU-2138 & \ok & component.statusClass deve essere 'status-indexing' \\
\hline
TU-2139 & \ok & component.statusLabel deve essere 'In Aggiornamento...' \\
\hline
TU-2140 & \ok & progressEl deve essere null \\
\hline
TU-2141 & \ok & progressEl deve essere truthy \\
\hline
TU-2142 & \ok & progressEl.nativeElement.textContent deve contenere '45\%' \\
\hline
TU-2143 & \ok & component.statusClass deve essere 'status-unknown' \\
\hline
TU-2144 & \ok & component.statusLabel deve essere 'CUSTOM\_STATE' \\
\hline
TU-2145 & \ok & component deve essere truthy \\
\hline
TU-2146 & \ok & component.fields.length deve essere 0 \\
\hline
TU-2147 & \ok & renderedText deve contenere 'Seleziona una tipologia' \\
\hline
TU-2148 & \ok & component.fields.length deve essere 3 \\
\hline
TU-2149 & \ok & component.form.contains('cognomePF') deve essere true \\
\hline
TU-2150 & \ok & component.form.contains('nomePF') deve essere true \\
\hline
TU-2151 & \ok & inputs.length deve essere 3 \\
\hline
TU-2152 & \ok & cognomeControl?.hasValidator deve essere truthy \\
\hline
TU-2153 & \ok & cognomeControl?.invalid deve essere true \\
\hline
TU-2154 & \ok & component.form.contains('codiceFiscalePartitaIvaPG') deve essere true \\
\hline
TU-2155 & \ok & component.form.contains('codiceFiscalePG') deve essere false \\
\hline
TU-2156 & \ok & component.form.contains('partitaIvaPG') deve essere false \\
\hline
TU-2157 & \ok & component.fields.some((f) => f.key === 'codiceFiscalePartitaIvaPG') deve essere true \\
\hline
TU-2158 & \ok & component.fields.some((f) => f.label === 'Codice Fiscale / Partita IVA') deve essere true \\
\hline
TU-2159 & \ok & component.form.value.cognomePF deve essere 'Rossi' \\
\hline
TU-2160 & \ok & component.form.value.nomePF deve essere 'Mario' \\
\hline
TU-2161 & \ok & () => { component.ngOnChanges({ details: new SimpleChange(null, { cognomePF: 'Rossi' }, true), }) } non deve lanciare un errore \\
\hline
TU-2163 & \ok & emitSpy deve essere stato chiamato con object containing { cognomePF: 'Bianchi' } \\
\hline
TU-2164 & \ok & nextSpy deve essere stato chiamato \\
\hline
TU-2165 & \ok & completeSpy deve essere stato chiamato \\
\hline
TU-2166 & \ok & component.currentStep() deve essere 0 \\
\hline
TU-2167 & \ok & component.selectedRole() deve essere null \\
\hline
TU-2168 & \ok & expectedRole deve essere definito \\
\hline
TU-2169 & \ok & setRoleSpy deve essere stato chiamato con expectedRole \\
\hline
TU-2170 & \ok & setTypeSpy deve essere stato chiamato con allowedTypes[idx] \\
\hline
TU-2171 & \ok & resetSpy deve essere stato chiamato \\
\hline
TU-2172 & \ok & resetSpy non deve essere stato chiamato \\
\hline
TU-2173 & \ok & component.selectedRole() deve essere SubjectRoleType.AUTORE \\
\hline
TU-2174 & \ok & component.currentStep() deve essere 2 \\
\hline
TU-2175 & \ok & component.selectedType() deve essere SubjectType.PAI \\
\hline
TU-2176 & \ok & component.currentDetails() deve essere null \\
\hline
TU-2177 & \ok & component.currentStep() deve essere 3 \\
\hline
TU-2178 & \ok & component.currentStep() deve essere 1 \\
\hline
TU-2179 & \ok & component.currentStep() deve essere 3 \\
\hline
TU-2180 & \ok & component.currentStep() deve essere 3 \\
\hline
TU-2181 & \ok & component.currentStep() deve essere 0 \\
\hline
TU-2182 & \ok & component.selectedRole() deve essere null \\
\hline
TU-2183 & \ok & component.selectedType() deve essere null \\
\hline
TU-2184 & \ok & component.subjectsList() deve essere uguale a mockSubject \\
\hline
TU-2185 & \ok & component.currentDetails() deve essere uguale a mockDetails \\
\hline
TU-2186 & \ok & emitSpy deve essere stato chiamato \\
\hline
TU-2187 & \ok & component.subjectsList().length deve essere 1 \\
\hline
TU-2188 & \ok & component.subjectsList()[0].role deve essere SubjectRoleType.AUTORE \\
\hline
TU-2189 & \ok & component.subjectsList()[0].details deve essere uguale a mockDetails \\
\hline
TU-2190 & \ok & component.currentStep() deve essere 0 \\
\hline
TU-2191 & \ok & component.subjectsList().length deve essere 1 \\
\hline
TU-2192 & \ok & component.subjectsList()[0].role deve essere 'DESTINATARIO' \\
\hline
TU-2193 & \ok & emitSpy deve essere stato chiamato con component.subjectsList() \\
\hline
TU-2194 & \ok & buttons.length deve essere component.availableRoles().length \\
\hline
TU-2195 & \ok & actionBtns.length deve essere 1 \\
\hline
TU-2196 & \ok & component.currentStep() deve essere 0 \\
\hline
TU-2197 & \ok & buttons.length deve essere component.allowedTypes().length \\
\hline
TU-2198 & \ok & actionBtns.length deve essere 2 \\
\hline
TU-2199 & \ok & component.currentStep() deve essere 1 \\
\hline
TU-2200 & \ok & component.currentStep() deve essere 0 \\
\hline
TU-2201 & \ok & detailForm deve essere truthy \\
\hline
TU-2202 & \ok & actionBtns.length deve essere 3 \\
\hline
TU-2203 & \ok & emitSpy deve essere stato chiamato con object containing { customMeta: mockEntries } \\
\hline
TU-2204 & \ok & emitSpy deve essere stato chiamato con object containing { subject: mockSubject } \\
\hline
TU-2205 & \ok & emitSubmitSpy non deve essere stato chiamato \\
\hline
TU-2206 & \ok & submitBtn.disabled deve essere true \\
\hline
TU-2207 & \ok & emitSubmitSpy deve essere stato chiamato \\
\hline
TU-2208 & \ok & mockValidatorFn deve essere stato chiamato \\
\hline
TU-2209 & \ok & emitValidationSpy deve essere stato chiamato \\
\hline
TU-2210 & \ok & emitFiltersSpy deve essere stato chiamato con object containing { common: expect.objectContaining({ testo: 'prova reattività' }), } \\
\hline
TU-2211 & \ok & component.currentValidationResult deve essere null \\
\hline
TU-2212 & \ok & emitResetSpy deve essere stato chiamato \\
\hline
TU-2213 & \ok & patchSpy deve essere stato chiamato con newFilters, { emitEvent: false } \\
\hline
TU-2214 & \ok & resetSpy deve essere stato chiamato con {}, { emitEvent: false } \\
\hline
TU-2215 & \ok & component.subjectResetCounter deve essere initialCounter + 1 \\
\hline
TU-2216 & \ok & resetSpy deve essere stato chiamato con {}, { emitEvent: false } \\
\hline
TU-2217 & \ok & () => { component.filters = { common: { t: 1 } } as any } non deve lanciare un errore \\
\hline
TU-2219 & \ok & state.disableAggregate deve essere false \\
\hline
TU-2220 & \ok & state.disableDiDai deve essere false \\
\hline
TU-2221 & \ok & state.motivoBloccoAggregate deve essere '' \\
\hline
TU-2222 & \ok & state.motivoBloccoDiDai deve essere '' \\
\hline
TU-2223 & \ok & state.disableAggregate deve essere true \\
\hline
TU-2224 & \ok & state.motivoBloccoAggregate deve essere 'Filtri bloccati: hai selezionato il Tipo Documento.' \\
\hline
TU-2225 & \ok & state.disableDiDai deve essere false \\
\hline
TU-2226 & \ok & state.disableAggregate deve essere true \\
\hline
TU-2227 & \ok & state.motivoBloccoAggregate deve essere 'Filtri bloccati: hai selezionato il Tipo Documento.' \\
\hline
TU-2228 & \ok & state.disableAggregate deve essere true \\
\hline
TU-2229 & \ok & state.motivoBloccoAggregate deve essere 'Filtri bloccati: stai compilando i parametri DiDai.' \\
\hline
TU-2230 & \ok & state.disableDiDai deve essere true \\
\hline
TU-2231 & \ok & state.motivoBloccoDiDai deve essere 'Filtri bloccati: hai selezionato il Tipo Fascicolo.' \\
\hline
TU-2232 & \ok & state.disableAggregate deve essere false \\
\hline
TU-2233 & \ok & state.disableDiDai deve essere true \\
\hline
TU-2234 & \ok & state.motivoBloccoDiDai deve essere 'Filtri bloccati: stai compilando i parametri Aggregazione.' \\
\hline
TU-2235 & \ok & FilterRulesManager.hasValoriCompilati(null) deve essere false \\
\hline
TU-2236 & \ok & FilterRulesManager.hasValoriCompilati(undefined) deve essere false \\
\hline
TU-2237 & \ok & FilterRulesManager.hasValoriCompilati('just a string') deve essere false \\
\hline
TU-2238 & \ok & FilterRulesManager.hasValoriCompilati(123) deve essere false \\
\hline
TU-2239 & \ok & FilterRulesManager.hasValoriCompilati(true) deve essere false \\
\hline
TU-2240 & \ok & FilterRulesManager.hasValoriCompilati(new Date()) deve essere true \\
\hline
TU-2241 & \ok & FilterRulesManager.hasValoriCompilati({ a: null, b: undefined, c: '', d: ' ' }) deve essere false \\
\hline
TU-2242 & \ok & FilterRulesManager.hasValoriCompilati({ a: false }) deve essere false \\
\hline
TU-2243 & \ok & FilterRulesManager.hasValoriCompilati({ items: [] }) deve essere false \\
\hline
TU-2244 & \ok & FilterRulesManager.hasValoriCompilati({ wrapper: { nested: { prop: '' } } }) deve essere false \\
\hline
TU-2245 & \ok & FilterRulesManager.hasValoriCompilati({ dateProp: new Date() }) deve essere true \\
\hline
TU-2246 & \ok & FilterRulesManager.hasValoriCompilati({ wrapper: { nested: { prop: 'valid' } } }) deve essere true \\
\hline
TU-2247 & \ok & FilterRulesManager.hasValoriCompilati({ numberProp: 0 }) deve essere true \\
\hline
TU-2248 & \ok & FilterRulesManager.hasValoriCompilati({ booleanProp: true }) deve essere true \\
\hline
TU-2249 & \ok & FilterRulesManager.hasValoriCompilati({ arrayProp: [1, 2, 3] }) deve essere true \\
\hline
TU-2250 & \ok & component deve essere truthy \\
\hline
TU-2251 & \ok & mockFacade.getState deve essere stato chiamato \\
\hline
TU-2252 & \ok & mockFacade.loadCustomMetadataKeys deve essere stato chiamato con \\
\hline
TU-2253 & \ok & mockFacade.setQuery deve essere stato chiamato \\
\hline
TU-2254 & \ok & mockFacade.search deve essere stato chiamato \\
\hline
TU-2255 & \ok & mockFacade.searchSemantic deve essere stato chiamato \\
\hline
TU-2256 & \ok & mockFacade.search non deve essere stato chiamato \\
\hline
TU-2257 & \ok & mockFacade.search deve essere stato chiamato \\
\hline
TU-2258 & \ok & mockFacade.searchSemantic non deve essere stato chiamato \\
\hline
TU-2259 & \ok & component.externalValidation?.isValid deve essere false \\
\hline
TU-2260 & \ok & component.externalValidation?.errors.get('aggregate.dataApertura')?.[0] deve essere uguale a validationError \\
\hline
TU-2261 & \ok & mockFacade.search deve essere stato chiamato \\
\hline
TU-2262 & \ok & mockFacade.validateFilters deve essere stato chiamato con filters \\
\hline
TU-2263 & \ok & mockFacade.setFilters deve essere stato chiamato con mockFilters \\
\hline
TU-2264 & \ok & mockFacade.setFilters deve essere stato chiamato con { common: {}, diDai: {}, aggregate: {}, customMeta: null, subject: [], } \\
\hline
TU-2265 & \ok & mockFacade.searchAdvanced deve essere stato chiamato con mockFilters \\
\hline
TU-2266 & \ok & mockFacade.search non deve essere stato chiamato \\
\hline
TU-2267 & \ok & component.externalValidation deve essere uguale a result \\
\hline
TU-2268 & \ok & loader deve essere truthy \\
\hline
TU-2269 & \ok & loader.nativeElement.textContent deve contenere 'Ricerca in corso' \\
\hline
TU-2270 & \ok & errorBox deve essere truthy \\
\hline
TU-2271 & \ok & errorBox.nativeElement.textContent deve contenere 'Errore di Rete' \\
\hline
TU-2272 & \ok & retryBtn deve essere truthy \\
\hline
TU-2273 & \ok & emptyState deve essere truthy \\
\hline
TU-2274 & \ok & emptyState.nativeElement.textContent deve contenere 'Nessun risultato trovato' \\
\hline
TU-2275 & \ok & resultsComponent deve essere truthy \\
\hline
TU-2276 & \ok & pageTitle.nativeElement.textContent deve contenere 'Trovati 1 risultati' \\
\hline
TU-2277 & \ok & getIntegrityClass(IntegrityStatusEnum.VALID) deve essere 'integrity--valid' \\
\hline
TU-2278 & \ok & getIntegrityClass(IntegrityStatusEnum.INVALID) deve essere 'integrity--invalid' \\
\hline
TU-2279 & \ok & getIntegrityClass(IntegrityStatusEnum.UNKNOWN) deve essere 'integrity--unknown' \\
\hline
TU-2280 & \ok & getIntegrityClass('NON\_EXISTENT\_STATUS' as IntegrityStatusEnum) deve essere 'integrity--unknown' \\
\hline
TU-2281 & \ok & strategies.length deve essere 5 \\
\hline
TU-2282 & \ok & result.isValid deve essere false \\
\hline
TU-2283 & \ok & result.errors.has('testField') deve essere true \\
\hline
TU-2284 & \ok & result.isValid deve essere true \\
\hline
TU-2285 & \ok & result.errors.size deve essere 0 \\
\hline
TU-2286 & \ok & result.size deve essere 0 \\
\hline
TU-2287 & \ok & result.size deve essere 0 \\
\hline
TU-2288 & \ok & result.size deve essere 0 \\
\hline
TU-2289 & \ok & result.size deve essere 1 \\
\hline
TU-2290 & \ok & errors deve essere definito \\
\hline
TU-2291 & \ok & errors?.[0].code deve essere 'AGG\_CNTR\_001' \\
\hline
TU-2292 & \ok & result.size deve essere 1 \\
\hline
TU-2293 & \ok & errors deve essere definito \\
\hline
TU-2294 & \ok & errors?.[0].code deve essere 'AGG\_CNTR\_002' \\
\hline
TU-2295 & \ok & result.size deve essere 1 \\
\hline
TU-2296 & \ok & errors deve essere definito \\
\hline
TU-2297 & \ok & errors?.[0].code deve essere 'AGG\_CNTR\_003' \\
\hline
TU-2298 & \ok & result.size deve essere 2 \\
\hline
TU-2299 & \ok & result.has('aggregate.tipoFascicolo') deve essere true \\
\hline
TU-2300 & \ok & result.has('aggregate.assegnazione') deve essere true \\
\hline
TU-2301 & \ok & result.has('diDai.registrazione.dataRegistrazione') deve essere true \\
\hline
TU-2302 & \ok & result.get('diDai.registrazione.dataRegistrazione')?.[0].code deve essere 'ERR\_CONF\_DT\_001' \\
\hline
TU-2303 & \ok & result.has('diDai.tracciatureModifiche.0.dataModifica') deve essere true \\
\hline
TU-2304 & \ok & result.get('diDai.tracciatureModifiche.0.dataModifica')?.[0].code deve essere 'ERR\_CONF\_DT\_002' \\
\hline
TU-2305 & \ok & result.size deve essere 0 \\
\hline
TU-2306 & \ok & result.size deve essere 0 \\
\hline
TU-2307 & \ok & result.size deve essere 0 \\
\hline
TU-2308 & \ok & result.size deve essere 0 \\
\hline
TU-2309 & \ok & result.has('diDai.verifica.conformitaCopie') deve essere true \\
\hline
TU-2310 & \ok & result.get('diDai.verifica.conformitaCopie')?.[0].code deve essere 'ERR\_CONF\_FORM\_001' \\
\hline
TU-2311 & \ok & result.has('diDai.verifica.conformitaCopie') deve essere true \\
\hline
TU-2312 & \ok & result.size deve essere 0 \\
\hline
TU-2313 & \ok & result.size deve essere 0 \\
\hline
TU-2314 & \ok & result.size deve essere 0 \\
\hline
TU-2315 & \ok & result.size deve essere 0 \\
\hline
TU-2316 & \ok & result.has('diDai.registrazione.codiceRegistro') deve essere true \\
\hline
TU-2317 & \ok & result.get('diDai.registrazione.codiceRegistro')?.[0].code deve essere 'ERR\_CONF\_REG\_001' \\
\hline
TU-2318 & \ok & result.has('diDai.registrazione.numeroRegistrazione') deve essere true \\
\hline
TU-2319 & \ok & result.get('diDai.registrazione.numeroRegistrazione')?.[0].code deve essere 'ERR\_CONF\_REG\_002' \\
\hline
TU-2320 & \ok & result.size deve essere 2 \\
\hline
TU-2321 & \ok & result.has('diDai.registrazione.codiceRegistro') deve essere true \\
\hline
TU-2322 & \ok & result.has('diDai.registrazione.numeroRegistrazione') deve essere true \\
\hline
TU-2323 & \ok & result.size deve essere 0 \\
\hline
TU-2324 & \ok & result.size deve essere 0 \\
\hline
TU-2325 & \ok & result.size deve essere 0 \\
\hline
TU-2326 & \ok & result.size deve essere 0 \\
\hline
TU-2327 & \ok & result.size deve essere 0 \\
\hline
TU-2328 & \ok & result.has('common.dataDa') deve essere true \\
\hline
TU-2329 & \ok & result.get('common.dataDa')?.[0].code deve essere 'ERR\_RANGE\_001' \\
\hline
TU-2330 & \ok & strategy.validate(filters1).size deve essere 0 \\
\hline
TU-2331 & \ok & strategy.validate(filters2).size deve essere 0 \\
\hline
TU-2332 & \ok & result.has('aggregate.dataApertura') deve essere true \\
\hline
TU-2333 & \ok & result.get('aggregate.dataApertura')?.[0].code deve essere 'ERR\_RANGE\_001' \\
\hline
TU-2334 & \ok & strategy.validate(filters1).size deve essere 0 \\
\hline
TU-2335 & \ok & strategy.validate(filters2).size deve essere 0 \\
\hline
TU-2336 & \ok & result.has('aggregate.assegnazione.dataInizioAssegn') deve essere true \\
\hline
TU-2337 & \ok & result.get('aggregate.assegnazione.dataInizioAssegn')?.[0].code deve essere 'ERR\_RANGE\_001' \\
\hline
TU-2338 & \ok & strategy.validate(filters1).size deve essere 0 \\
\hline
TU-2339 & \ok & strategy.validate(filters2).size deve essere 0 \\
\hline
TU-2340 & \ok & result.has('aggregate.procedimento.fasi.0.dataInizioFase') deve essere true \\
\hline
TU-2341 & \ok & result.get('aggregate.procedimento.fasi.0.dataInizioFase')?.[0].code deve essere 'ERR\_RANGE\_001' \\
\hline
TU-2342 & \ok & strategy.validate(filters1).size deve essere 0 \\
\hline
TU-2343 & \ok & strategy.validate(filters2).size deve essere 0 \\
\hline
TU-2344 & \ok & result.has('common.numeroAllegatiMin') deve essere true \\
\hline
TU-2345 & \ok & result.get('common.numeroAllegatiMin')?.[0].code deve essere 'ERR\_RANGE\_002' \\
\hline
TU-2346 & \ok & strategy.validate(filters1).size deve essere 0 \\
\hline
TU-2347 & \ok & strategy.validate(filters2).size deve essere 0 \\
\hline
TU-2348 & \ok & gateway deve essere truthy \\
\hline
TU-2349 & \ok & gateway.checkDipIntegrity(1) deve lanciare `Electron API not available for channel IpcChannels.CHECK\_DIP\_INTEGRITY\_STATUS` \\
\hline
TU-2350 & \ok & mockInvoke deve essere stato chiamato con IpcChannels.CHECK\_DIP\_INTEGRITY\_STATUS, 1 \\
\hline
TU-2351 & \ok & result deve essere IntegrityStatusEnum.VALID \\
\hline
TU-2352 & \ok & mockInvoke deve essere stato chiamato con IpcChannels.CHECK\_DIP\_INTEGRITY\_STATUS, 10 \\
\hline
TU-2353 & \ok & res deve essere IntegrityStatusEnum.VALID \\
\hline
TU-2354 & \ok & mockInvoke deve essere stato chiamato con IpcChannels.CHECK\_DOCUMENT\_CLASS\_INTEGRITY\_STATUS, 11 \\
\hline
TU-2355 & \ok & res deve essere IntegrityStatusEnum.INVALID \\
\hline
TU-2356 & \ok & mockInvoke deve essere stato chiamato con IpcChannels.CHECK\_PROCESS\_INTEGRITY\_STATUS, 12 \\
\hline
TU-2357 & \ok & res deve essere IntegrityStatusEnum.UNKNOWN \\
\hline
TU-2358 & \ok & mockInvoke deve essere stato chiamato con IpcChannels.CHECK\_DOCUMENT\_INTEGRITY\_STATUS, 13 \\
\hline
TU-2359 & \ok & res deve essere IntegrityStatusEnum.VALID \\
\hline
TU-2360 & \ok & mockInvoke deve essere stato chiamato con IpcChannels.CHECK\_FILE\_INTEGRITY\_STATUS, 14 \\
\hline
TU-2361 & \ok & res deve essere IntegrityStatusEnum.INVALID \\
\hline
TU-2362 & \ok & mockInvoke deve essere stato chiamato con IpcChannels.BROWSE\_GET\_DOCUMENT\_CLASS\_BY\_DIP\_ID, 20 \\
\hline
TU-2363 & \ok & res deve essere uguale a mockResult \\
\hline
TU-2364 & \ok & mockInvoke deve essere stato chiamato con IpcChannels.BROWSE\_GET\_PROCESS\_BY\_DOCUMENT\_CLASS, 21 \\
\hline
TU-2365 & \ok & res deve essere uguale a mockResult \\
\hline
TU-2366 & \ok & mockInvoke deve essere stato chiamato con IpcChannels.BROWSE\_GET\_DOCUMENTS\_BY\_PROCESS, 22 \\
\hline
TU-2367 & \ok & res deve essere uguale a mockResult \\
\hline
TU-2368 & \ok & mockInvoke deve essere stato chiamato con IpcChannels.BROWSE\_GET\_FILE\_BY\_DOCUMENT, 23 \\
\hline
TU-2369 & \ok & res deve essere uguale a mockResult \\
\hline
TU-2370 & \ok & mockCache.invalidatePrefix deve essere stato chiamato con 'aggregate:' \\
\hline
TU-2371 & \ok & mockCache.invalidatePrefix deve essere stato chiamato con 'process:' \\
\hline
TU-2372 & \ok & mockCache.invalidatePrefix deve essere stato chiamato con 'document:' \\
\hline
TU-2373 & \ok & mockCache.invalidatePrefix deve essere stato chiamato con 'node-fallback:' \\
\hline
TU-2374 & \ok & facade.currentDipStatus()() deve essere IntegrityStatusEnum.VALID \\
\hline
TU-2375 & \ok & result deve essere IntegrityStatusEnum.VALID \\
\hline
TU-2376 & \ok & mockGateway.checkDocumentIntegrity deve essere stato chiamato con 12 \\
\hline
TU-2377 & \ok & mockCache.invalidate deve essere stato chiamato con 'document:12' \\
\hline
TU-2378 & \ok & result deve essere IntegrityStatusEnum.INVALID \\
\hline
TU-2379 & \ok & mockGateway.checkProcessIntegrity deve essere stato chiamato con 31 \\
\hline
TU-2380 & \ok & mockCache.invalidate deve essere stato chiamato con 'process:31' \\
\hline
TU-2381 & \ok & mockCache.invalidate deve essere stato chiamato con 'aggregate:31' \\
\hline
TU-2382 & \ok & mockCache.invalidatePrefix deve essere stato chiamato con 'document:' \\
\hline
TU-2383 & \ok & result deve essere IntegrityStatusEnum.UNKNOWN \\
\hline
TU-2384 & \ok & mockGateway.checkDocumentClassIntegrity deve essere stato chiamato con 22 \\
\hline
TU-2385 & \ok & mockCache.invalidate deve essere stato chiamato con 'node-fallback:DOCUMENT\_CLASS:22' \\
\hline
TU-2386 & \ok & mockCache.invalidatePrefix deve essere stato chiamato con 'process:' \\
\hline
TU-2387 & \ok & mockCache.invalidatePrefix deve essere stato chiamato con 'aggregate:' \\
\hline
TU-2388 & \ok & mockCache.invalidatePrefix deve essere stato chiamato con 'document:' \\
\hline
TU-2389 & \ok & result deve essere IntegrityStatusEnum.VALID \\
\hline
TU-2390 & \ok & mockGateway.checkProcessIntegrity deve essere stato chiamato con 44 \\
\hline
TU-2391 & \ok & mockCache.invalidate deve essere stato chiamato con 'aggregate:44' \\
\hline
TU-2392 & \ok & mockCache.invalidate deve essere stato chiamato con 'process:44' \\
\hline
TU-2393 & \ok & mockCache.invalidatePrefix deve essere stato chiamato con 'document:' \\
\hline
TU-2394 & \ok & result deve essere 'UNKNOWN' \\
\hline
TU-2395 & \ok & mockErrorHandler.handle deve essere stato chiamato 1 volte \\
\hline
TU-2396 & \ok & handledError.source deve essere 'IntegrityFacade.verifyItem' \\
\hline
TU-2397 & \ok & handledError.context deve essere uguale a { itemId: '55', itemType: 'DOCUMENT', action: 'CHECK\_ITEM\_INTEGRITY', } \\
\hline
TU-2398 & \ok & facade.error()() deve essere 'errore test' \\
\hline
TU-2399 & \ok & facade.currentDipStatus()() deve essere IntegrityStatusEnum.INVALID \\
\hline
TU-2400 & \ok & facade.dipClasses()() deve essere uguale a [{ id: 1, name: 'Classe A' }] \\
\hline
TU-2401 & \ok & facade.currentDipStatus()() deve essere null \\
\hline
TU-2402 & \ok & facade.dipClasses()() deve essere uguale a [] \\
\hline
TU-2403 & \ok & facade.error()() deve essere null \\
\hline
TU-2404 & \ok & facade.currentDipStatus()() deve essere IntegrityStatusEnum.VALID \\
\hline
TU-2405 & \ok & facade.dipClasses()() deve essere uguale a [{ id: 2, name: 'Classe B' }] \\
\hline
TU-2406 & \ok & component deve essere truthy \\
\hline
TU-2407 & \ok & validCard?.textContent deve contenere '5' \\
\hline
TU-2408 & \ok & validCard?.textContent deve contenere 'Elementi Integri' \\
\hline
TU-2409 & \ok & invalidCard?.textContent deve contenere '3' \\
\hline
TU-2410 & \ok & invalidCard?.textContent deve contenere 'Elementi Corrotti' \\
\hline
TU-2411 & \ok & unknownCard?.textContent deve contenere '2' \\
\hline
TU-2412 & \ok & unknownCard?.textContent deve contenere 'In Attesa di Verifica' \\
\hline
TU-2413 & \ok & component deve essere truthy \\
\hline
TU-2414 & \ok & compiled.querySelector('.valid-container') deve essere null \\
\hline
TU-2415 & \ok & compiled.querySelector('.valid-container') deve essere truthy \\
\hline
TU-2416 & \ok & items.length deve essere 1 \\
\hline
TU-2417 & \ok & items[0].textContent deve contenere 'Classe' \\
\hline
TU-2418 & \ok & items[0].textContent deve contenere 'Node Valid 1' \\
\hline
TU-2419 & \ok & items[0].textContent deve contenere 'Valido' \\
\hline
TU-2420 & \ok & compiled.querySelector('.status-pill')?.textContent deve contenere 'Invalido' \\
\hline
TU-2421 & \ok & compiled.querySelector('.status-pill--invalid') deve essere truthy \\
\hline
TU-2422 & \ok & compiled.querySelector('.error-group') deve essere truthy \\
\hline
TU-2423 & \ok & compiled.textContent deve contenere 'Processo A' \\
\hline
TU-2424 & \ok & compiled.textContent deve contenere 'Doc 1' \\
\hline
TU-2425 & \ok & compiled.querySelector('.error-group') deve essere null \\
\hline
TU-2426 & \ok & component.formatType('CLASS') deve essere 'Classe' \\
\hline
TU-2427 & \ok & component.formatType('PROCESS') deve essere 'Processo' \\
\hline
TU-2428 & \ok & component.formatType('DOCUMENT') deve essere 'Documento' \\
\hline
TU-2429 & \ok & component.formatType('SOME\_UNKNOWN\_TYPE') deve essere 'Documento' \\
\hline
TU-2430 & \ok & component.getStatusLabel(IntegrityStatusEnum.VALID) deve essere 'Valido' \\
\hline
TU-2431 & \ok & component.getStatusLabel(IntegrityStatusEnum.INVALID) deve essere 'Invalido' \\
\hline
TU-2432 & \ok & component.getStatusLabel(IntegrityStatusEnum.UNKNOWN) deve essere 'Non verificato' \\
\hline
TU-2433 & \ok & component.getStatusClass(IntegrityStatusEnum.VALID) deve essere 'status-pill status-pill--valid' \\
\hline
TU-2434 & \ok & component.getStatusClass(IntegrityStatusEnum.INVALID) deve essere 'status-pill status-pill--invalid' \\
\hline
TU-2435 & \ok & component.getStatusClass(IntegrityStatusEnum.UNKNOWN) deve essere 'status-pill status-pill--unknown' \\
\hline
TU-2436 & \ok & computedNodes.length deve essere 3 \\
\hline
TU-2437 & \ok & computedNodes[0].name deve essere 'Delta' \\
\hline
TU-2438 & \ok & computedNodes[1].name deve essere 'Alpha' \\
\hline
TU-2439 & \ok & computedNodes[2].name deve essere 'Zebra' \\
\hline
TU-2440 & \ok & errors.get('Classe non disponibile')?.processes[0].processName deve essere 'Processo non disponibile' \\
\hline
TU-2441 & \ok & errors.get('C1')?.processes[0].processName deve essere 'P1' \\
\hline
TU-2442 & \ok & errors.get('C1')?.totalDocuments deve essere 1 \\
\hline
TU-2443 & \ok & errors.get('C2')?.processes[0].processName deve essere 'P2' \\
\hline
TU-2444 & \ok & errors.get('C3')?.processes[0].processName deve essere 'Processo non disponibile' \\
\hline
TU-2445 & \ok & component deve essere truthy \\
\hline
TU-2446 & \ok & mockFacade.loadOverview deve essere stato chiamato con component.currentDipId \\
\hline
TU-2447 & \ok & compiled.querySelector('.progress-section') deve essere truthy \\
\hline
TU-2448 & \ok & compiled.querySelector('app-integrity-summary-cards') deve essere null \\
\hline
TU-2449 & \ok & compiled.querySelector('app-integrity-valid-panel') deve essere null \\
\hline
TU-2450 & \ok & compiled.querySelector('.progress-section') deve essere null \\
\hline
TU-2451 & \ok & compiled.querySelector('app-integrity-summary-cards') deve essere truthy \\
\hline
TU-2452 & \ok & compiled.querySelector('app-integrity-valid-panel') deve essere truthy \\
\hline
TU-2453 & \ok & mockFacade.verifyDip deve essere stato chiamato con component.currentDipId \\
\hline
TU-2454 & \ok & mockFacade.loadOverview deve essere stato chiamato 2 volte \\
\hline
TU-2455 & \ok & component.isFatal() deve essere false \\
\hline
TU-2456 & \ok & component.getIcon() deve essere 'bi-x-octagon-fill' \\
\hline
TU-2457 & \ok & component.getTitle() deve essere 'Si è verificato un errore' \\
\hline
TU-2458 & \ok & headerTitle.textContent.trim() deve essere 'Si è verificato un errore' \\
\hline
TU-2459 & \ok & dialog.classList.contains('fatal-dialog') deve essere false \\
\hline
TU-2460 & \ok & component.isFatal() deve essere true \\
\hline
TU-2461 & \ok & component.getIcon() deve essere 'bi-bug-fill' \\
\hline
TU-2462 & \ok & component.getTitle() deve essere 'Errore Critico del Sistema' \\
\hline
TU-2463 & \ok & dialog.classList.contains('fatal-dialog') deve essere true \\
\hline
TU-2464 & \ok & component.getIcon() deve essere 'bi-exclamation-triangle-fill' \\
\hline
TU-2465 & \ok & component.getTitle() deve essere 'Attenzione' \\
\hline
TU-2466 & \ok & component.getIcon() deve essere 'bi-info-circle-fill' \\
\hline
TU-2467 & \ok & component.getTitle() deve essere 'Avviso' \\
\hline
TU-2468 & \ok & monoElements.length deve essere 0 \\
\hline
TU-2469 & \ok & techDetails.length deve essere 0 \\
\hline
TU-2470 & \ok & retryButtons.length deve essere 0 \\
\hline
TU-2471 & \ok & errorMeta.textContent deve contenere 'SEARCH\_ENGINE\_ERROR' \\
\hline
TU-2472 & \ok & errorMeta.textContent deve contenere 'Backend API' \\
\hline
TU-2473 & \ok & techDetails.textContent deve contenere 'Stack trace details...' \\
\hline
TU-2474 & \ok & closeSpy deve essere stato chiamato 1 volte \\
\hline
TU-2475 & \ok & retrySpy deve essere stato chiamato 1 volte \\
\hline
TU-2476 & \ok & mockBridge.invoke deve essere stato chiamato con 'ipc:log:write', mockLogEntry \\
\hline
TU-2477 & \ok & errorResult deve essere mockError \\
\hline
TU-2478 & \ok & mockBridge.invoke deve essere stato chiamato 1 volte \\
\hline
TU-2479 & \ok & error.code deve essere ErrorCode.DOC\_BLOB\_LOAD\_ERROR \\
\hline
TU-2480 & \ok & error.message deve essere 'File non trovato' \\
\hline
TU-2481 & \ok & error.source deve essere 'DocFacade' \\
\hline
TU-2482 & \ok & error.category deve essere ErrorCategory.IO \\
\hline
TU-2483 & \ok & error.severity deve essere ErrorSeverity.ERROR \\
\hline
TU-2484 & \ok & error.recoverable deve essere true \\
\hline
TU-2485 & \ok & error.severity deve essere ErrorSeverity.FATAL \\
\hline
TU-2486 & \ok & error.recoverable deve essere false \\
\hline
TU-2487 & \ok & result.code deve essere ErrorCode.DOC\_FORMAT\_UNSUPPORTED \\
\hline
TU-2488 & \ok & result.message deve essere 'Formato .xyz non supportato' \\
\hline
TU-2489 & \ok & result.source deve essere 'ElectronRenderer' \\
\hline
TU-2490 & \ok & result.category deve essere ErrorCategory.IO \\
\hline
TU-2491 & \ok & result.code deve essere ErrorCode.SEARCH\_ENGINE\_ERROR \\
\hline
TU-2492 & \ok & result.message deve essere 'Errore misterioso nel motore di ricerca' \\
\hline
TU-2493 & \ok & result.detail deve essere rawError.stack \\
\hline
TU-2494 & \ok & result.code deve essere ErrorCode.IPC\_ERROR \\
\hline
TU-2495 & \ok & result.category deve essere ErrorCategory.IPC \\
\hline
TU-2496 & \ok & result.code deve essere ErrorCode.SEARCH\_ENGINE\_ERROR \\
\hline
TU-2497 & \ok & result.message deve essere 'Qualcosa è andato storto' \\
\hline
TU-2498 & \ok & result.source deve essere 'IPC Gateway' \\
\hline
TU-2499 & \ok & error.category deve essere ErrorCategory.IPC \\
\hline
TU-2500 & \ok & error.severity deve essere ErrorSeverity.ERROR \\
\hline
TU-2501 & \ok & error.recoverable deve essere false \\
\hline
TU-2502 & \ok & result.code deve essere ErrorCode.SEARCH\_ENGINE\_ERROR \\
\hline
TU-2503 & \ok & result deve essere uguale a { id: 1, name: 'Documento' } \\
\hline
TU-2504 & \ok & result deve essere null \\
\hline
TU-2505 & \ok & result deve essere null \\
\hline
TU-2506 & \ok & cacheService.get('manual-key') deve essere null \\
\hline
TU-2507 & \ok & cacheService.get('search/123') deve essere null \\
\hline
TU-2508 & \ok & cacheService.get('search/456') deve essere null \\
\hline
TU-2509 & \ok & cacheService.get('doc/789') deve essere 'dati documento' \\
\hline
TU-2510 & \ok & mockLoggingChannel.writeLog deve essere stato chiamato con { level: LogLevel.INFO, event: TelemetryEvent.SEARCH\_EXECUTED, timestamp: MOCK\_TIME, payload: null, } \\
\hline
TU-2511 & \ok & mockLoggingChannel.writeLog deve essere stato chiamato con { level: LogLevel.INFO, event: TelemetryMetric.SEARCH\_LATENCY\_MS, timestamp: MOCK\_TIME, payload: { durationMs: 350 }, } \\
\hline
TU-2512 & \ok & mockLoggingChannel.writeLog deve essere stato chiamato con { level: LogLevel.ERROR, event: 'APP\_ERROR', timestamp: MOCK\_TIME, payload: mockError, } \\
\hline
TU-2513 & \ok & () => service.trackEvent(TelemetryEvent.SEARCH\_EXECUTED) non deve lanciare un errore \\
\hline
TU-2514 & \ok & console.error deve essere stato chiamato con 'Fallimento scrittura telemetria via IPC:', logError \\
\hline
TU-2515 & \ok & mockBridge.invoke deve essere stato chiamato con 'ipc:log:write', mockLogEntry \\
\hline
TU-2516 & \ok & errorResult deve essere mockError \\
\hline
TU-2517 & \ok & mockBridge.invoke deve essere stato chiamato 1 volte \\
\hline
TU-2518 & \ok & error.code deve essere ErrorCode.DOC\_BLOB\_LOAD\_ERROR \\
\hline
TU-2519 & \ok & error.message deve essere 'File non trovato' \\
\hline
TU-2520 & \ok & error.source deve essere 'DocFacade' \\
\hline
TU-2521 & \ok & error.category deve essere ErrorCategory.IO \\
\hline
TU-2522 & \ok & error.severity deve essere ErrorSeverity.ERROR \\
\hline
TU-2523 & \ok & error.recoverable deve essere true \\
\hline
TU-2524 & \ok & error.severity deve essere ErrorSeverity.FATAL \\
\hline
TU-2525 & \ok & error.recoverable deve essere false \\
\hline
TU-2526 & \ok & result.code deve essere ErrorCode.DOC\_FORMAT\_UNSUPPORTED \\
\hline
TU-2527 & \ok & result.message deve essere 'Formato .xyz non supportato' \\
\hline
TU-2528 & \ok & result.source deve essere 'ElectronRenderer' \\
\hline
TU-2529 & \ok & result.category deve essere ErrorCategory.IO \\
\hline
TU-2530 & \ok & result.code deve essere ErrorCode.SEARCH\_ENGINE\_ERROR \\
\hline
TU-2531 & \ok & result.message deve essere 'Errore misterioso nel motore di ricerca' \\
\hline
TU-2532 & \ok & result.detail deve essere rawError.stack \\
\hline
TU-2533 & \ok & result.code deve essere ErrorCode.IPC\_ERROR \\
\hline
TU-2534 & \ok & result.category deve essere ErrorCategory.IPC \\
\hline
TU-2535 & \ok & result.code deve essere ErrorCode.SEARCH\_ENGINE\_ERROR \\
\hline
TU-2536 & \ok & result.message deve essere 'Qualcosa è andato storto' \\
\hline
TU-2537 & \ok & result.source deve essere 'IPC Gateway' \\
\hline
TU-2538 & \ok & error.category deve essere ErrorCategory.IPC \\
\hline
TU-2539 & \ok & error.severity deve essere ErrorSeverity.ERROR \\
\hline
TU-2540 & \ok & error.recoverable deve essere false \\
\hline
TU-2541 & \ok & result.code deve essere ErrorCode.SEARCH\_ENGINE\_ERROR \\
\hline
TU-2542 & \ok & result deve essere uguale a { id: 1, name: 'Documento' } \\
\hline
TU-2543 & \ok & result deve essere null \\
\hline
TU-2544 & \ok & result deve essere null \\
\hline
TU-2545 & \ok & cacheService.get('manual-key') deve essere null \\
\hline
TU-2546 & \ok & cacheService.get('search/123') deve essere null \\
\hline
TU-2547 & \ok & cacheService.get('search/456') deve essere null \\
\hline
TU-2548 & \ok & cacheService.get('doc/789') deve essere 'dati documento' \\
\hline
TU-2549 & \ok & mockCdkAnnouncer.announce deve essere stato chiamato con 'Ricerca completata', 'polite' \\
\hline
TU-2550 & \ok & mockCdkAnnouncer.announce deve essere stato chiamato con 'Ricerca completata', 'polite' \\
\hline
TU-2551 & \ok & mockCdkAnnouncer.announce deve essere stato chiamato con 'Errore di sistema critico', 'assertive' \\
\hline
TU-2552 & \ok & mockLoggingChannel.writeLog deve essere stato chiamato con { level: LogLevel.INFO, event: TelemetryEvent.SEARCH\_EXECUTED, timestamp: MOCK\_TIME, payload: null, } \\
\hline
TU-2553 & \ok & mockLoggingChannel.writeLog deve essere stato chiamato con { level: LogLevel.INFO, event: TelemetryMetric.SEARCH\_LATENCY\_MS, timestamp: MOCK\_TIME, payload: { durationMs: 350 }, } \\
\hline
TU-2554 & \ok & mockLoggingChannel.writeLog deve essere stato chiamato con { level: LogLevel.ERROR, event: 'APP\_ERROR', timestamp: MOCK\_TIME, payload: mockError, } \\
\hline
TU-2555 & \ok & () => service.trackEvent(TelemetryEvent.SEARCH\_EXECUTED) non deve lanciare un errore \\
\hline
TU-2556 & \ok & console.error deve essere stato chiamato con 'Fallimento scrittura telemetria via IPC:', logError \\
\hline
TU-2557 & \ok & component deve essere truthy \\
\hline
TU-2558 & \ok & simplifyCustomMetadataLabel('Process.PreservationSession.DocumentsStats.DocumentsFilesCount') deve essere 'Documents Files Count' \\
\hline
TU-2559 & \ok & simplifyCustomMetadataLabel('Process.End.Date') deve essere 'End Date' \\
\hline
TU-2560 & \ok & simplifyCustomMetadataLabel('Root.CustomMetadata.Process.Status') deve essere 'Status' \\
\hline
TU-2561 & \ok & simplifyCustomMetadataLabel('DocumentsFilesCount') deve essere 'Documents Files Count' \\
\hline
TU-2562 & \ok & toSpacedMetadataLabel('SoggettoAutoreDellaModifica') deve essere 'Soggetto Autore Della Modifica' \\
\hline
TU-2563 & \ok & toSpacedMetadataLabel('CodiceFiscale\_PartitaIva') deve essere 'Codice Fiscale Partita Iva' \\
\hline
TU-2564 & \ok & normalizeDisplayFileName('/869e1069-e50d-48e6-8191-1c677f3053a2/Allegato\_1.pdf') deve essere 'Allegato\_1.pdf' \\
\hline
TU-2565 & \ok & normalizeDisplayFileName('./Allegato\_1.pdf') deve essere 'Allegato\_1.pdf' \\
\hline
TU-2566 & \ok & normalizeDisplayFileName('folder\\sub\\FilePrincipale.pdf') deve essere 'FilePrincipale.pdf' \\
\hline
TU-2567 & \ok & normalizeDisplayFileName('documento.pdf') deve essere 'documento.pdf' \\
\hline
TU-2568 & \ok & normalizeDisplayFileName(' ') deve essere '' \\
\hline
TU-2569 & \ok & label deve essere 'Contratto.pdf' \\
\hline
TU-2570 & \ok & label deve essere 'DOC-2' \\
\hline
TU-2571 & \ok & entries deve essere uguale a [ { tipoDocumento: 'Documento', identificativo: 'Documento Wrapper', routeId: '10', }, ] \\
\hline
TU-2572 & \ok & mapDocumentsToIndexEntries(null) deve essere uguale a [] \\
\hline
TU-2573 & \ok & result deve avere lunghezza 2 \\
\hline
TU-2574 & \ok & result[0].name deve essere 'Oggetto' \\
\hline
TU-2575 & \ok & result deve essere uguale a [{ name: 'NomeDelDocumento', value: 'Doc.pdf' }] \\
\hline
TU-2576 & \ok & normalizeMetadataNodes(null) deve essere uguale a [] \\
\hline
TU-2577 & \ok & normalizeMetadataNodes(undefined) deve essere uguale a [] \\
\hline
TU-2578 & \ok & normalizeMetadataNodes({ name: 'x', value: 'y' }) deve essere uguale a [] \\
\hline
\rowcolor{white}
\caption{Report test di unità}
\label{tab:test-unita}
\end{longtable}

\end{document}